\documentclass[11pt,letterpaper]{article}

% Compile with XeLaTeX. Two sections from one source:
%   xelatex -jobname=EC111_FA2026_Syllabus_10H "\def\honors{1}\documentclass[11pt,letterpaper]{article}

% Compile with XeLaTeX. Two sections from one source:
%   xelatex -jobname=EC111_FA2026_Syllabus_10H "\def\honors{1}\documentclass[11pt,letterpaper]{article}

% Compile with XeLaTeX. Two sections from one source:
%   xelatex -jobname=EC111_FA2026_Syllabus_10H "\def\honors{1}\documentclass[11pt,letterpaper]{article}

% Compile with XeLaTeX. Two sections from one source:
%   xelatex -jobname=EC111_FA2026_Syllabus_10H "\def\honors{1}\input{EC111_FA2026_Syllabus.tex}"
%   xelatex -jobname=EC111_FA2026_Syllabus_20  "\input{EC111_FA2026_Syllabus.tex}"
\ifdefined\honors
  \newcommand{\sectioncode}{EC 111-10-H}
  \newcommand{\sectionname}{EC 111-10-H Principles of Microeconomics (Honors)}
  \newcommand{\meetingtime}{Monday/Wednesday, 3:30 PM -- 4:50 PM}
  \newcommand{\classroom}{Smith Technology Center 103}
  \newcommand{\finaltime}{Friday, December 11, 2026, 6:00 PM -- 8:00 PM}
  \newcommand{\mindtapsection}{Section 10-H of EC 111 Principles of Microeconomics Fall 2026}
  \newcommand{\mindtaplink}{https://student.cengage.com/course-link/MTPNS57GW6XN}
  \newcommand{\mindtapaccess}{the end of Sunday, September 13, 2026}
\else
  \newcommand{\sectioncode}{EC 111-20}
  \newcommand{\sectionname}{EC 111-20 Principles of Microeconomics}
  \newcommand{\meetingtime}{Monday/Wednesday, 9:30 AM -- 10:50 AM}
  \newcommand{\classroom}{Jennison Hall 412}
  \newcommand{\finaltime}{Friday, December 11, 2026, 8:00 AM -- 10:00 AM}
  \newcommand{\mindtapsection}{Section 20 of EC 111 Principles of Microeconomics Fall 2026}
  \newcommand{\mindtaplink}{https://student.cengage.com/course-link/MTPP8N7GWX7L}
  \newcommand{\mindtapaccess}{the end of Sunday, September 13, 2026}
\fi

% XeLaTeX for Turkish character support
\usepackage{fontspec}
\setmainfont{DejaVuSerif}[Extension=.ttf, BoldFont=DejaVuSerif-Bold, ItalicFont=DejaVuSerif-Italic, BoldItalicFont=DejaVuSerif-BoldItalic]
\setsansfont{DejaVuSans}[Extension=.ttf, BoldFont=DejaVuSans-Bold, ItalicFont=DejaVuSans-Oblique, BoldItalicFont=DejaVuSans-BoldOblique]

% Page layout
\usepackage[margin=1in]{geometry}
\usepackage{setspace}
\onehalfspacing

% Colors and styling
\usepackage[dvipsnames]{xcolor}
\definecolor{bentleyblue}{RGB}{0, 51, 102}
\definecolor{linkblue}{RGB}{0, 102, 204}

% Headers and sections
\usepackage{titlesec}
\titleformat{\section}{\Large\bfseries\color{bentleyblue}}{\thesection}{1em}{}[\titlerule]
\titleformat{\subsection}{\large\bfseries\color{bentleyblue!80}}{\thesubsection}{1em}{}
\titleformat{\subsubsection}{\normalsize\bfseries\color{bentleyblue!60}}{\thesubsubsection}{1em}{}

% Lists
\usepackage{enumitem}
\setlist[itemize]{leftmargin=*, itemsep=2pt, parsep=0pt}
\setlist[enumerate]{leftmargin=*, itemsep=2pt, parsep=0pt}

% Tables
\usepackage{booktabs}
\usepackage{tabularx}

% Links
\usepackage{hyperref}
\hypersetup{
    colorlinks=true,
    linkcolor=linkblue,
    urlcolor=linkblue,
    pdfauthor={Onur Altındağ},
    pdftitle={EC111 Principles of Microeconomics - Fall 2026 Syllabus}
}

% Header/Footer
\usepackage{fancyhdr}
\pagestyle{fancy}
\fancyhf{}
\fancyhead[L]{\small\color{gray}\sectioncode{} Principles of Microeconomics}
\fancyhead[R]{\small\color{gray}Fall 2026}
\fancyfoot[C]{\thepage}
\renewcommand{\headrulewidth}{0.4pt}

% Remove paragraph indent, add space between paragraphs
\setlength{\parindent}{0pt}
\setlength{\parskip}{6pt}

\begin{document}

% Title
\begin{center}
{\LARGE\bfseries\color{bentleyblue} EC 111: Principles of Microeconomics}\\[0.3cm]
{\Large Fall 2026 Syllabus}\\[0.5cm]
{\large Bentley University}
\end{center}

\vspace{0.5cm}

%-----------------------------------------------------------
\section*{Course Description}
%-----------------------------------------------------------

This course introduces the fundamental principles and tools of microeconomics: how individuals and firms make decisions when resources are scarce, and how markets coordinate those decisions. We study demand, supply, and market equilibrium; elasticity; the effects of government policies such as taxes and price controls; the efficiency of markets and the sources of market failure; the costs of production; and firm behavior in different market structures, from perfect competition to monopoly, monopolistic competition, and oligopoly.

The emphasis throughout is on economic reasoning: thinking at the margin, recognizing trade-offs, and using simple graphical and numerical models to make sense of real-world questions in business and policy.

\subsection*{Knowledge}
\begin{itemize}
    \item Scarcity, opportunity cost, trade-offs, and the gains from trade.
    \item Demand, supply, market equilibrium, and how markets respond to shocks.
    \item Elasticity and its role in pricing, revenue, and tax incidence.
    \item Consumer surplus, producer surplus, market efficiency, and the deadweight loss of taxes and price controls.
    \item Externalities, public goods, and the case for government intervention.
    \item Costs of production and firm behavior under perfect competition, monopoly, monopolistic competition, and oligopoly.
\end{itemize}

\subsection*{Skills}
\begin{itemize}
    \item Use supply and demand diagrams to predict the effects of changes in market conditions and policy.
    \item Compute and interpret elasticities, surpluses, costs, and profit-maximizing output and price.
    \item Read and evaluate economic arguments critically and communicate them clearly.
\end{itemize}

\subsection*{Perspectives}
\begin{itemize}
    \item Understand marginal decision making as the core of economic reasoning.
    \item Understand when markets work well, when they fail, and what policy can and cannot fix.
\end{itemize}

%-----------------------------------------------------------
\section*{Class Information}
%-----------------------------------------------------------

\begin{tabularx}{\textwidth}{@{}lX@{}}
\textbf{Instructor:} & Onur Altındağ \\
\textbf{Office:} & RAU 019F \\
\textbf{Email:} & \href{mailto:oaltindag@bentley.edu}{oaltindag@bentley.edu} \\
\textbf{Web:} & \url{https://www.onuraltindag.info/} \\
\textbf{Course page:} & \url{https://www.onuraltindag.info/teaching/ec111/} \\
\textbf{Password:} & \texttt{GoFalcons!} \\[0.3cm]
\textbf{Course:} & \sectionname \\
\textbf{Semester:} & Fall 2026 (August 31, 2026 – December 15, 2026) \\
\textbf{Meeting Time:} & \meetingtime \\
\textbf{Location:} & \classroom \\
\textbf{Delivery Mode:} & In-Person Lecture \\
\end{tabularx}

\vspace{0.3cm}

Office hours are on Mondays and Wednesdays, 11:00 AM -- 12:00 PM and 1:00 -- 3:00 PM, by appointment. Please \href{https://outlook.office.com/bookwithme/user/1c4e4f54e0f0493397e615b2c80a59c7@bentley.edu/meetingtype/rwPThvM0MkK5oGlK_t37nw2?bookingcode=17c9a7b2-69fa-4c48-85ea-ad404642334b&anonymous&ismsaljsauthenabled&ep=mlink}{book a 30-minute slot} to schedule a meeting. If you need to meet urgently or no slots are available, email me directly.

Announcements come by email from me to your Bentley address; check it regularly. Slides, lecture notes, handouts, and mock exams are posted on the course page above after each class. When you email me, put \textbf{EC111} in the subject line.

%-----------------------------------------------------------
\section*{Important Dates and Evaluation}
%-----------------------------------------------------------

\begin{tabularx}{\textwidth}{@{}Xr@{}}
\toprule
\textbf{Component} & \textbf{Weight} \\
\midrule
Weekly assignments + participation & 20\% \\
\textbf{Wednesday, October 7, 2026} – First Midterm & 25\% \\
\textbf{Wednesday, November 11, 2026} – Second Midterm & 25\% \\
\textbf{\finaltime} – Final Exam & 30\% \\
\bottomrule
\end{tabularx}

%-----------------------------------------------------------
\section*{Grading}
%-----------------------------------------------------------

You \textbf{MUST} attend all the midterms and the final as there will be no make-up exams in this course. The midterms are \textbf{NOT} cumulative. If you miss or are likely to miss a midterm due to an emergency, please contact me as soon as possible. You will need to provide supporting documentation/verification of your absence. I will re-weight your final exam accordingly if you have a valid excuse. Please note that family vacations or missing the school shuttle are not valid excuses. If a serious illness, a death in the family, or another difficult circumstance keeps you from an exam or assignment, contact the Office of Student Success (JEN 336, 781.891.2803); they will notify all of your instructors on your behalf.

The final exam \textbf{is} cumulative. If you miss the final exam due to an emergency, you will receive an incomplete for this course. Do not take this class if you know that you will not be able to attend the final exam.

Exams are closed book and closed notes. A basic (non-graphing) calculator is allowed; phones, graphing calculators, and any other device are not, and calculators may not be shared.

You will have \textbf{10 homework assignments} on MindTap, roughly one per week (none in exam weeks), due \textbf{Sunday 11:00 PM}. Your \textbf{two lowest scores are dropped}; the remaining 8 count equally. Each question allows three attempts; an assignment scoring at least \textbf{50\%} receives full credit, otherwise zero. The submission system is automated and will not accept anything after the deadline (even 5 minutes), and there are no exceptions. The questions are algorithmic and slightly different for everyone, so you are encouraged to work together.

\subsection*{Letter grades}
\begin{center}
\begin{tabular}{@{}llllll@{}}
\toprule
A & 94--100 & B+ & 87--89 & C+ & 77--79 \\
A$-$ & 90--93 & B & 83--86 & C & 73--76 \\
 & & B$-$ & 80--82 & C$-$ & 70--72 \\
\midrule
D+ & 67--69 & D & 63--66 & D$-$ & 60--62 \\
F & below 60 & & & & \\
\bottomrule
\end{tabular}
\end{center}
I may lower these cutoffs at the end of the semester; I will not raise them. I do not respond to requests for individual extra credit or for grade adjustments after the fact.

%-----------------------------------------------------------
\section*{Textbook and MindTap}
%-----------------------------------------------------------

The course uses \textbf{MindTap for Mankiw, \textit{Principles of Microeconomics}} (Cengage). The e-book, your homework assignments, and practice questions all live there; a hard copy of the textbook is optional. You \textbf{MUST} register, using your Bentley email, through the link for your section:

\begin{center}
\textbf{\mindtapsection}\\
\url{\mindtaplink}
\end{center}

You have free temporary access until \mindtapaccess; after that you will be prompted to purchase (if your course materials are included in your tuition, no purchase is needed). Registration help: \url{https://startstrong.cengage.com}; technical support: \url{https://support.cengage.com}. Register in the first week; MindTap problems are not an excuse for a missed assignment.

\ifdefined\honors
%-----------------------------------------------------------
\section*{Honors Reading List}
%-----------------------------------------------------------

In addition to the textbook, honors students read and discuss a short set of articles and book chapters over the semester. The list and discussion dates will be posted on the course page in the first weeks of the semester.

\fi
%-----------------------------------------------------------
\section*{Academic Integrity}
%-----------------------------------------------------------

All students are expected to adhere to Bentley's Academic Integrity policy, which includes Bentley's Honor Code (details are in the Undergraduate Student Handbook, the \href{https://catalog.bentley.edu/undergraduate/academic-policies-procedures/}{academic policies section of the catalog}, and the academic integrity course page on Brightspace in which all students and faculty are enrolled). The essence of the policy is that you should not represent someone else's work as your own: no plagiarism, no cheating on exams, no illicit collaboration. Failure to adhere to the policy has serious consequences, including course failure, suspension, or expulsion from the university. The best way to avoid a problem is to consult with me before taking an action that might constitute a violation.

\subsection*{Generative AI}
Nothing in this course is graded on writing, and the homework is algorithmic practice, so the line is simple. You may use generative AI tools (ChatGPT, Claude, Gemini, Copilot, and the like) to study: to explain a concept, to check your reasoning on a practice problem, or to work through a homework question you are stuck on. Treat them as a tutor, not a substitute; the exams are closed book and closed device, and the only way to pass them is to be able to do the work yourself. Any use of AI or any other device during an exam is a violation of the academic integrity policy.

%-----------------------------------------------------------
\section*{Inclusion}
%-----------------------------------------------------------

Everyone in this class has different life experiences and perspectives, and all are valid. As the instructor, I will behave maturely and respectfully in all our class-related engagements, and I expect the same from you. We all adhere to Bentley's standards of conduct, the \href{https://www.bentley.edu/about/mission-and-values}{Bentley Core Values}, in class and in all course work. If you feel that I or anyone in this class has acted outside those values, come and talk to me. If you would rather not, speak with someone in the Office of Student Success (JEN 336, 781.891.2803). My roster shows your name as it appears in Workday; tell me early in the semester if you go by a different name or pronoun.

%-----------------------------------------------------------
\section*{Support and Resources}
%-----------------------------------------------------------

\subsection*{Student Accessibility Services}
Bentley University abides by Section 504 of the Rehabilitation Act of 1973 and the Americans with Disabilities Act of 1990, which stipulate that no student shall be denied the benefits of an education solely by reason of a disability. If you have a hidden or visible disability that may require classroom accommodations, contact \href{https://www.bentley.edu/offices/student-accessibility-services}{Student Accessibility Services} (Office of Student Success, JEN 336, 781.891.2004) within the first four weeks of the semester. Send me your accommodation plan as soon as you have it; accommodations are arranged at least one week before each exam.

\subsection*{The Howard A. Winer '58 Lab for Economics, Accounting and Finance (LEAF)}
The LEAF provides a welcoming and inclusive learning environment where students are encouraged to seek academic support for their Economics, Accounting and Finance courses. Students utilizing the LEAF will find knowledgeable peer tutors ready to help their colleagues to thrive in the Bentley business curriculum. LEAF's free, drop-in tutoring takes place in-person at Morison 220D (rear entrance), opening September 14. For additional information visit \url{https://www.bentley.edu/centers/leaf}. LEAF is a supplement to, not a substitute for, office hours.

\subsection*{Bentley policies and procedures}
All courses at Bentley reflect the university's commitment to a set of core values and practices. You are expected to be familiar with:
\begin{itemize}
    \item \href{https://catalog.bentley.edu/undergraduate/academic-policies-procedures/}{Bentley's honor code and academic integrity system}
    \item \href{https://www.bentley.edu/equity-reporting}{Equity and bias reporting forms and procedures} (Bias Incident Response Team, Office of Institutional Equity)
    \item \href{https://www.bentley.edu/offices/student-accessibility-services}{Student disability accommodations}
    \item \href{https://catalog.bentley.edu/undergraduate/academic-policies-procedures/}{Religious observances policy}
    \item \href{https://www.bentley.edu/about/mission-and-values}{Bentley's core values}
\end{itemize}

\subsection*{Course materials}
Slides, notes, handouts, and exams for this course are for your own use during this semester. They are protected by copyright; do not post or distribute them outside the class without my written consent. The course page is password-protected; do not share the password outside the class.

%-----------------------------------------------------------
\section*{Attendance Policy}
%-----------------------------------------------------------

All students must attend the in-person classes. In-person attendance is essential because it allows for real-time interaction, immediate feedback, and the kind of dialogue that makes learning stick. I do not record lectures, so if you miss a class, you miss the material.

Remote attendance is not an option. This course follows Bentley's \href{https://catalog.bentley.edu/undergraduate/academic-policies-procedures/}{Academic Engagement/Attendance Policy}: students who do not attend or engage during the first week are dropped as no-shows, and two or more consecutive weeks of non-participation can lead to administrative withdrawal. If you must miss class, you are responsible for the material; read the chapter and the posted notes, then bring specific questions to office hours.

\textbf{Weather and closures.} If the university cancels classes, I will email you whether there is an alternative assignment or arrangement for that day.

\textbf{Electronic devices.} Phones away and laptops and tablets closed for the whole class. The work in this course is drawing and reading diagrams, which you learn by doing them by hand, and a screen in your line of sight costs the people around you as much as it costs you. Using a device during class counts against your participation grade. Exceptions are by my permission only; if you have one, you never need to explain it to anyone in the room. Photographing the board at the end of class is fine.

\textbf{Religious observances.} Bentley provides reasonable academic adjustments when a sincerely held religious observance conflicts with a class, exam, or deadline (see the \href{https://downloads.bentley.edu/general/Religious\%20Observances\%202026-2027.pdf}{religious observances calendar} and the \href{https://catalog.bentley.edu/undergraduate/academic-policies-procedures/}{policy}). If you anticipate a conflict, tell me in advance, ideally in the first two weeks of the semester, so we can agree on an adjustment.

%-----------------------------------------------------------
\section*{Tentative Schedule}
%-----------------------------------------------------------

\begin{itemize}
    \item \textbf{Week 1 (Aug 31, Sep 2):} Ten Principles of Economics – Ch. 1 \& Thinking Like an Economist – Ch. 2
    \item \textbf{Week 2 (Sep 9):} Interdependence and the Gains from Trade – Ch. 3 \textit{(no class Sep 7, Labor Day)}
    \item \textbf{Week 3 (Sep 14, 16):} The Market Forces of Supply and Demand – Ch. 4
    \item \textbf{Week 4 (Sep 21, 23):} Elasticity and Its Application – Ch. 5
    \item \textbf{Week 5 (Sep 28, 30):} Supply, Demand, and Government Policies – Ch. 6
    \item \textbf{Week 6 (Oct 5, 7):} Review / \textbf{First Midterm (Oct 7)}
    \item \textbf{Week 7 (Oct 14):} Consumers, Producers, and the Efficiency of Markets – Ch. 7 \textit{(no class Oct 12, Fall Break)}
    \item \textbf{Week 8 (Oct 19, 21):} The Costs of Taxation – Ch. 8 \& International Trade – Ch. 9
    \item \textbf{Week 9 (Oct 26, 28):} Externalities – Ch. 10 \& Public Goods and Common Resources – Ch. 11
    \item \textbf{Week 10 (Nov 2, 4):} The Costs of Production – Ch. 14
    \item \textbf{Week 11 (Nov 9, 11):} Review / \textbf{Second Midterm (Nov 11)}
    \item \textbf{Week 12 (Nov 16, 18):} Firms in Competitive Markets – Ch. 15
    \item \textbf{Week 13 (Nov 23):} Monopoly – Ch. 16 \textit{(no class Nov 25, Thanksgiving Break)}
    \item \textbf{Week 14 (Nov 30, Dec 2):} Monopolistic Competition – Ch. 17 \& Oligopoly – Ch. 18
    \item \textbf{Week 15 (Dec 7):} Review
    \item \textbf{Final Exam:} \finaltime
\end{itemize}

\textit{Note: Chapter references are from Mankiw's \textit{Principles of Microeconomics}, 10th edition (the MindTap e-book). The schedule, and if necessary other parts of this syllabus, may change during the semester; changes are announced in class and by email.}

\end{document}
"
%   xelatex -jobname=EC111_FA2026_Syllabus_20  "\documentclass[11pt,letterpaper]{article}

% Compile with XeLaTeX. Two sections from one source:
%   xelatex -jobname=EC111_FA2026_Syllabus_10H "\def\honors{1}\input{EC111_FA2026_Syllabus.tex}"
%   xelatex -jobname=EC111_FA2026_Syllabus_20  "\input{EC111_FA2026_Syllabus.tex}"
\ifdefined\honors
  \newcommand{\sectioncode}{EC 111-10-H}
  \newcommand{\sectionname}{EC 111-10-H Principles of Microeconomics (Honors)}
  \newcommand{\meetingtime}{Monday/Wednesday, 3:30 PM -- 4:50 PM}
  \newcommand{\classroom}{Smith Technology Center 103}
  \newcommand{\finaltime}{Friday, December 11, 2026, 6:00 PM -- 8:00 PM}
  \newcommand{\mindtapsection}{Section 10-H of EC 111 Principles of Microeconomics Fall 2026}
  \newcommand{\mindtaplink}{https://student.cengage.com/course-link/MTPNS57GW6XN}
  \newcommand{\mindtapaccess}{the end of Sunday, September 13, 2026}
\else
  \newcommand{\sectioncode}{EC 111-20}
  \newcommand{\sectionname}{EC 111-20 Principles of Microeconomics}
  \newcommand{\meetingtime}{Monday/Wednesday, 9:30 AM -- 10:50 AM}
  \newcommand{\classroom}{Jennison Hall 412}
  \newcommand{\finaltime}{Friday, December 11, 2026, 8:00 AM -- 10:00 AM}
  \newcommand{\mindtapsection}{Section 20 of EC 111 Principles of Microeconomics Fall 2026}
  \newcommand{\mindtaplink}{https://student.cengage.com/course-link/MTPP8N7GWX7L}
  \newcommand{\mindtapaccess}{the end of Sunday, September 13, 2026}
\fi

% XeLaTeX for Turkish character support
\usepackage{fontspec}
\setmainfont{DejaVuSerif}[Extension=.ttf, BoldFont=DejaVuSerif-Bold, ItalicFont=DejaVuSerif-Italic, BoldItalicFont=DejaVuSerif-BoldItalic]
\setsansfont{DejaVuSans}[Extension=.ttf, BoldFont=DejaVuSans-Bold, ItalicFont=DejaVuSans-Oblique, BoldItalicFont=DejaVuSans-BoldOblique]

% Page layout
\usepackage[margin=1in]{geometry}
\usepackage{setspace}
\onehalfspacing

% Colors and styling
\usepackage[dvipsnames]{xcolor}
\definecolor{bentleyblue}{RGB}{0, 51, 102}
\definecolor{linkblue}{RGB}{0, 102, 204}

% Headers and sections
\usepackage{titlesec}
\titleformat{\section}{\Large\bfseries\color{bentleyblue}}{\thesection}{1em}{}[\titlerule]
\titleformat{\subsection}{\large\bfseries\color{bentleyblue!80}}{\thesubsection}{1em}{}
\titleformat{\subsubsection}{\normalsize\bfseries\color{bentleyblue!60}}{\thesubsubsection}{1em}{}

% Lists
\usepackage{enumitem}
\setlist[itemize]{leftmargin=*, itemsep=2pt, parsep=0pt}
\setlist[enumerate]{leftmargin=*, itemsep=2pt, parsep=0pt}

% Tables
\usepackage{booktabs}
\usepackage{tabularx}

% Links
\usepackage{hyperref}
\hypersetup{
    colorlinks=true,
    linkcolor=linkblue,
    urlcolor=linkblue,
    pdfauthor={Onur Altındağ},
    pdftitle={EC111 Principles of Microeconomics - Fall 2026 Syllabus}
}

% Header/Footer
\usepackage{fancyhdr}
\pagestyle{fancy}
\fancyhf{}
\fancyhead[L]{\small\color{gray}\sectioncode{} Principles of Microeconomics}
\fancyhead[R]{\small\color{gray}Fall 2026}
\fancyfoot[C]{\thepage}
\renewcommand{\headrulewidth}{0.4pt}

% Remove paragraph indent, add space between paragraphs
\setlength{\parindent}{0pt}
\setlength{\parskip}{6pt}

\begin{document}

% Title
\begin{center}
{\LARGE\bfseries\color{bentleyblue} EC 111: Principles of Microeconomics}\\[0.3cm]
{\Large Fall 2026 Syllabus}\\[0.5cm]
{\large Bentley University}
\end{center}

\vspace{0.5cm}

%-----------------------------------------------------------
\section*{Course Description}
%-----------------------------------------------------------

This course introduces the fundamental principles and tools of microeconomics: how individuals and firms make decisions when resources are scarce, and how markets coordinate those decisions. We study demand, supply, and market equilibrium; elasticity; the effects of government policies such as taxes and price controls; the efficiency of markets and the sources of market failure; the costs of production; and firm behavior in different market structures, from perfect competition to monopoly, monopolistic competition, and oligopoly.

The emphasis throughout is on economic reasoning: thinking at the margin, recognizing trade-offs, and using simple graphical and numerical models to make sense of real-world questions in business and policy.

\subsection*{Knowledge}
\begin{itemize}
    \item Scarcity, opportunity cost, trade-offs, and the gains from trade.
    \item Demand, supply, market equilibrium, and how markets respond to shocks.
    \item Elasticity and its role in pricing, revenue, and tax incidence.
    \item Consumer surplus, producer surplus, market efficiency, and the deadweight loss of taxes and price controls.
    \item Externalities, public goods, and the case for government intervention.
    \item Costs of production and firm behavior under perfect competition, monopoly, monopolistic competition, and oligopoly.
\end{itemize}

\subsection*{Skills}
\begin{itemize}
    \item Use supply and demand diagrams to predict the effects of changes in market conditions and policy.
    \item Compute and interpret elasticities, surpluses, costs, and profit-maximizing output and price.
    \item Read and evaluate economic arguments critically and communicate them clearly.
\end{itemize}

\subsection*{Perspectives}
\begin{itemize}
    \item Understand marginal decision making as the core of economic reasoning.
    \item Understand when markets work well, when they fail, and what policy can and cannot fix.
\end{itemize}

%-----------------------------------------------------------
\section*{Class Information}
%-----------------------------------------------------------

\begin{tabularx}{\textwidth}{@{}lX@{}}
\textbf{Instructor:} & Onur Altındağ \\
\textbf{Office:} & RAU 019F \\
\textbf{Email:} & \href{mailto:oaltindag@bentley.edu}{oaltindag@bentley.edu} \\
\textbf{Web:} & \url{https://www.onuraltindag.info/} \\
\textbf{Course page:} & \url{https://www.onuraltindag.info/teaching/ec111/} \\
\textbf{Password:} & \texttt{GoFalcons!} \\[0.3cm]
\textbf{Course:} & \sectionname \\
\textbf{Semester:} & Fall 2026 (August 31, 2026 – December 15, 2026) \\
\textbf{Meeting Time:} & \meetingtime \\
\textbf{Location:} & \classroom \\
\textbf{Delivery Mode:} & In-Person Lecture \\
\end{tabularx}

\vspace{0.3cm}

Office hours are on Mondays and Wednesdays, 11:00 AM -- 12:00 PM and 1:00 -- 3:00 PM, by appointment. Please \href{https://outlook.office.com/bookwithme/user/1c4e4f54e0f0493397e615b2c80a59c7@bentley.edu/meetingtype/rwPThvM0MkK5oGlK_t37nw2?bookingcode=17c9a7b2-69fa-4c48-85ea-ad404642334b&anonymous&ismsaljsauthenabled&ep=mlink}{book a 30-minute slot} to schedule a meeting. If you need to meet urgently or no slots are available, email me directly.

Announcements come by email from me to your Bentley address; check it regularly. Slides, lecture notes, handouts, and mock exams are posted on the course page above after each class. When you email me, put \textbf{EC111} in the subject line.

%-----------------------------------------------------------
\section*{Important Dates and Evaluation}
%-----------------------------------------------------------

\begin{tabularx}{\textwidth}{@{}Xr@{}}
\toprule
\textbf{Component} & \textbf{Weight} \\
\midrule
Weekly assignments + participation & 20\% \\
\textbf{Wednesday, October 7, 2026} – First Midterm & 25\% \\
\textbf{Wednesday, November 11, 2026} – Second Midterm & 25\% \\
\textbf{\finaltime} – Final Exam & 30\% \\
\bottomrule
\end{tabularx}

%-----------------------------------------------------------
\section*{Grading}
%-----------------------------------------------------------

You \textbf{MUST} attend all the midterms and the final as there will be no make-up exams in this course. The midterms are \textbf{NOT} cumulative. If you miss or are likely to miss a midterm due to an emergency, please contact me as soon as possible. You will need to provide supporting documentation/verification of your absence. I will re-weight your final exam accordingly if you have a valid excuse. Please note that family vacations or missing the school shuttle are not valid excuses. If a serious illness, a death in the family, or another difficult circumstance keeps you from an exam or assignment, contact the Office of Student Success (JEN 336, 781.891.2803); they will notify all of your instructors on your behalf.

The final exam \textbf{is} cumulative. If you miss the final exam due to an emergency, you will receive an incomplete for this course. Do not take this class if you know that you will not be able to attend the final exam.

Exams are closed book and closed notes. A basic (non-graphing) calculator is allowed; phones, graphing calculators, and any other device are not, and calculators may not be shared.

You will have \textbf{10 homework assignments} on MindTap, roughly one per week (none in exam weeks), due \textbf{Sunday 11:00 PM}. Your \textbf{two lowest scores are dropped}; the remaining 8 count equally. Each question allows three attempts; an assignment scoring at least \textbf{50\%} receives full credit, otherwise zero. The submission system is automated and will not accept anything after the deadline (even 5 minutes), and there are no exceptions. The questions are algorithmic and slightly different for everyone, so you are encouraged to work together.

\subsection*{Letter grades}
\begin{center}
\begin{tabular}{@{}llllll@{}}
\toprule
A & 94--100 & B+ & 87--89 & C+ & 77--79 \\
A$-$ & 90--93 & B & 83--86 & C & 73--76 \\
 & & B$-$ & 80--82 & C$-$ & 70--72 \\
\midrule
D+ & 67--69 & D & 63--66 & D$-$ & 60--62 \\
F & below 60 & & & & \\
\bottomrule
\end{tabular}
\end{center}
I may lower these cutoffs at the end of the semester; I will not raise them. I do not respond to requests for individual extra credit or for grade adjustments after the fact.

%-----------------------------------------------------------
\section*{Textbook and MindTap}
%-----------------------------------------------------------

The course uses \textbf{MindTap for Mankiw, \textit{Principles of Microeconomics}} (Cengage). The e-book, your homework assignments, and practice questions all live there; a hard copy of the textbook is optional. You \textbf{MUST} register, using your Bentley email, through the link for your section:

\begin{center}
\textbf{\mindtapsection}\\
\url{\mindtaplink}
\end{center}

You have free temporary access until \mindtapaccess; after that you will be prompted to purchase (if your course materials are included in your tuition, no purchase is needed). Registration help: \url{https://startstrong.cengage.com}; technical support: \url{https://support.cengage.com}. Register in the first week; MindTap problems are not an excuse for a missed assignment.

\ifdefined\honors
%-----------------------------------------------------------
\section*{Honors Reading List}
%-----------------------------------------------------------

In addition to the textbook, honors students read and discuss a short set of articles and book chapters over the semester. The list and discussion dates will be posted on the course page in the first weeks of the semester.

\fi
%-----------------------------------------------------------
\section*{Academic Integrity}
%-----------------------------------------------------------

All students are expected to adhere to Bentley's Academic Integrity policy, which includes Bentley's Honor Code (details are in the Undergraduate Student Handbook, the \href{https://catalog.bentley.edu/undergraduate/academic-policies-procedures/}{academic policies section of the catalog}, and the academic integrity course page on Brightspace in which all students and faculty are enrolled). The essence of the policy is that you should not represent someone else's work as your own: no plagiarism, no cheating on exams, no illicit collaboration. Failure to adhere to the policy has serious consequences, including course failure, suspension, or expulsion from the university. The best way to avoid a problem is to consult with me before taking an action that might constitute a violation.

\subsection*{Generative AI}
Nothing in this course is graded on writing, and the homework is algorithmic practice, so the line is simple. You may use generative AI tools (ChatGPT, Claude, Gemini, Copilot, and the like) to study: to explain a concept, to check your reasoning on a practice problem, or to work through a homework question you are stuck on. Treat them as a tutor, not a substitute; the exams are closed book and closed device, and the only way to pass them is to be able to do the work yourself. Any use of AI or any other device during an exam is a violation of the academic integrity policy.

%-----------------------------------------------------------
\section*{Inclusion}
%-----------------------------------------------------------

Everyone in this class has different life experiences and perspectives, and all are valid. As the instructor, I will behave maturely and respectfully in all our class-related engagements, and I expect the same from you. We all adhere to Bentley's standards of conduct, the \href{https://www.bentley.edu/about/mission-and-values}{Bentley Core Values}, in class and in all course work. If you feel that I or anyone in this class has acted outside those values, come and talk to me. If you would rather not, speak with someone in the Office of Student Success (JEN 336, 781.891.2803). My roster shows your name as it appears in Workday; tell me early in the semester if you go by a different name or pronoun.

%-----------------------------------------------------------
\section*{Support and Resources}
%-----------------------------------------------------------

\subsection*{Student Accessibility Services}
Bentley University abides by Section 504 of the Rehabilitation Act of 1973 and the Americans with Disabilities Act of 1990, which stipulate that no student shall be denied the benefits of an education solely by reason of a disability. If you have a hidden or visible disability that may require classroom accommodations, contact \href{https://www.bentley.edu/offices/student-accessibility-services}{Student Accessibility Services} (Office of Student Success, JEN 336, 781.891.2004) within the first four weeks of the semester. Send me your accommodation plan as soon as you have it; accommodations are arranged at least one week before each exam.

\subsection*{The Howard A. Winer '58 Lab for Economics, Accounting and Finance (LEAF)}
The LEAF provides a welcoming and inclusive learning environment where students are encouraged to seek academic support for their Economics, Accounting and Finance courses. Students utilizing the LEAF will find knowledgeable peer tutors ready to help their colleagues to thrive in the Bentley business curriculum. LEAF's free, drop-in tutoring takes place in-person at Morison 220D (rear entrance), opening September 14. For additional information visit \url{https://www.bentley.edu/centers/leaf}. LEAF is a supplement to, not a substitute for, office hours.

\subsection*{Bentley policies and procedures}
All courses at Bentley reflect the university's commitment to a set of core values and practices. You are expected to be familiar with:
\begin{itemize}
    \item \href{https://catalog.bentley.edu/undergraduate/academic-policies-procedures/}{Bentley's honor code and academic integrity system}
    \item \href{https://www.bentley.edu/equity-reporting}{Equity and bias reporting forms and procedures} (Bias Incident Response Team, Office of Institutional Equity)
    \item \href{https://www.bentley.edu/offices/student-accessibility-services}{Student disability accommodations}
    \item \href{https://catalog.bentley.edu/undergraduate/academic-policies-procedures/}{Religious observances policy}
    \item \href{https://www.bentley.edu/about/mission-and-values}{Bentley's core values}
\end{itemize}

\subsection*{Course materials}
Slides, notes, handouts, and exams for this course are for your own use during this semester. They are protected by copyright; do not post or distribute them outside the class without my written consent. The course page is password-protected; do not share the password outside the class.

%-----------------------------------------------------------
\section*{Attendance Policy}
%-----------------------------------------------------------

All students must attend the in-person classes. In-person attendance is essential because it allows for real-time interaction, immediate feedback, and the kind of dialogue that makes learning stick. I do not record lectures, so if you miss a class, you miss the material.

Remote attendance is not an option. This course follows Bentley's \href{https://catalog.bentley.edu/undergraduate/academic-policies-procedures/}{Academic Engagement/Attendance Policy}: students who do not attend or engage during the first week are dropped as no-shows, and two or more consecutive weeks of non-participation can lead to administrative withdrawal. If you must miss class, you are responsible for the material; read the chapter and the posted notes, then bring specific questions to office hours.

\textbf{Weather and closures.} If the university cancels classes, I will email you whether there is an alternative assignment or arrangement for that day.

\textbf{Electronic devices.} Phones away and laptops and tablets closed for the whole class. The work in this course is drawing and reading diagrams, which you learn by doing them by hand, and a screen in your line of sight costs the people around you as much as it costs you. Using a device during class counts against your participation grade. Exceptions are by my permission only; if you have one, you never need to explain it to anyone in the room. Photographing the board at the end of class is fine.

\textbf{Religious observances.} Bentley provides reasonable academic adjustments when a sincerely held religious observance conflicts with a class, exam, or deadline (see the \href{https://downloads.bentley.edu/general/Religious\%20Observances\%202026-2027.pdf}{religious observances calendar} and the \href{https://catalog.bentley.edu/undergraduate/academic-policies-procedures/}{policy}). If you anticipate a conflict, tell me in advance, ideally in the first two weeks of the semester, so we can agree on an adjustment.

%-----------------------------------------------------------
\section*{Tentative Schedule}
%-----------------------------------------------------------

\begin{itemize}
    \item \textbf{Week 1 (Aug 31, Sep 2):} Ten Principles of Economics – Ch. 1 \& Thinking Like an Economist – Ch. 2
    \item \textbf{Week 2 (Sep 9):} Interdependence and the Gains from Trade – Ch. 3 \textit{(no class Sep 7, Labor Day)}
    \item \textbf{Week 3 (Sep 14, 16):} The Market Forces of Supply and Demand – Ch. 4
    \item \textbf{Week 4 (Sep 21, 23):} Elasticity and Its Application – Ch. 5
    \item \textbf{Week 5 (Sep 28, 30):} Supply, Demand, and Government Policies – Ch. 6
    \item \textbf{Week 6 (Oct 5, 7):} Review / \textbf{First Midterm (Oct 7)}
    \item \textbf{Week 7 (Oct 14):} Consumers, Producers, and the Efficiency of Markets – Ch. 7 \textit{(no class Oct 12, Fall Break)}
    \item \textbf{Week 8 (Oct 19, 21):} The Costs of Taxation – Ch. 8 \& International Trade – Ch. 9
    \item \textbf{Week 9 (Oct 26, 28):} Externalities – Ch. 10 \& Public Goods and Common Resources – Ch. 11
    \item \textbf{Week 10 (Nov 2, 4):} The Costs of Production – Ch. 14
    \item \textbf{Week 11 (Nov 9, 11):} Review / \textbf{Second Midterm (Nov 11)}
    \item \textbf{Week 12 (Nov 16, 18):} Firms in Competitive Markets – Ch. 15
    \item \textbf{Week 13 (Nov 23):} Monopoly – Ch. 16 \textit{(no class Nov 25, Thanksgiving Break)}
    \item \textbf{Week 14 (Nov 30, Dec 2):} Monopolistic Competition – Ch. 17 \& Oligopoly – Ch. 18
    \item \textbf{Week 15 (Dec 7):} Review
    \item \textbf{Final Exam:} \finaltime
\end{itemize}

\textit{Note: Chapter references are from Mankiw's \textit{Principles of Microeconomics}, 10th edition (the MindTap e-book). The schedule, and if necessary other parts of this syllabus, may change during the semester; changes are announced in class and by email.}

\end{document}
"
\ifdefined\honors
  \newcommand{\sectioncode}{EC 111-10-H}
  \newcommand{\sectionname}{EC 111-10-H Principles of Microeconomics (Honors)}
  \newcommand{\meetingtime}{Monday/Wednesday, 3:30 PM -- 4:50 PM}
  \newcommand{\classroom}{Smith Technology Center 103}
  \newcommand{\finaltime}{Friday, December 11, 2026, 6:00 PM -- 8:00 PM}
  \newcommand{\mindtapsection}{Section 10-H of EC 111 Principles of Microeconomics Fall 2026}
  \newcommand{\mindtaplink}{https://student.cengage.com/course-link/MTPNS57GW6XN}
  \newcommand{\mindtapaccess}{the end of Sunday, September 13, 2026}
\else
  \newcommand{\sectioncode}{EC 111-20}
  \newcommand{\sectionname}{EC 111-20 Principles of Microeconomics}
  \newcommand{\meetingtime}{Monday/Wednesday, 9:30 AM -- 10:50 AM}
  \newcommand{\classroom}{Jennison Hall 412}
  \newcommand{\finaltime}{Friday, December 11, 2026, 8:00 AM -- 10:00 AM}
  \newcommand{\mindtapsection}{Section 20 of EC 111 Principles of Microeconomics Fall 2026}
  \newcommand{\mindtaplink}{https://student.cengage.com/course-link/MTPP8N7GWX7L}
  \newcommand{\mindtapaccess}{the end of Sunday, September 13, 2026}
\fi

% XeLaTeX for Turkish character support
\usepackage{fontspec}
\setmainfont{DejaVuSerif}[Extension=.ttf, BoldFont=DejaVuSerif-Bold, ItalicFont=DejaVuSerif-Italic, BoldItalicFont=DejaVuSerif-BoldItalic]
\setsansfont{DejaVuSans}[Extension=.ttf, BoldFont=DejaVuSans-Bold, ItalicFont=DejaVuSans-Oblique, BoldItalicFont=DejaVuSans-BoldOblique]

% Page layout
\usepackage[margin=1in]{geometry}
\usepackage{setspace}
\onehalfspacing

% Colors and styling
\usepackage[dvipsnames]{xcolor}
\definecolor{bentleyblue}{RGB}{0, 51, 102}
\definecolor{linkblue}{RGB}{0, 102, 204}

% Headers and sections
\usepackage{titlesec}
\titleformat{\section}{\Large\bfseries\color{bentleyblue}}{\thesection}{1em}{}[\titlerule]
\titleformat{\subsection}{\large\bfseries\color{bentleyblue!80}}{\thesubsection}{1em}{}
\titleformat{\subsubsection}{\normalsize\bfseries\color{bentleyblue!60}}{\thesubsubsection}{1em}{}

% Lists
\usepackage{enumitem}
\setlist[itemize]{leftmargin=*, itemsep=2pt, parsep=0pt}
\setlist[enumerate]{leftmargin=*, itemsep=2pt, parsep=0pt}

% Tables
\usepackage{booktabs}
\usepackage{tabularx}

% Links
\usepackage{hyperref}
\hypersetup{
    colorlinks=true,
    linkcolor=linkblue,
    urlcolor=linkblue,
    pdfauthor={Onur Altındağ},
    pdftitle={EC111 Principles of Microeconomics - Fall 2026 Syllabus}
}

% Header/Footer
\usepackage{fancyhdr}
\pagestyle{fancy}
\fancyhf{}
\fancyhead[L]{\small\color{gray}\sectioncode{} Principles of Microeconomics}
\fancyhead[R]{\small\color{gray}Fall 2026}
\fancyfoot[C]{\thepage}
\renewcommand{\headrulewidth}{0.4pt}

% Remove paragraph indent, add space between paragraphs
\setlength{\parindent}{0pt}
\setlength{\parskip}{6pt}

\begin{document}

% Title
\begin{center}
{\LARGE\bfseries\color{bentleyblue} EC 111: Principles of Microeconomics}\\[0.3cm]
{\Large Fall 2026 Syllabus}\\[0.5cm]
{\large Bentley University}
\end{center}

\vspace{0.5cm}

%-----------------------------------------------------------
\section*{Course Description}
%-----------------------------------------------------------

This course introduces the fundamental principles and tools of microeconomics: how individuals and firms make decisions when resources are scarce, and how markets coordinate those decisions. We study demand, supply, and market equilibrium; elasticity; the effects of government policies such as taxes and price controls; the efficiency of markets and the sources of market failure; the costs of production; and firm behavior in different market structures, from perfect competition to monopoly, monopolistic competition, and oligopoly.

The emphasis throughout is on economic reasoning: thinking at the margin, recognizing trade-offs, and using simple graphical and numerical models to make sense of real-world questions in business and policy.

\subsection*{Knowledge}
\begin{itemize}
    \item Scarcity, opportunity cost, trade-offs, and the gains from trade.
    \item Demand, supply, market equilibrium, and how markets respond to shocks.
    \item Elasticity and its role in pricing, revenue, and tax incidence.
    \item Consumer surplus, producer surplus, market efficiency, and the deadweight loss of taxes and price controls.
    \item Externalities, public goods, and the case for government intervention.
    \item Costs of production and firm behavior under perfect competition, monopoly, monopolistic competition, and oligopoly.
\end{itemize}

\subsection*{Skills}
\begin{itemize}
    \item Use supply and demand diagrams to predict the effects of changes in market conditions and policy.
    \item Compute and interpret elasticities, surpluses, costs, and profit-maximizing output and price.
    \item Read and evaluate economic arguments critically and communicate them clearly.
\end{itemize}

\subsection*{Perspectives}
\begin{itemize}
    \item Understand marginal decision making as the core of economic reasoning.
    \item Understand when markets work well, when they fail, and what policy can and cannot fix.
\end{itemize}

%-----------------------------------------------------------
\section*{Class Information}
%-----------------------------------------------------------

\begin{tabularx}{\textwidth}{@{}lX@{}}
\textbf{Instructor:} & Onur Altındağ \\
\textbf{Office:} & RAU 019F \\
\textbf{Email:} & \href{mailto:oaltindag@bentley.edu}{oaltindag@bentley.edu} \\
\textbf{Web:} & \url{https://www.onuraltindag.info/} \\
\textbf{Course page:} & \url{https://www.onuraltindag.info/teaching/ec111/} \\
\textbf{Password:} & \texttt{GoFalcons!} \\[0.3cm]
\textbf{Course:} & \sectionname \\
\textbf{Semester:} & Fall 2026 (August 31, 2026 – December 15, 2026) \\
\textbf{Meeting Time:} & \meetingtime \\
\textbf{Location:} & \classroom \\
\textbf{Delivery Mode:} & In-Person Lecture \\
\end{tabularx}

\vspace{0.3cm}

Office hours are on Mondays and Wednesdays, 11:00 AM -- 12:00 PM and 1:00 -- 3:00 PM, by appointment. Please \href{https://outlook.office.com/bookwithme/user/1c4e4f54e0f0493397e615b2c80a59c7@bentley.edu/meetingtype/rwPThvM0MkK5oGlK_t37nw2?bookingcode=17c9a7b2-69fa-4c48-85ea-ad404642334b&anonymous&ismsaljsauthenabled&ep=mlink}{book a 30-minute slot} to schedule a meeting. If you need to meet urgently or no slots are available, email me directly.

Announcements come by email from me to your Bentley address; check it regularly. Slides, lecture notes, handouts, and mock exams are posted on the course page above after each class. When you email me, put \textbf{EC111} in the subject line.

%-----------------------------------------------------------
\section*{Important Dates and Evaluation}
%-----------------------------------------------------------

\begin{tabularx}{\textwidth}{@{}Xr@{}}
\toprule
\textbf{Component} & \textbf{Weight} \\
\midrule
Weekly assignments + participation & 20\% \\
\textbf{Wednesday, October 7, 2026} – First Midterm & 25\% \\
\textbf{Wednesday, November 11, 2026} – Second Midterm & 25\% \\
\textbf{\finaltime} – Final Exam & 30\% \\
\bottomrule
\end{tabularx}

%-----------------------------------------------------------
\section*{Grading}
%-----------------------------------------------------------

You \textbf{MUST} attend all the midterms and the final as there will be no make-up exams in this course. The midterms are \textbf{NOT} cumulative. If you miss or are likely to miss a midterm due to an emergency, please contact me as soon as possible. You will need to provide supporting documentation/verification of your absence. I will re-weight your final exam accordingly if you have a valid excuse. Please note that family vacations or missing the school shuttle are not valid excuses. If a serious illness, a death in the family, or another difficult circumstance keeps you from an exam or assignment, contact the Office of Student Success (JEN 336, 781.891.2803); they will notify all of your instructors on your behalf.

The final exam \textbf{is} cumulative. If you miss the final exam due to an emergency, you will receive an incomplete for this course. Do not take this class if you know that you will not be able to attend the final exam.

Exams are closed book and closed notes. A basic (non-graphing) calculator is allowed; phones, graphing calculators, and any other device are not, and calculators may not be shared.

You will have \textbf{10 homework assignments} on MindTap, roughly one per week (none in exam weeks), due \textbf{Sunday 11:00 PM}. Your \textbf{two lowest scores are dropped}; the remaining 8 count equally. Each question allows three attempts; an assignment scoring at least \textbf{50\%} receives full credit, otherwise zero. The submission system is automated and will not accept anything after the deadline (even 5 minutes), and there are no exceptions. The questions are algorithmic and slightly different for everyone, so you are encouraged to work together.

\subsection*{Letter grades}
\begin{center}
\begin{tabular}{@{}llllll@{}}
\toprule
A & 94--100 & B+ & 87--89 & C+ & 77--79 \\
A$-$ & 90--93 & B & 83--86 & C & 73--76 \\
 & & B$-$ & 80--82 & C$-$ & 70--72 \\
\midrule
D+ & 67--69 & D & 63--66 & D$-$ & 60--62 \\
F & below 60 & & & & \\
\bottomrule
\end{tabular}
\end{center}
I may lower these cutoffs at the end of the semester; I will not raise them. I do not respond to requests for individual extra credit or for grade adjustments after the fact.

%-----------------------------------------------------------
\section*{Textbook and MindTap}
%-----------------------------------------------------------

The course uses \textbf{MindTap for Mankiw, \textit{Principles of Microeconomics}} (Cengage). The e-book, your homework assignments, and practice questions all live there; a hard copy of the textbook is optional. You \textbf{MUST} register, using your Bentley email, through the link for your section:

\begin{center}
\textbf{\mindtapsection}\\
\url{\mindtaplink}
\end{center}

You have free temporary access until \mindtapaccess; after that you will be prompted to purchase (if your course materials are included in your tuition, no purchase is needed). Registration help: \url{https://startstrong.cengage.com}; technical support: \url{https://support.cengage.com}. Register in the first week; MindTap problems are not an excuse for a missed assignment.

\ifdefined\honors
%-----------------------------------------------------------
\section*{Honors Reading List}
%-----------------------------------------------------------

In addition to the textbook, honors students read and discuss a short set of articles and book chapters over the semester. The list and discussion dates will be posted on the course page in the first weeks of the semester.

\fi
%-----------------------------------------------------------
\section*{Academic Integrity}
%-----------------------------------------------------------

All students are expected to adhere to Bentley's Academic Integrity policy, which includes Bentley's Honor Code (details are in the Undergraduate Student Handbook, the \href{https://catalog.bentley.edu/undergraduate/academic-policies-procedures/}{academic policies section of the catalog}, and the academic integrity course page on Brightspace in which all students and faculty are enrolled). The essence of the policy is that you should not represent someone else's work as your own: no plagiarism, no cheating on exams, no illicit collaboration. Failure to adhere to the policy has serious consequences, including course failure, suspension, or expulsion from the university. The best way to avoid a problem is to consult with me before taking an action that might constitute a violation.

\subsection*{Generative AI}
Nothing in this course is graded on writing, and the homework is algorithmic practice, so the line is simple. You may use generative AI tools (ChatGPT, Claude, Gemini, Copilot, and the like) to study: to explain a concept, to check your reasoning on a practice problem, or to work through a homework question you are stuck on. Treat them as a tutor, not a substitute; the exams are closed book and closed device, and the only way to pass them is to be able to do the work yourself. Any use of AI or any other device during an exam is a violation of the academic integrity policy.

%-----------------------------------------------------------
\section*{Inclusion}
%-----------------------------------------------------------

Everyone in this class has different life experiences and perspectives, and all are valid. As the instructor, I will behave maturely and respectfully in all our class-related engagements, and I expect the same from you. We all adhere to Bentley's standards of conduct, the \href{https://www.bentley.edu/about/mission-and-values}{Bentley Core Values}, in class and in all course work. If you feel that I or anyone in this class has acted outside those values, come and talk to me. If you would rather not, speak with someone in the Office of Student Success (JEN 336, 781.891.2803). My roster shows your name as it appears in Workday; tell me early in the semester if you go by a different name or pronoun.

%-----------------------------------------------------------
\section*{Support and Resources}
%-----------------------------------------------------------

\subsection*{Student Accessibility Services}
Bentley University abides by Section 504 of the Rehabilitation Act of 1973 and the Americans with Disabilities Act of 1990, which stipulate that no student shall be denied the benefits of an education solely by reason of a disability. If you have a hidden or visible disability that may require classroom accommodations, contact \href{https://www.bentley.edu/offices/student-accessibility-services}{Student Accessibility Services} (Office of Student Success, JEN 336, 781.891.2004) within the first four weeks of the semester. Send me your accommodation plan as soon as you have it; accommodations are arranged at least one week before each exam.

\subsection*{The Howard A. Winer '58 Lab for Economics, Accounting and Finance (LEAF)}
The LEAF provides a welcoming and inclusive learning environment where students are encouraged to seek academic support for their Economics, Accounting and Finance courses. Students utilizing the LEAF will find knowledgeable peer tutors ready to help their colleagues to thrive in the Bentley business curriculum. LEAF's free, drop-in tutoring takes place in-person at Morison 220D (rear entrance), opening September 14. For additional information visit \url{https://www.bentley.edu/centers/leaf}. LEAF is a supplement to, not a substitute for, office hours.

\subsection*{Bentley policies and procedures}
All courses at Bentley reflect the university's commitment to a set of core values and practices. You are expected to be familiar with:
\begin{itemize}
    \item \href{https://catalog.bentley.edu/undergraduate/academic-policies-procedures/}{Bentley's honor code and academic integrity system}
    \item \href{https://www.bentley.edu/equity-reporting}{Equity and bias reporting forms and procedures} (Bias Incident Response Team, Office of Institutional Equity)
    \item \href{https://www.bentley.edu/offices/student-accessibility-services}{Student disability accommodations}
    \item \href{https://catalog.bentley.edu/undergraduate/academic-policies-procedures/}{Religious observances policy}
    \item \href{https://www.bentley.edu/about/mission-and-values}{Bentley's core values}
\end{itemize}

\subsection*{Course materials}
Slides, notes, handouts, and exams for this course are for your own use during this semester. They are protected by copyright; do not post or distribute them outside the class without my written consent. The course page is password-protected; do not share the password outside the class.

%-----------------------------------------------------------
\section*{Attendance Policy}
%-----------------------------------------------------------

All students must attend the in-person classes. In-person attendance is essential because it allows for real-time interaction, immediate feedback, and the kind of dialogue that makes learning stick. I do not record lectures, so if you miss a class, you miss the material.

Remote attendance is not an option. This course follows Bentley's \href{https://catalog.bentley.edu/undergraduate/academic-policies-procedures/}{Academic Engagement/Attendance Policy}: students who do not attend or engage during the first week are dropped as no-shows, and two or more consecutive weeks of non-participation can lead to administrative withdrawal. If you must miss class, you are responsible for the material; read the chapter and the posted notes, then bring specific questions to office hours.

\textbf{Weather and closures.} If the university cancels classes, I will email you whether there is an alternative assignment or arrangement for that day.

\textbf{Electronic devices.} Phones away and laptops and tablets closed for the whole class. The work in this course is drawing and reading diagrams, which you learn by doing them by hand, and a screen in your line of sight costs the people around you as much as it costs you. Using a device during class counts against your participation grade. Exceptions are by my permission only; if you have one, you never need to explain it to anyone in the room. Photographing the board at the end of class is fine.

\textbf{Religious observances.} Bentley provides reasonable academic adjustments when a sincerely held religious observance conflicts with a class, exam, or deadline (see the \href{https://downloads.bentley.edu/general/Religious\%20Observances\%202026-2027.pdf}{religious observances calendar} and the \href{https://catalog.bentley.edu/undergraduate/academic-policies-procedures/}{policy}). If you anticipate a conflict, tell me in advance, ideally in the first two weeks of the semester, so we can agree on an adjustment.

%-----------------------------------------------------------
\section*{Tentative Schedule}
%-----------------------------------------------------------

\begin{itemize}
    \item \textbf{Week 1 (Aug 31, Sep 2):} Ten Principles of Economics – Ch. 1 \& Thinking Like an Economist – Ch. 2
    \item \textbf{Week 2 (Sep 9):} Interdependence and the Gains from Trade – Ch. 3 \textit{(no class Sep 7, Labor Day)}
    \item \textbf{Week 3 (Sep 14, 16):} The Market Forces of Supply and Demand – Ch. 4
    \item \textbf{Week 4 (Sep 21, 23):} Elasticity and Its Application – Ch. 5
    \item \textbf{Week 5 (Sep 28, 30):} Supply, Demand, and Government Policies – Ch. 6
    \item \textbf{Week 6 (Oct 5, 7):} Review / \textbf{First Midterm (Oct 7)}
    \item \textbf{Week 7 (Oct 14):} Consumers, Producers, and the Efficiency of Markets – Ch. 7 \textit{(no class Oct 12, Fall Break)}
    \item \textbf{Week 8 (Oct 19, 21):} The Costs of Taxation – Ch. 8 \& International Trade – Ch. 9
    \item \textbf{Week 9 (Oct 26, 28):} Externalities – Ch. 10 \& Public Goods and Common Resources – Ch. 11
    \item \textbf{Week 10 (Nov 2, 4):} The Costs of Production – Ch. 14
    \item \textbf{Week 11 (Nov 9, 11):} Review / \textbf{Second Midterm (Nov 11)}
    \item \textbf{Week 12 (Nov 16, 18):} Firms in Competitive Markets – Ch. 15
    \item \textbf{Week 13 (Nov 23):} Monopoly – Ch. 16 \textit{(no class Nov 25, Thanksgiving Break)}
    \item \textbf{Week 14 (Nov 30, Dec 2):} Monopolistic Competition – Ch. 17 \& Oligopoly – Ch. 18
    \item \textbf{Week 15 (Dec 7):} Review
    \item \textbf{Final Exam:} \finaltime
\end{itemize}

\textit{Note: Chapter references are from Mankiw's \textit{Principles of Microeconomics}, 10th edition (the MindTap e-book). The schedule, and if necessary other parts of this syllabus, may change during the semester; changes are announced in class and by email.}

\end{document}
"
%   xelatex -jobname=EC111_FA2026_Syllabus_20  "\documentclass[11pt,letterpaper]{article}

% Compile with XeLaTeX. Two sections from one source:
%   xelatex -jobname=EC111_FA2026_Syllabus_10H "\def\honors{1}\documentclass[11pt,letterpaper]{article}

% Compile with XeLaTeX. Two sections from one source:
%   xelatex -jobname=EC111_FA2026_Syllabus_10H "\def\honors{1}\input{EC111_FA2026_Syllabus.tex}"
%   xelatex -jobname=EC111_FA2026_Syllabus_20  "\input{EC111_FA2026_Syllabus.tex}"
\ifdefined\honors
  \newcommand{\sectioncode}{EC 111-10-H}
  \newcommand{\sectionname}{EC 111-10-H Principles of Microeconomics (Honors)}
  \newcommand{\meetingtime}{Monday/Wednesday, 3:30 PM -- 4:50 PM}
  \newcommand{\classroom}{Smith Technology Center 103}
  \newcommand{\finaltime}{Friday, December 11, 2026, 6:00 PM -- 8:00 PM}
  \newcommand{\mindtapsection}{Section 10-H of EC 111 Principles of Microeconomics Fall 2026}
  \newcommand{\mindtaplink}{https://student.cengage.com/course-link/MTPNS57GW6XN}
  \newcommand{\mindtapaccess}{the end of Sunday, September 13, 2026}
\else
  \newcommand{\sectioncode}{EC 111-20}
  \newcommand{\sectionname}{EC 111-20 Principles of Microeconomics}
  \newcommand{\meetingtime}{Monday/Wednesday, 9:30 AM -- 10:50 AM}
  \newcommand{\classroom}{Jennison Hall 412}
  \newcommand{\finaltime}{Friday, December 11, 2026, 8:00 AM -- 10:00 AM}
  \newcommand{\mindtapsection}{Section 20 of EC 111 Principles of Microeconomics Fall 2026}
  \newcommand{\mindtaplink}{https://student.cengage.com/course-link/MTPP8N7GWX7L}
  \newcommand{\mindtapaccess}{the end of Sunday, September 13, 2026}
\fi

% XeLaTeX for Turkish character support
\usepackage{fontspec}
\setmainfont{DejaVuSerif}[Extension=.ttf, BoldFont=DejaVuSerif-Bold, ItalicFont=DejaVuSerif-Italic, BoldItalicFont=DejaVuSerif-BoldItalic]
\setsansfont{DejaVuSans}[Extension=.ttf, BoldFont=DejaVuSans-Bold, ItalicFont=DejaVuSans-Oblique, BoldItalicFont=DejaVuSans-BoldOblique]

% Page layout
\usepackage[margin=1in]{geometry}
\usepackage{setspace}
\onehalfspacing

% Colors and styling
\usepackage[dvipsnames]{xcolor}
\definecolor{bentleyblue}{RGB}{0, 51, 102}
\definecolor{linkblue}{RGB}{0, 102, 204}

% Headers and sections
\usepackage{titlesec}
\titleformat{\section}{\Large\bfseries\color{bentleyblue}}{\thesection}{1em}{}[\titlerule]
\titleformat{\subsection}{\large\bfseries\color{bentleyblue!80}}{\thesubsection}{1em}{}
\titleformat{\subsubsection}{\normalsize\bfseries\color{bentleyblue!60}}{\thesubsubsection}{1em}{}

% Lists
\usepackage{enumitem}
\setlist[itemize]{leftmargin=*, itemsep=2pt, parsep=0pt}
\setlist[enumerate]{leftmargin=*, itemsep=2pt, parsep=0pt}

% Tables
\usepackage{booktabs}
\usepackage{tabularx}

% Links
\usepackage{hyperref}
\hypersetup{
    colorlinks=true,
    linkcolor=linkblue,
    urlcolor=linkblue,
    pdfauthor={Onur Altındağ},
    pdftitle={EC111 Principles of Microeconomics - Fall 2026 Syllabus}
}

% Header/Footer
\usepackage{fancyhdr}
\pagestyle{fancy}
\fancyhf{}
\fancyhead[L]{\small\color{gray}\sectioncode{} Principles of Microeconomics}
\fancyhead[R]{\small\color{gray}Fall 2026}
\fancyfoot[C]{\thepage}
\renewcommand{\headrulewidth}{0.4pt}

% Remove paragraph indent, add space between paragraphs
\setlength{\parindent}{0pt}
\setlength{\parskip}{6pt}

\begin{document}

% Title
\begin{center}
{\LARGE\bfseries\color{bentleyblue} EC 111: Principles of Microeconomics}\\[0.3cm]
{\Large Fall 2026 Syllabus}\\[0.5cm]
{\large Bentley University}
\end{center}

\vspace{0.5cm}

%-----------------------------------------------------------
\section*{Course Description}
%-----------------------------------------------------------

This course introduces the fundamental principles and tools of microeconomics: how individuals and firms make decisions when resources are scarce, and how markets coordinate those decisions. We study demand, supply, and market equilibrium; elasticity; the effects of government policies such as taxes and price controls; the efficiency of markets and the sources of market failure; the costs of production; and firm behavior in different market structures, from perfect competition to monopoly, monopolistic competition, and oligopoly.

The emphasis throughout is on economic reasoning: thinking at the margin, recognizing trade-offs, and using simple graphical and numerical models to make sense of real-world questions in business and policy.

\subsection*{Knowledge}
\begin{itemize}
    \item Scarcity, opportunity cost, trade-offs, and the gains from trade.
    \item Demand, supply, market equilibrium, and how markets respond to shocks.
    \item Elasticity and its role in pricing, revenue, and tax incidence.
    \item Consumer surplus, producer surplus, market efficiency, and the deadweight loss of taxes and price controls.
    \item Externalities, public goods, and the case for government intervention.
    \item Costs of production and firm behavior under perfect competition, monopoly, monopolistic competition, and oligopoly.
\end{itemize}

\subsection*{Skills}
\begin{itemize}
    \item Use supply and demand diagrams to predict the effects of changes in market conditions and policy.
    \item Compute and interpret elasticities, surpluses, costs, and profit-maximizing output and price.
    \item Read and evaluate economic arguments critically and communicate them clearly.
\end{itemize}

\subsection*{Perspectives}
\begin{itemize}
    \item Understand marginal decision making as the core of economic reasoning.
    \item Understand when markets work well, when they fail, and what policy can and cannot fix.
\end{itemize}

%-----------------------------------------------------------
\section*{Class Information}
%-----------------------------------------------------------

\begin{tabularx}{\textwidth}{@{}lX@{}}
\textbf{Instructor:} & Onur Altındağ \\
\textbf{Office:} & RAU 019F \\
\textbf{Email:} & \href{mailto:oaltindag@bentley.edu}{oaltindag@bentley.edu} \\
\textbf{Web:} & \url{https://www.onuraltindag.info/} \\
\textbf{Course page:} & \url{https://www.onuraltindag.info/teaching/ec111/} \\
\textbf{Password:} & \texttt{GoFalcons!} \\[0.3cm]
\textbf{Course:} & \sectionname \\
\textbf{Semester:} & Fall 2026 (August 31, 2026 – December 15, 2026) \\
\textbf{Meeting Time:} & \meetingtime \\
\textbf{Location:} & \classroom \\
\textbf{Delivery Mode:} & In-Person Lecture \\
\end{tabularx}

\vspace{0.3cm}

Office hours are on Mondays and Wednesdays, 11:00 AM -- 12:00 PM and 1:00 -- 3:00 PM, by appointment. Please \href{https://outlook.office.com/bookwithme/user/1c4e4f54e0f0493397e615b2c80a59c7@bentley.edu/meetingtype/rwPThvM0MkK5oGlK_t37nw2?bookingcode=17c9a7b2-69fa-4c48-85ea-ad404642334b&anonymous&ismsaljsauthenabled&ep=mlink}{book a 30-minute slot} to schedule a meeting. If you need to meet urgently or no slots are available, email me directly.

Announcements come by email from me to your Bentley address; check it regularly. Slides, lecture notes, handouts, and mock exams are posted on the course page above after each class. When you email me, put \textbf{EC111} in the subject line.

%-----------------------------------------------------------
\section*{Important Dates and Evaluation}
%-----------------------------------------------------------

\begin{tabularx}{\textwidth}{@{}Xr@{}}
\toprule
\textbf{Component} & \textbf{Weight} \\
\midrule
Weekly assignments + participation & 20\% \\
\textbf{Wednesday, October 7, 2026} – First Midterm & 25\% \\
\textbf{Wednesday, November 11, 2026} – Second Midterm & 25\% \\
\textbf{\finaltime} – Final Exam & 30\% \\
\bottomrule
\end{tabularx}

%-----------------------------------------------------------
\section*{Grading}
%-----------------------------------------------------------

You \textbf{MUST} attend all the midterms and the final as there will be no make-up exams in this course. The midterms are \textbf{NOT} cumulative. If you miss or are likely to miss a midterm due to an emergency, please contact me as soon as possible. You will need to provide supporting documentation/verification of your absence. I will re-weight your final exam accordingly if you have a valid excuse. Please note that family vacations or missing the school shuttle are not valid excuses. If a serious illness, a death in the family, or another difficult circumstance keeps you from an exam or assignment, contact the Office of Student Success (JEN 336, 781.891.2803); they will notify all of your instructors on your behalf.

The final exam \textbf{is} cumulative. If you miss the final exam due to an emergency, you will receive an incomplete for this course. Do not take this class if you know that you will not be able to attend the final exam.

Exams are closed book and closed notes. A basic (non-graphing) calculator is allowed; phones, graphing calculators, and any other device are not, and calculators may not be shared.

You will have \textbf{10 homework assignments} on MindTap, roughly one per week (none in exam weeks), due \textbf{Sunday 11:00 PM}. Your \textbf{two lowest scores are dropped}; the remaining 8 count equally. Each question allows three attempts; an assignment scoring at least \textbf{50\%} receives full credit, otherwise zero. The submission system is automated and will not accept anything after the deadline (even 5 minutes), and there are no exceptions. The questions are algorithmic and slightly different for everyone, so you are encouraged to work together.

\subsection*{Letter grades}
\begin{center}
\begin{tabular}{@{}llllll@{}}
\toprule
A & 94--100 & B+ & 87--89 & C+ & 77--79 \\
A$-$ & 90--93 & B & 83--86 & C & 73--76 \\
 & & B$-$ & 80--82 & C$-$ & 70--72 \\
\midrule
D+ & 67--69 & D & 63--66 & D$-$ & 60--62 \\
F & below 60 & & & & \\
\bottomrule
\end{tabular}
\end{center}
I may lower these cutoffs at the end of the semester; I will not raise them. I do not respond to requests for individual extra credit or for grade adjustments after the fact.

%-----------------------------------------------------------
\section*{Textbook and MindTap}
%-----------------------------------------------------------

The course uses \textbf{MindTap for Mankiw, \textit{Principles of Microeconomics}} (Cengage). The e-book, your homework assignments, and practice questions all live there; a hard copy of the textbook is optional. You \textbf{MUST} register, using your Bentley email, through the link for your section:

\begin{center}
\textbf{\mindtapsection}\\
\url{\mindtaplink}
\end{center}

You have free temporary access until \mindtapaccess; after that you will be prompted to purchase (if your course materials are included in your tuition, no purchase is needed). Registration help: \url{https://startstrong.cengage.com}; technical support: \url{https://support.cengage.com}. Register in the first week; MindTap problems are not an excuse for a missed assignment.

\ifdefined\honors
%-----------------------------------------------------------
\section*{Honors Reading List}
%-----------------------------------------------------------

In addition to the textbook, honors students read and discuss a short set of articles and book chapters over the semester. The list and discussion dates will be posted on the course page in the first weeks of the semester.

\fi
%-----------------------------------------------------------
\section*{Academic Integrity}
%-----------------------------------------------------------

All students are expected to adhere to Bentley's Academic Integrity policy, which includes Bentley's Honor Code (details are in the Undergraduate Student Handbook, the \href{https://catalog.bentley.edu/undergraduate/academic-policies-procedures/}{academic policies section of the catalog}, and the academic integrity course page on Brightspace in which all students and faculty are enrolled). The essence of the policy is that you should not represent someone else's work as your own: no plagiarism, no cheating on exams, no illicit collaboration. Failure to adhere to the policy has serious consequences, including course failure, suspension, or expulsion from the university. The best way to avoid a problem is to consult with me before taking an action that might constitute a violation.

\subsection*{Generative AI}
Nothing in this course is graded on writing, and the homework is algorithmic practice, so the line is simple. You may use generative AI tools (ChatGPT, Claude, Gemini, Copilot, and the like) to study: to explain a concept, to check your reasoning on a practice problem, or to work through a homework question you are stuck on. Treat them as a tutor, not a substitute; the exams are closed book and closed device, and the only way to pass them is to be able to do the work yourself. Any use of AI or any other device during an exam is a violation of the academic integrity policy.

%-----------------------------------------------------------
\section*{Inclusion}
%-----------------------------------------------------------

Everyone in this class has different life experiences and perspectives, and all are valid. As the instructor, I will behave maturely and respectfully in all our class-related engagements, and I expect the same from you. We all adhere to Bentley's standards of conduct, the \href{https://www.bentley.edu/about/mission-and-values}{Bentley Core Values}, in class and in all course work. If you feel that I or anyone in this class has acted outside those values, come and talk to me. If you would rather not, speak with someone in the Office of Student Success (JEN 336, 781.891.2803). My roster shows your name as it appears in Workday; tell me early in the semester if you go by a different name or pronoun.

%-----------------------------------------------------------
\section*{Support and Resources}
%-----------------------------------------------------------

\subsection*{Student Accessibility Services}
Bentley University abides by Section 504 of the Rehabilitation Act of 1973 and the Americans with Disabilities Act of 1990, which stipulate that no student shall be denied the benefits of an education solely by reason of a disability. If you have a hidden or visible disability that may require classroom accommodations, contact \href{https://www.bentley.edu/offices/student-accessibility-services}{Student Accessibility Services} (Office of Student Success, JEN 336, 781.891.2004) within the first four weeks of the semester. Send me your accommodation plan as soon as you have it; accommodations are arranged at least one week before each exam.

\subsection*{The Howard A. Winer '58 Lab for Economics, Accounting and Finance (LEAF)}
The LEAF provides a welcoming and inclusive learning environment where students are encouraged to seek academic support for their Economics, Accounting and Finance courses. Students utilizing the LEAF will find knowledgeable peer tutors ready to help their colleagues to thrive in the Bentley business curriculum. LEAF's free, drop-in tutoring takes place in-person at Morison 220D (rear entrance), opening September 14. For additional information visit \url{https://www.bentley.edu/centers/leaf}. LEAF is a supplement to, not a substitute for, office hours.

\subsection*{Bentley policies and procedures}
All courses at Bentley reflect the university's commitment to a set of core values and practices. You are expected to be familiar with:
\begin{itemize}
    \item \href{https://catalog.bentley.edu/undergraduate/academic-policies-procedures/}{Bentley's honor code and academic integrity system}
    \item \href{https://www.bentley.edu/equity-reporting}{Equity and bias reporting forms and procedures} (Bias Incident Response Team, Office of Institutional Equity)
    \item \href{https://www.bentley.edu/offices/student-accessibility-services}{Student disability accommodations}
    \item \href{https://catalog.bentley.edu/undergraduate/academic-policies-procedures/}{Religious observances policy}
    \item \href{https://www.bentley.edu/about/mission-and-values}{Bentley's core values}
\end{itemize}

\subsection*{Course materials}
Slides, notes, handouts, and exams for this course are for your own use during this semester. They are protected by copyright; do not post or distribute them outside the class without my written consent. The course page is password-protected; do not share the password outside the class.

%-----------------------------------------------------------
\section*{Attendance Policy}
%-----------------------------------------------------------

All students must attend the in-person classes. In-person attendance is essential because it allows for real-time interaction, immediate feedback, and the kind of dialogue that makes learning stick. I do not record lectures, so if you miss a class, you miss the material.

Remote attendance is not an option. This course follows Bentley's \href{https://catalog.bentley.edu/undergraduate/academic-policies-procedures/}{Academic Engagement/Attendance Policy}: students who do not attend or engage during the first week are dropped as no-shows, and two or more consecutive weeks of non-participation can lead to administrative withdrawal. If you must miss class, you are responsible for the material; read the chapter and the posted notes, then bring specific questions to office hours.

\textbf{Weather and closures.} If the university cancels classes, I will email you whether there is an alternative assignment or arrangement for that day.

\textbf{Electronic devices.} Phones away and laptops and tablets closed for the whole class. The work in this course is drawing and reading diagrams, which you learn by doing them by hand, and a screen in your line of sight costs the people around you as much as it costs you. Using a device during class counts against your participation grade. Exceptions are by my permission only; if you have one, you never need to explain it to anyone in the room. Photographing the board at the end of class is fine.

\textbf{Religious observances.} Bentley provides reasonable academic adjustments when a sincerely held religious observance conflicts with a class, exam, or deadline (see the \href{https://downloads.bentley.edu/general/Religious\%20Observances\%202026-2027.pdf}{religious observances calendar} and the \href{https://catalog.bentley.edu/undergraduate/academic-policies-procedures/}{policy}). If you anticipate a conflict, tell me in advance, ideally in the first two weeks of the semester, so we can agree on an adjustment.

%-----------------------------------------------------------
\section*{Tentative Schedule}
%-----------------------------------------------------------

\begin{itemize}
    \item \textbf{Week 1 (Aug 31, Sep 2):} Ten Principles of Economics – Ch. 1 \& Thinking Like an Economist – Ch. 2
    \item \textbf{Week 2 (Sep 9):} Interdependence and the Gains from Trade – Ch. 3 \textit{(no class Sep 7, Labor Day)}
    \item \textbf{Week 3 (Sep 14, 16):} The Market Forces of Supply and Demand – Ch. 4
    \item \textbf{Week 4 (Sep 21, 23):} Elasticity and Its Application – Ch. 5
    \item \textbf{Week 5 (Sep 28, 30):} Supply, Demand, and Government Policies – Ch. 6
    \item \textbf{Week 6 (Oct 5, 7):} Review / \textbf{First Midterm (Oct 7)}
    \item \textbf{Week 7 (Oct 14):} Consumers, Producers, and the Efficiency of Markets – Ch. 7 \textit{(no class Oct 12, Fall Break)}
    \item \textbf{Week 8 (Oct 19, 21):} The Costs of Taxation – Ch. 8 \& International Trade – Ch. 9
    \item \textbf{Week 9 (Oct 26, 28):} Externalities – Ch. 10 \& Public Goods and Common Resources – Ch. 11
    \item \textbf{Week 10 (Nov 2, 4):} The Costs of Production – Ch. 14
    \item \textbf{Week 11 (Nov 9, 11):} Review / \textbf{Second Midterm (Nov 11)}
    \item \textbf{Week 12 (Nov 16, 18):} Firms in Competitive Markets – Ch. 15
    \item \textbf{Week 13 (Nov 23):} Monopoly – Ch. 16 \textit{(no class Nov 25, Thanksgiving Break)}
    \item \textbf{Week 14 (Nov 30, Dec 2):} Monopolistic Competition – Ch. 17 \& Oligopoly – Ch. 18
    \item \textbf{Week 15 (Dec 7):} Review
    \item \textbf{Final Exam:} \finaltime
\end{itemize}

\textit{Note: Chapter references are from Mankiw's \textit{Principles of Microeconomics}, 10th edition (the MindTap e-book). The schedule, and if necessary other parts of this syllabus, may change during the semester; changes are announced in class and by email.}

\end{document}
"
%   xelatex -jobname=EC111_FA2026_Syllabus_20  "\documentclass[11pt,letterpaper]{article}

% Compile with XeLaTeX. Two sections from one source:
%   xelatex -jobname=EC111_FA2026_Syllabus_10H "\def\honors{1}\input{EC111_FA2026_Syllabus.tex}"
%   xelatex -jobname=EC111_FA2026_Syllabus_20  "\input{EC111_FA2026_Syllabus.tex}"
\ifdefined\honors
  \newcommand{\sectioncode}{EC 111-10-H}
  \newcommand{\sectionname}{EC 111-10-H Principles of Microeconomics (Honors)}
  \newcommand{\meetingtime}{Monday/Wednesday, 3:30 PM -- 4:50 PM}
  \newcommand{\classroom}{Smith Technology Center 103}
  \newcommand{\finaltime}{Friday, December 11, 2026, 6:00 PM -- 8:00 PM}
  \newcommand{\mindtapsection}{Section 10-H of EC 111 Principles of Microeconomics Fall 2026}
  \newcommand{\mindtaplink}{https://student.cengage.com/course-link/MTPNS57GW6XN}
  \newcommand{\mindtapaccess}{the end of Sunday, September 13, 2026}
\else
  \newcommand{\sectioncode}{EC 111-20}
  \newcommand{\sectionname}{EC 111-20 Principles of Microeconomics}
  \newcommand{\meetingtime}{Monday/Wednesday, 9:30 AM -- 10:50 AM}
  \newcommand{\classroom}{Jennison Hall 412}
  \newcommand{\finaltime}{Friday, December 11, 2026, 8:00 AM -- 10:00 AM}
  \newcommand{\mindtapsection}{Section 20 of EC 111 Principles of Microeconomics Fall 2026}
  \newcommand{\mindtaplink}{https://student.cengage.com/course-link/MTPP8N7GWX7L}
  \newcommand{\mindtapaccess}{the end of Sunday, September 13, 2026}
\fi

% XeLaTeX for Turkish character support
\usepackage{fontspec}
\setmainfont{DejaVuSerif}[Extension=.ttf, BoldFont=DejaVuSerif-Bold, ItalicFont=DejaVuSerif-Italic, BoldItalicFont=DejaVuSerif-BoldItalic]
\setsansfont{DejaVuSans}[Extension=.ttf, BoldFont=DejaVuSans-Bold, ItalicFont=DejaVuSans-Oblique, BoldItalicFont=DejaVuSans-BoldOblique]

% Page layout
\usepackage[margin=1in]{geometry}
\usepackage{setspace}
\onehalfspacing

% Colors and styling
\usepackage[dvipsnames]{xcolor}
\definecolor{bentleyblue}{RGB}{0, 51, 102}
\definecolor{linkblue}{RGB}{0, 102, 204}

% Headers and sections
\usepackage{titlesec}
\titleformat{\section}{\Large\bfseries\color{bentleyblue}}{\thesection}{1em}{}[\titlerule]
\titleformat{\subsection}{\large\bfseries\color{bentleyblue!80}}{\thesubsection}{1em}{}
\titleformat{\subsubsection}{\normalsize\bfseries\color{bentleyblue!60}}{\thesubsubsection}{1em}{}

% Lists
\usepackage{enumitem}
\setlist[itemize]{leftmargin=*, itemsep=2pt, parsep=0pt}
\setlist[enumerate]{leftmargin=*, itemsep=2pt, parsep=0pt}

% Tables
\usepackage{booktabs}
\usepackage{tabularx}

% Links
\usepackage{hyperref}
\hypersetup{
    colorlinks=true,
    linkcolor=linkblue,
    urlcolor=linkblue,
    pdfauthor={Onur Altındağ},
    pdftitle={EC111 Principles of Microeconomics - Fall 2026 Syllabus}
}

% Header/Footer
\usepackage{fancyhdr}
\pagestyle{fancy}
\fancyhf{}
\fancyhead[L]{\small\color{gray}\sectioncode{} Principles of Microeconomics}
\fancyhead[R]{\small\color{gray}Fall 2026}
\fancyfoot[C]{\thepage}
\renewcommand{\headrulewidth}{0.4pt}

% Remove paragraph indent, add space between paragraphs
\setlength{\parindent}{0pt}
\setlength{\parskip}{6pt}

\begin{document}

% Title
\begin{center}
{\LARGE\bfseries\color{bentleyblue} EC 111: Principles of Microeconomics}\\[0.3cm]
{\Large Fall 2026 Syllabus}\\[0.5cm]
{\large Bentley University}
\end{center}

\vspace{0.5cm}

%-----------------------------------------------------------
\section*{Course Description}
%-----------------------------------------------------------

This course introduces the fundamental principles and tools of microeconomics: how individuals and firms make decisions when resources are scarce, and how markets coordinate those decisions. We study demand, supply, and market equilibrium; elasticity; the effects of government policies such as taxes and price controls; the efficiency of markets and the sources of market failure; the costs of production; and firm behavior in different market structures, from perfect competition to monopoly, monopolistic competition, and oligopoly.

The emphasis throughout is on economic reasoning: thinking at the margin, recognizing trade-offs, and using simple graphical and numerical models to make sense of real-world questions in business and policy.

\subsection*{Knowledge}
\begin{itemize}
    \item Scarcity, opportunity cost, trade-offs, and the gains from trade.
    \item Demand, supply, market equilibrium, and how markets respond to shocks.
    \item Elasticity and its role in pricing, revenue, and tax incidence.
    \item Consumer surplus, producer surplus, market efficiency, and the deadweight loss of taxes and price controls.
    \item Externalities, public goods, and the case for government intervention.
    \item Costs of production and firm behavior under perfect competition, monopoly, monopolistic competition, and oligopoly.
\end{itemize}

\subsection*{Skills}
\begin{itemize}
    \item Use supply and demand diagrams to predict the effects of changes in market conditions and policy.
    \item Compute and interpret elasticities, surpluses, costs, and profit-maximizing output and price.
    \item Read and evaluate economic arguments critically and communicate them clearly.
\end{itemize}

\subsection*{Perspectives}
\begin{itemize}
    \item Understand marginal decision making as the core of economic reasoning.
    \item Understand when markets work well, when they fail, and what policy can and cannot fix.
\end{itemize}

%-----------------------------------------------------------
\section*{Class Information}
%-----------------------------------------------------------

\begin{tabularx}{\textwidth}{@{}lX@{}}
\textbf{Instructor:} & Onur Altındağ \\
\textbf{Office:} & RAU 019F \\
\textbf{Email:} & \href{mailto:oaltindag@bentley.edu}{oaltindag@bentley.edu} \\
\textbf{Web:} & \url{https://www.onuraltindag.info/} \\
\textbf{Course page:} & \url{https://www.onuraltindag.info/teaching/ec111/} \\
\textbf{Password:} & \texttt{GoFalcons!} \\[0.3cm]
\textbf{Course:} & \sectionname \\
\textbf{Semester:} & Fall 2026 (August 31, 2026 – December 15, 2026) \\
\textbf{Meeting Time:} & \meetingtime \\
\textbf{Location:} & \classroom \\
\textbf{Delivery Mode:} & In-Person Lecture \\
\end{tabularx}

\vspace{0.3cm}

Office hours are on Mondays and Wednesdays, 11:00 AM -- 12:00 PM and 1:00 -- 3:00 PM, by appointment. Please \href{https://outlook.office.com/bookwithme/user/1c4e4f54e0f0493397e615b2c80a59c7@bentley.edu/meetingtype/rwPThvM0MkK5oGlK_t37nw2?bookingcode=17c9a7b2-69fa-4c48-85ea-ad404642334b&anonymous&ismsaljsauthenabled&ep=mlink}{book a 30-minute slot} to schedule a meeting. If you need to meet urgently or no slots are available, email me directly.

Announcements come by email from me to your Bentley address; check it regularly. Slides, lecture notes, handouts, and mock exams are posted on the course page above after each class. When you email me, put \textbf{EC111} in the subject line.

%-----------------------------------------------------------
\section*{Important Dates and Evaluation}
%-----------------------------------------------------------

\begin{tabularx}{\textwidth}{@{}Xr@{}}
\toprule
\textbf{Component} & \textbf{Weight} \\
\midrule
Weekly assignments + participation & 20\% \\
\textbf{Wednesday, October 7, 2026} – First Midterm & 25\% \\
\textbf{Wednesday, November 11, 2026} – Second Midterm & 25\% \\
\textbf{\finaltime} – Final Exam & 30\% \\
\bottomrule
\end{tabularx}

%-----------------------------------------------------------
\section*{Grading}
%-----------------------------------------------------------

You \textbf{MUST} attend all the midterms and the final as there will be no make-up exams in this course. The midterms are \textbf{NOT} cumulative. If you miss or are likely to miss a midterm due to an emergency, please contact me as soon as possible. You will need to provide supporting documentation/verification of your absence. I will re-weight your final exam accordingly if you have a valid excuse. Please note that family vacations or missing the school shuttle are not valid excuses. If a serious illness, a death in the family, or another difficult circumstance keeps you from an exam or assignment, contact the Office of Student Success (JEN 336, 781.891.2803); they will notify all of your instructors on your behalf.

The final exam \textbf{is} cumulative. If you miss the final exam due to an emergency, you will receive an incomplete for this course. Do not take this class if you know that you will not be able to attend the final exam.

Exams are closed book and closed notes. A basic (non-graphing) calculator is allowed; phones, graphing calculators, and any other device are not, and calculators may not be shared.

You will have \textbf{10 homework assignments} on MindTap, roughly one per week (none in exam weeks), due \textbf{Sunday 11:00 PM}. Your \textbf{two lowest scores are dropped}; the remaining 8 count equally. Each question allows three attempts; an assignment scoring at least \textbf{50\%} receives full credit, otherwise zero. The submission system is automated and will not accept anything after the deadline (even 5 minutes), and there are no exceptions. The questions are algorithmic and slightly different for everyone, so you are encouraged to work together.

\subsection*{Letter grades}
\begin{center}
\begin{tabular}{@{}llllll@{}}
\toprule
A & 94--100 & B+ & 87--89 & C+ & 77--79 \\
A$-$ & 90--93 & B & 83--86 & C & 73--76 \\
 & & B$-$ & 80--82 & C$-$ & 70--72 \\
\midrule
D+ & 67--69 & D & 63--66 & D$-$ & 60--62 \\
F & below 60 & & & & \\
\bottomrule
\end{tabular}
\end{center}
I may lower these cutoffs at the end of the semester; I will not raise them. I do not respond to requests for individual extra credit or for grade adjustments after the fact.

%-----------------------------------------------------------
\section*{Textbook and MindTap}
%-----------------------------------------------------------

The course uses \textbf{MindTap for Mankiw, \textit{Principles of Microeconomics}} (Cengage). The e-book, your homework assignments, and practice questions all live there; a hard copy of the textbook is optional. You \textbf{MUST} register, using your Bentley email, through the link for your section:

\begin{center}
\textbf{\mindtapsection}\\
\url{\mindtaplink}
\end{center}

You have free temporary access until \mindtapaccess; after that you will be prompted to purchase (if your course materials are included in your tuition, no purchase is needed). Registration help: \url{https://startstrong.cengage.com}; technical support: \url{https://support.cengage.com}. Register in the first week; MindTap problems are not an excuse for a missed assignment.

\ifdefined\honors
%-----------------------------------------------------------
\section*{Honors Reading List}
%-----------------------------------------------------------

In addition to the textbook, honors students read and discuss a short set of articles and book chapters over the semester. The list and discussion dates will be posted on the course page in the first weeks of the semester.

\fi
%-----------------------------------------------------------
\section*{Academic Integrity}
%-----------------------------------------------------------

All students are expected to adhere to Bentley's Academic Integrity policy, which includes Bentley's Honor Code (details are in the Undergraduate Student Handbook, the \href{https://catalog.bentley.edu/undergraduate/academic-policies-procedures/}{academic policies section of the catalog}, and the academic integrity course page on Brightspace in which all students and faculty are enrolled). The essence of the policy is that you should not represent someone else's work as your own: no plagiarism, no cheating on exams, no illicit collaboration. Failure to adhere to the policy has serious consequences, including course failure, suspension, or expulsion from the university. The best way to avoid a problem is to consult with me before taking an action that might constitute a violation.

\subsection*{Generative AI}
Nothing in this course is graded on writing, and the homework is algorithmic practice, so the line is simple. You may use generative AI tools (ChatGPT, Claude, Gemini, Copilot, and the like) to study: to explain a concept, to check your reasoning on a practice problem, or to work through a homework question you are stuck on. Treat them as a tutor, not a substitute; the exams are closed book and closed device, and the only way to pass them is to be able to do the work yourself. Any use of AI or any other device during an exam is a violation of the academic integrity policy.

%-----------------------------------------------------------
\section*{Inclusion}
%-----------------------------------------------------------

Everyone in this class has different life experiences and perspectives, and all are valid. As the instructor, I will behave maturely and respectfully in all our class-related engagements, and I expect the same from you. We all adhere to Bentley's standards of conduct, the \href{https://www.bentley.edu/about/mission-and-values}{Bentley Core Values}, in class and in all course work. If you feel that I or anyone in this class has acted outside those values, come and talk to me. If you would rather not, speak with someone in the Office of Student Success (JEN 336, 781.891.2803). My roster shows your name as it appears in Workday; tell me early in the semester if you go by a different name or pronoun.

%-----------------------------------------------------------
\section*{Support and Resources}
%-----------------------------------------------------------

\subsection*{Student Accessibility Services}
Bentley University abides by Section 504 of the Rehabilitation Act of 1973 and the Americans with Disabilities Act of 1990, which stipulate that no student shall be denied the benefits of an education solely by reason of a disability. If you have a hidden or visible disability that may require classroom accommodations, contact \href{https://www.bentley.edu/offices/student-accessibility-services}{Student Accessibility Services} (Office of Student Success, JEN 336, 781.891.2004) within the first four weeks of the semester. Send me your accommodation plan as soon as you have it; accommodations are arranged at least one week before each exam.

\subsection*{The Howard A. Winer '58 Lab for Economics, Accounting and Finance (LEAF)}
The LEAF provides a welcoming and inclusive learning environment where students are encouraged to seek academic support for their Economics, Accounting and Finance courses. Students utilizing the LEAF will find knowledgeable peer tutors ready to help their colleagues to thrive in the Bentley business curriculum. LEAF's free, drop-in tutoring takes place in-person at Morison 220D (rear entrance), opening September 14. For additional information visit \url{https://www.bentley.edu/centers/leaf}. LEAF is a supplement to, not a substitute for, office hours.

\subsection*{Bentley policies and procedures}
All courses at Bentley reflect the university's commitment to a set of core values and practices. You are expected to be familiar with:
\begin{itemize}
    \item \href{https://catalog.bentley.edu/undergraduate/academic-policies-procedures/}{Bentley's honor code and academic integrity system}
    \item \href{https://www.bentley.edu/equity-reporting}{Equity and bias reporting forms and procedures} (Bias Incident Response Team, Office of Institutional Equity)
    \item \href{https://www.bentley.edu/offices/student-accessibility-services}{Student disability accommodations}
    \item \href{https://catalog.bentley.edu/undergraduate/academic-policies-procedures/}{Religious observances policy}
    \item \href{https://www.bentley.edu/about/mission-and-values}{Bentley's core values}
\end{itemize}

\subsection*{Course materials}
Slides, notes, handouts, and exams for this course are for your own use during this semester. They are protected by copyright; do not post or distribute them outside the class without my written consent. The course page is password-protected; do not share the password outside the class.

%-----------------------------------------------------------
\section*{Attendance Policy}
%-----------------------------------------------------------

All students must attend the in-person classes. In-person attendance is essential because it allows for real-time interaction, immediate feedback, and the kind of dialogue that makes learning stick. I do not record lectures, so if you miss a class, you miss the material.

Remote attendance is not an option. This course follows Bentley's \href{https://catalog.bentley.edu/undergraduate/academic-policies-procedures/}{Academic Engagement/Attendance Policy}: students who do not attend or engage during the first week are dropped as no-shows, and two or more consecutive weeks of non-participation can lead to administrative withdrawal. If you must miss class, you are responsible for the material; read the chapter and the posted notes, then bring specific questions to office hours.

\textbf{Weather and closures.} If the university cancels classes, I will email you whether there is an alternative assignment or arrangement for that day.

\textbf{Electronic devices.} Phones away and laptops and tablets closed for the whole class. The work in this course is drawing and reading diagrams, which you learn by doing them by hand, and a screen in your line of sight costs the people around you as much as it costs you. Using a device during class counts against your participation grade. Exceptions are by my permission only; if you have one, you never need to explain it to anyone in the room. Photographing the board at the end of class is fine.

\textbf{Religious observances.} Bentley provides reasonable academic adjustments when a sincerely held religious observance conflicts with a class, exam, or deadline (see the \href{https://downloads.bentley.edu/general/Religious\%20Observances\%202026-2027.pdf}{religious observances calendar} and the \href{https://catalog.bentley.edu/undergraduate/academic-policies-procedures/}{policy}). If you anticipate a conflict, tell me in advance, ideally in the first two weeks of the semester, so we can agree on an adjustment.

%-----------------------------------------------------------
\section*{Tentative Schedule}
%-----------------------------------------------------------

\begin{itemize}
    \item \textbf{Week 1 (Aug 31, Sep 2):} Ten Principles of Economics – Ch. 1 \& Thinking Like an Economist – Ch. 2
    \item \textbf{Week 2 (Sep 9):} Interdependence and the Gains from Trade – Ch. 3 \textit{(no class Sep 7, Labor Day)}
    \item \textbf{Week 3 (Sep 14, 16):} The Market Forces of Supply and Demand – Ch. 4
    \item \textbf{Week 4 (Sep 21, 23):} Elasticity and Its Application – Ch. 5
    \item \textbf{Week 5 (Sep 28, 30):} Supply, Demand, and Government Policies – Ch. 6
    \item \textbf{Week 6 (Oct 5, 7):} Review / \textbf{First Midterm (Oct 7)}
    \item \textbf{Week 7 (Oct 14):} Consumers, Producers, and the Efficiency of Markets – Ch. 7 \textit{(no class Oct 12, Fall Break)}
    \item \textbf{Week 8 (Oct 19, 21):} The Costs of Taxation – Ch. 8 \& International Trade – Ch. 9
    \item \textbf{Week 9 (Oct 26, 28):} Externalities – Ch. 10 \& Public Goods and Common Resources – Ch. 11
    \item \textbf{Week 10 (Nov 2, 4):} The Costs of Production – Ch. 14
    \item \textbf{Week 11 (Nov 9, 11):} Review / \textbf{Second Midterm (Nov 11)}
    \item \textbf{Week 12 (Nov 16, 18):} Firms in Competitive Markets – Ch. 15
    \item \textbf{Week 13 (Nov 23):} Monopoly – Ch. 16 \textit{(no class Nov 25, Thanksgiving Break)}
    \item \textbf{Week 14 (Nov 30, Dec 2):} Monopolistic Competition – Ch. 17 \& Oligopoly – Ch. 18
    \item \textbf{Week 15 (Dec 7):} Review
    \item \textbf{Final Exam:} \finaltime
\end{itemize}

\textit{Note: Chapter references are from Mankiw's \textit{Principles of Microeconomics}, 10th edition (the MindTap e-book). The schedule, and if necessary other parts of this syllabus, may change during the semester; changes are announced in class and by email.}

\end{document}
"
\ifdefined\honors
  \newcommand{\sectioncode}{EC 111-10-H}
  \newcommand{\sectionname}{EC 111-10-H Principles of Microeconomics (Honors)}
  \newcommand{\meetingtime}{Monday/Wednesday, 3:30 PM -- 4:50 PM}
  \newcommand{\classroom}{Smith Technology Center 103}
  \newcommand{\finaltime}{Friday, December 11, 2026, 6:00 PM -- 8:00 PM}
  \newcommand{\mindtapsection}{Section 10-H of EC 111 Principles of Microeconomics Fall 2026}
  \newcommand{\mindtaplink}{https://student.cengage.com/course-link/MTPNS57GW6XN}
  \newcommand{\mindtapaccess}{the end of Sunday, September 13, 2026}
\else
  \newcommand{\sectioncode}{EC 111-20}
  \newcommand{\sectionname}{EC 111-20 Principles of Microeconomics}
  \newcommand{\meetingtime}{Monday/Wednesday, 9:30 AM -- 10:50 AM}
  \newcommand{\classroom}{Jennison Hall 412}
  \newcommand{\finaltime}{Friday, December 11, 2026, 8:00 AM -- 10:00 AM}
  \newcommand{\mindtapsection}{Section 20 of EC 111 Principles of Microeconomics Fall 2026}
  \newcommand{\mindtaplink}{https://student.cengage.com/course-link/MTPP8N7GWX7L}
  \newcommand{\mindtapaccess}{the end of Sunday, September 13, 2026}
\fi

% XeLaTeX for Turkish character support
\usepackage{fontspec}
\setmainfont{DejaVuSerif}[Extension=.ttf, BoldFont=DejaVuSerif-Bold, ItalicFont=DejaVuSerif-Italic, BoldItalicFont=DejaVuSerif-BoldItalic]
\setsansfont{DejaVuSans}[Extension=.ttf, BoldFont=DejaVuSans-Bold, ItalicFont=DejaVuSans-Oblique, BoldItalicFont=DejaVuSans-BoldOblique]

% Page layout
\usepackage[margin=1in]{geometry}
\usepackage{setspace}
\onehalfspacing

% Colors and styling
\usepackage[dvipsnames]{xcolor}
\definecolor{bentleyblue}{RGB}{0, 51, 102}
\definecolor{linkblue}{RGB}{0, 102, 204}

% Headers and sections
\usepackage{titlesec}
\titleformat{\section}{\Large\bfseries\color{bentleyblue}}{\thesection}{1em}{}[\titlerule]
\titleformat{\subsection}{\large\bfseries\color{bentleyblue!80}}{\thesubsection}{1em}{}
\titleformat{\subsubsection}{\normalsize\bfseries\color{bentleyblue!60}}{\thesubsubsection}{1em}{}

% Lists
\usepackage{enumitem}
\setlist[itemize]{leftmargin=*, itemsep=2pt, parsep=0pt}
\setlist[enumerate]{leftmargin=*, itemsep=2pt, parsep=0pt}

% Tables
\usepackage{booktabs}
\usepackage{tabularx}

% Links
\usepackage{hyperref}
\hypersetup{
    colorlinks=true,
    linkcolor=linkblue,
    urlcolor=linkblue,
    pdfauthor={Onur Altındağ},
    pdftitle={EC111 Principles of Microeconomics - Fall 2026 Syllabus}
}

% Header/Footer
\usepackage{fancyhdr}
\pagestyle{fancy}
\fancyhf{}
\fancyhead[L]{\small\color{gray}\sectioncode{} Principles of Microeconomics}
\fancyhead[R]{\small\color{gray}Fall 2026}
\fancyfoot[C]{\thepage}
\renewcommand{\headrulewidth}{0.4pt}

% Remove paragraph indent, add space between paragraphs
\setlength{\parindent}{0pt}
\setlength{\parskip}{6pt}

\begin{document}

% Title
\begin{center}
{\LARGE\bfseries\color{bentleyblue} EC 111: Principles of Microeconomics}\\[0.3cm]
{\Large Fall 2026 Syllabus}\\[0.5cm]
{\large Bentley University}
\end{center}

\vspace{0.5cm}

%-----------------------------------------------------------
\section*{Course Description}
%-----------------------------------------------------------

This course introduces the fundamental principles and tools of microeconomics: how individuals and firms make decisions when resources are scarce, and how markets coordinate those decisions. We study demand, supply, and market equilibrium; elasticity; the effects of government policies such as taxes and price controls; the efficiency of markets and the sources of market failure; the costs of production; and firm behavior in different market structures, from perfect competition to monopoly, monopolistic competition, and oligopoly.

The emphasis throughout is on economic reasoning: thinking at the margin, recognizing trade-offs, and using simple graphical and numerical models to make sense of real-world questions in business and policy.

\subsection*{Knowledge}
\begin{itemize}
    \item Scarcity, opportunity cost, trade-offs, and the gains from trade.
    \item Demand, supply, market equilibrium, and how markets respond to shocks.
    \item Elasticity and its role in pricing, revenue, and tax incidence.
    \item Consumer surplus, producer surplus, market efficiency, and the deadweight loss of taxes and price controls.
    \item Externalities, public goods, and the case for government intervention.
    \item Costs of production and firm behavior under perfect competition, monopoly, monopolistic competition, and oligopoly.
\end{itemize}

\subsection*{Skills}
\begin{itemize}
    \item Use supply and demand diagrams to predict the effects of changes in market conditions and policy.
    \item Compute and interpret elasticities, surpluses, costs, and profit-maximizing output and price.
    \item Read and evaluate economic arguments critically and communicate them clearly.
\end{itemize}

\subsection*{Perspectives}
\begin{itemize}
    \item Understand marginal decision making as the core of economic reasoning.
    \item Understand when markets work well, when they fail, and what policy can and cannot fix.
\end{itemize}

%-----------------------------------------------------------
\section*{Class Information}
%-----------------------------------------------------------

\begin{tabularx}{\textwidth}{@{}lX@{}}
\textbf{Instructor:} & Onur Altındağ \\
\textbf{Office:} & RAU 019F \\
\textbf{Email:} & \href{mailto:oaltindag@bentley.edu}{oaltindag@bentley.edu} \\
\textbf{Web:} & \url{https://www.onuraltindag.info/} \\
\textbf{Course page:} & \url{https://www.onuraltindag.info/teaching/ec111/} \\
\textbf{Password:} & \texttt{GoFalcons!} \\[0.3cm]
\textbf{Course:} & \sectionname \\
\textbf{Semester:} & Fall 2026 (August 31, 2026 – December 15, 2026) \\
\textbf{Meeting Time:} & \meetingtime \\
\textbf{Location:} & \classroom \\
\textbf{Delivery Mode:} & In-Person Lecture \\
\end{tabularx}

\vspace{0.3cm}

Office hours are on Mondays and Wednesdays, 11:00 AM -- 12:00 PM and 1:00 -- 3:00 PM, by appointment. Please \href{https://outlook.office.com/bookwithme/user/1c4e4f54e0f0493397e615b2c80a59c7@bentley.edu/meetingtype/rwPThvM0MkK5oGlK_t37nw2?bookingcode=17c9a7b2-69fa-4c48-85ea-ad404642334b&anonymous&ismsaljsauthenabled&ep=mlink}{book a 30-minute slot} to schedule a meeting. If you need to meet urgently or no slots are available, email me directly.

Announcements come by email from me to your Bentley address; check it regularly. Slides, lecture notes, handouts, and mock exams are posted on the course page above after each class. When you email me, put \textbf{EC111} in the subject line.

%-----------------------------------------------------------
\section*{Important Dates and Evaluation}
%-----------------------------------------------------------

\begin{tabularx}{\textwidth}{@{}Xr@{}}
\toprule
\textbf{Component} & \textbf{Weight} \\
\midrule
Weekly assignments + participation & 20\% \\
\textbf{Wednesday, October 7, 2026} – First Midterm & 25\% \\
\textbf{Wednesday, November 11, 2026} – Second Midterm & 25\% \\
\textbf{\finaltime} – Final Exam & 30\% \\
\bottomrule
\end{tabularx}

%-----------------------------------------------------------
\section*{Grading}
%-----------------------------------------------------------

You \textbf{MUST} attend all the midterms and the final as there will be no make-up exams in this course. The midterms are \textbf{NOT} cumulative. If you miss or are likely to miss a midterm due to an emergency, please contact me as soon as possible. You will need to provide supporting documentation/verification of your absence. I will re-weight your final exam accordingly if you have a valid excuse. Please note that family vacations or missing the school shuttle are not valid excuses. If a serious illness, a death in the family, or another difficult circumstance keeps you from an exam or assignment, contact the Office of Student Success (JEN 336, 781.891.2803); they will notify all of your instructors on your behalf.

The final exam \textbf{is} cumulative. If you miss the final exam due to an emergency, you will receive an incomplete for this course. Do not take this class if you know that you will not be able to attend the final exam.

Exams are closed book and closed notes. A basic (non-graphing) calculator is allowed; phones, graphing calculators, and any other device are not, and calculators may not be shared.

You will have \textbf{10 homework assignments} on MindTap, roughly one per week (none in exam weeks), due \textbf{Sunday 11:00 PM}. Your \textbf{two lowest scores are dropped}; the remaining 8 count equally. Each question allows three attempts; an assignment scoring at least \textbf{50\%} receives full credit, otherwise zero. The submission system is automated and will not accept anything after the deadline (even 5 minutes), and there are no exceptions. The questions are algorithmic and slightly different for everyone, so you are encouraged to work together.

\subsection*{Letter grades}
\begin{center}
\begin{tabular}{@{}llllll@{}}
\toprule
A & 94--100 & B+ & 87--89 & C+ & 77--79 \\
A$-$ & 90--93 & B & 83--86 & C & 73--76 \\
 & & B$-$ & 80--82 & C$-$ & 70--72 \\
\midrule
D+ & 67--69 & D & 63--66 & D$-$ & 60--62 \\
F & below 60 & & & & \\
\bottomrule
\end{tabular}
\end{center}
I may lower these cutoffs at the end of the semester; I will not raise them. I do not respond to requests for individual extra credit or for grade adjustments after the fact.

%-----------------------------------------------------------
\section*{Textbook and MindTap}
%-----------------------------------------------------------

The course uses \textbf{MindTap for Mankiw, \textit{Principles of Microeconomics}} (Cengage). The e-book, your homework assignments, and practice questions all live there; a hard copy of the textbook is optional. You \textbf{MUST} register, using your Bentley email, through the link for your section:

\begin{center}
\textbf{\mindtapsection}\\
\url{\mindtaplink}
\end{center}

You have free temporary access until \mindtapaccess; after that you will be prompted to purchase (if your course materials are included in your tuition, no purchase is needed). Registration help: \url{https://startstrong.cengage.com}; technical support: \url{https://support.cengage.com}. Register in the first week; MindTap problems are not an excuse for a missed assignment.

\ifdefined\honors
%-----------------------------------------------------------
\section*{Honors Reading List}
%-----------------------------------------------------------

In addition to the textbook, honors students read and discuss a short set of articles and book chapters over the semester. The list and discussion dates will be posted on the course page in the first weeks of the semester.

\fi
%-----------------------------------------------------------
\section*{Academic Integrity}
%-----------------------------------------------------------

All students are expected to adhere to Bentley's Academic Integrity policy, which includes Bentley's Honor Code (details are in the Undergraduate Student Handbook, the \href{https://catalog.bentley.edu/undergraduate/academic-policies-procedures/}{academic policies section of the catalog}, and the academic integrity course page on Brightspace in which all students and faculty are enrolled). The essence of the policy is that you should not represent someone else's work as your own: no plagiarism, no cheating on exams, no illicit collaboration. Failure to adhere to the policy has serious consequences, including course failure, suspension, or expulsion from the university. The best way to avoid a problem is to consult with me before taking an action that might constitute a violation.

\subsection*{Generative AI}
Nothing in this course is graded on writing, and the homework is algorithmic practice, so the line is simple. You may use generative AI tools (ChatGPT, Claude, Gemini, Copilot, and the like) to study: to explain a concept, to check your reasoning on a practice problem, or to work through a homework question you are stuck on. Treat them as a tutor, not a substitute; the exams are closed book and closed device, and the only way to pass them is to be able to do the work yourself. Any use of AI or any other device during an exam is a violation of the academic integrity policy.

%-----------------------------------------------------------
\section*{Inclusion}
%-----------------------------------------------------------

Everyone in this class has different life experiences and perspectives, and all are valid. As the instructor, I will behave maturely and respectfully in all our class-related engagements, and I expect the same from you. We all adhere to Bentley's standards of conduct, the \href{https://www.bentley.edu/about/mission-and-values}{Bentley Core Values}, in class and in all course work. If you feel that I or anyone in this class has acted outside those values, come and talk to me. If you would rather not, speak with someone in the Office of Student Success (JEN 336, 781.891.2803). My roster shows your name as it appears in Workday; tell me early in the semester if you go by a different name or pronoun.

%-----------------------------------------------------------
\section*{Support and Resources}
%-----------------------------------------------------------

\subsection*{Student Accessibility Services}
Bentley University abides by Section 504 of the Rehabilitation Act of 1973 and the Americans with Disabilities Act of 1990, which stipulate that no student shall be denied the benefits of an education solely by reason of a disability. If you have a hidden or visible disability that may require classroom accommodations, contact \href{https://www.bentley.edu/offices/student-accessibility-services}{Student Accessibility Services} (Office of Student Success, JEN 336, 781.891.2004) within the first four weeks of the semester. Send me your accommodation plan as soon as you have it; accommodations are arranged at least one week before each exam.

\subsection*{The Howard A. Winer '58 Lab for Economics, Accounting and Finance (LEAF)}
The LEAF provides a welcoming and inclusive learning environment where students are encouraged to seek academic support for their Economics, Accounting and Finance courses. Students utilizing the LEAF will find knowledgeable peer tutors ready to help their colleagues to thrive in the Bentley business curriculum. LEAF's free, drop-in tutoring takes place in-person at Morison 220D (rear entrance), opening September 14. For additional information visit \url{https://www.bentley.edu/centers/leaf}. LEAF is a supplement to, not a substitute for, office hours.

\subsection*{Bentley policies and procedures}
All courses at Bentley reflect the university's commitment to a set of core values and practices. You are expected to be familiar with:
\begin{itemize}
    \item \href{https://catalog.bentley.edu/undergraduate/academic-policies-procedures/}{Bentley's honor code and academic integrity system}
    \item \href{https://www.bentley.edu/equity-reporting}{Equity and bias reporting forms and procedures} (Bias Incident Response Team, Office of Institutional Equity)
    \item \href{https://www.bentley.edu/offices/student-accessibility-services}{Student disability accommodations}
    \item \href{https://catalog.bentley.edu/undergraduate/academic-policies-procedures/}{Religious observances policy}
    \item \href{https://www.bentley.edu/about/mission-and-values}{Bentley's core values}
\end{itemize}

\subsection*{Course materials}
Slides, notes, handouts, and exams for this course are for your own use during this semester. They are protected by copyright; do not post or distribute them outside the class without my written consent. The course page is password-protected; do not share the password outside the class.

%-----------------------------------------------------------
\section*{Attendance Policy}
%-----------------------------------------------------------

All students must attend the in-person classes. In-person attendance is essential because it allows for real-time interaction, immediate feedback, and the kind of dialogue that makes learning stick. I do not record lectures, so if you miss a class, you miss the material.

Remote attendance is not an option. This course follows Bentley's \href{https://catalog.bentley.edu/undergraduate/academic-policies-procedures/}{Academic Engagement/Attendance Policy}: students who do not attend or engage during the first week are dropped as no-shows, and two or more consecutive weeks of non-participation can lead to administrative withdrawal. If you must miss class, you are responsible for the material; read the chapter and the posted notes, then bring specific questions to office hours.

\textbf{Weather and closures.} If the university cancels classes, I will email you whether there is an alternative assignment or arrangement for that day.

\textbf{Electronic devices.} Phones away and laptops and tablets closed for the whole class. The work in this course is drawing and reading diagrams, which you learn by doing them by hand, and a screen in your line of sight costs the people around you as much as it costs you. Using a device during class counts against your participation grade. Exceptions are by my permission only; if you have one, you never need to explain it to anyone in the room. Photographing the board at the end of class is fine.

\textbf{Religious observances.} Bentley provides reasonable academic adjustments when a sincerely held religious observance conflicts with a class, exam, or deadline (see the \href{https://downloads.bentley.edu/general/Religious\%20Observances\%202026-2027.pdf}{religious observances calendar} and the \href{https://catalog.bentley.edu/undergraduate/academic-policies-procedures/}{policy}). If you anticipate a conflict, tell me in advance, ideally in the first two weeks of the semester, so we can agree on an adjustment.

%-----------------------------------------------------------
\section*{Tentative Schedule}
%-----------------------------------------------------------

\begin{itemize}
    \item \textbf{Week 1 (Aug 31, Sep 2):} Ten Principles of Economics – Ch. 1 \& Thinking Like an Economist – Ch. 2
    \item \textbf{Week 2 (Sep 9):} Interdependence and the Gains from Trade – Ch. 3 \textit{(no class Sep 7, Labor Day)}
    \item \textbf{Week 3 (Sep 14, 16):} The Market Forces of Supply and Demand – Ch. 4
    \item \textbf{Week 4 (Sep 21, 23):} Elasticity and Its Application – Ch. 5
    \item \textbf{Week 5 (Sep 28, 30):} Supply, Demand, and Government Policies – Ch. 6
    \item \textbf{Week 6 (Oct 5, 7):} Review / \textbf{First Midterm (Oct 7)}
    \item \textbf{Week 7 (Oct 14):} Consumers, Producers, and the Efficiency of Markets – Ch. 7 \textit{(no class Oct 12, Fall Break)}
    \item \textbf{Week 8 (Oct 19, 21):} The Costs of Taxation – Ch. 8 \& International Trade – Ch. 9
    \item \textbf{Week 9 (Oct 26, 28):} Externalities – Ch. 10 \& Public Goods and Common Resources – Ch. 11
    \item \textbf{Week 10 (Nov 2, 4):} The Costs of Production – Ch. 14
    \item \textbf{Week 11 (Nov 9, 11):} Review / \textbf{Second Midterm (Nov 11)}
    \item \textbf{Week 12 (Nov 16, 18):} Firms in Competitive Markets – Ch. 15
    \item \textbf{Week 13 (Nov 23):} Monopoly – Ch. 16 \textit{(no class Nov 25, Thanksgiving Break)}
    \item \textbf{Week 14 (Nov 30, Dec 2):} Monopolistic Competition – Ch. 17 \& Oligopoly – Ch. 18
    \item \textbf{Week 15 (Dec 7):} Review
    \item \textbf{Final Exam:} \finaltime
\end{itemize}

\textit{Note: Chapter references are from Mankiw's \textit{Principles of Microeconomics}, 10th edition (the MindTap e-book). The schedule, and if necessary other parts of this syllabus, may change during the semester; changes are announced in class and by email.}

\end{document}
"
\ifdefined\honors
  \newcommand{\sectioncode}{EC 111-10-H}
  \newcommand{\sectionname}{EC 111-10-H Principles of Microeconomics (Honors)}
  \newcommand{\meetingtime}{Monday/Wednesday, 3:30 PM -- 4:50 PM}
  \newcommand{\classroom}{Smith Technology Center 103}
  \newcommand{\finaltime}{Friday, December 11, 2026, 6:00 PM -- 8:00 PM}
  \newcommand{\mindtapsection}{Section 10-H of EC 111 Principles of Microeconomics Fall 2026}
  \newcommand{\mindtaplink}{https://student.cengage.com/course-link/MTPNS57GW6XN}
  \newcommand{\mindtapaccess}{the end of Sunday, September 13, 2026}
\else
  \newcommand{\sectioncode}{EC 111-20}
  \newcommand{\sectionname}{EC 111-20 Principles of Microeconomics}
  \newcommand{\meetingtime}{Monday/Wednesday, 9:30 AM -- 10:50 AM}
  \newcommand{\classroom}{Jennison Hall 412}
  \newcommand{\finaltime}{Friday, December 11, 2026, 8:00 AM -- 10:00 AM}
  \newcommand{\mindtapsection}{Section 20 of EC 111 Principles of Microeconomics Fall 2026}
  \newcommand{\mindtaplink}{https://student.cengage.com/course-link/MTPP8N7GWX7L}
  \newcommand{\mindtapaccess}{the end of Sunday, September 13, 2026}
\fi

% XeLaTeX for Turkish character support
\usepackage{fontspec}
\setmainfont{DejaVuSerif}[Extension=.ttf, BoldFont=DejaVuSerif-Bold, ItalicFont=DejaVuSerif-Italic, BoldItalicFont=DejaVuSerif-BoldItalic]
\setsansfont{DejaVuSans}[Extension=.ttf, BoldFont=DejaVuSans-Bold, ItalicFont=DejaVuSans-Oblique, BoldItalicFont=DejaVuSans-BoldOblique]

% Page layout
\usepackage[margin=1in]{geometry}
\usepackage{setspace}
\onehalfspacing

% Colors and styling
\usepackage[dvipsnames]{xcolor}
\definecolor{bentleyblue}{RGB}{0, 51, 102}
\definecolor{linkblue}{RGB}{0, 102, 204}

% Headers and sections
\usepackage{titlesec}
\titleformat{\section}{\Large\bfseries\color{bentleyblue}}{\thesection}{1em}{}[\titlerule]
\titleformat{\subsection}{\large\bfseries\color{bentleyblue!80}}{\thesubsection}{1em}{}
\titleformat{\subsubsection}{\normalsize\bfseries\color{bentleyblue!60}}{\thesubsubsection}{1em}{}

% Lists
\usepackage{enumitem}
\setlist[itemize]{leftmargin=*, itemsep=2pt, parsep=0pt}
\setlist[enumerate]{leftmargin=*, itemsep=2pt, parsep=0pt}

% Tables
\usepackage{booktabs}
\usepackage{tabularx}

% Links
\usepackage{hyperref}
\hypersetup{
    colorlinks=true,
    linkcolor=linkblue,
    urlcolor=linkblue,
    pdfauthor={Onur Altındağ},
    pdftitle={EC111 Principles of Microeconomics - Fall 2026 Syllabus}
}

% Header/Footer
\usepackage{fancyhdr}
\pagestyle{fancy}
\fancyhf{}
\fancyhead[L]{\small\color{gray}\sectioncode{} Principles of Microeconomics}
\fancyhead[R]{\small\color{gray}Fall 2026}
\fancyfoot[C]{\thepage}
\renewcommand{\headrulewidth}{0.4pt}

% Remove paragraph indent, add space between paragraphs
\setlength{\parindent}{0pt}
\setlength{\parskip}{6pt}

\begin{document}

% Title
\begin{center}
{\LARGE\bfseries\color{bentleyblue} EC 111: Principles of Microeconomics}\\[0.3cm]
{\Large Fall 2026 Syllabus}\\[0.5cm]
{\large Bentley University}
\end{center}

\vspace{0.5cm}

%-----------------------------------------------------------
\section*{Course Description}
%-----------------------------------------------------------

This course introduces the fundamental principles and tools of microeconomics: how individuals and firms make decisions when resources are scarce, and how markets coordinate those decisions. We study demand, supply, and market equilibrium; elasticity; the effects of government policies such as taxes and price controls; the efficiency of markets and the sources of market failure; the costs of production; and firm behavior in different market structures, from perfect competition to monopoly, monopolistic competition, and oligopoly.

The emphasis throughout is on economic reasoning: thinking at the margin, recognizing trade-offs, and using simple graphical and numerical models to make sense of real-world questions in business and policy.

\subsection*{Knowledge}
\begin{itemize}
    \item Scarcity, opportunity cost, trade-offs, and the gains from trade.
    \item Demand, supply, market equilibrium, and how markets respond to shocks.
    \item Elasticity and its role in pricing, revenue, and tax incidence.
    \item Consumer surplus, producer surplus, market efficiency, and the deadweight loss of taxes and price controls.
    \item Externalities, public goods, and the case for government intervention.
    \item Costs of production and firm behavior under perfect competition, monopoly, monopolistic competition, and oligopoly.
\end{itemize}

\subsection*{Skills}
\begin{itemize}
    \item Use supply and demand diagrams to predict the effects of changes in market conditions and policy.
    \item Compute and interpret elasticities, surpluses, costs, and profit-maximizing output and price.
    \item Read and evaluate economic arguments critically and communicate them clearly.
\end{itemize}

\subsection*{Perspectives}
\begin{itemize}
    \item Understand marginal decision making as the core of economic reasoning.
    \item Understand when markets work well, when they fail, and what policy can and cannot fix.
\end{itemize}

%-----------------------------------------------------------
\section*{Class Information}
%-----------------------------------------------------------

\begin{tabularx}{\textwidth}{@{}lX@{}}
\textbf{Instructor:} & Onur Altındağ \\
\textbf{Office:} & RAU 019F \\
\textbf{Email:} & \href{mailto:oaltindag@bentley.edu}{oaltindag@bentley.edu} \\
\textbf{Web:} & \url{https://www.onuraltindag.info/} \\
\textbf{Course page:} & \url{https://www.onuraltindag.info/teaching/ec111/} \\
\textbf{Password:} & \texttt{GoFalcons!} \\[0.3cm]
\textbf{Course:} & \sectionname \\
\textbf{Semester:} & Fall 2026 (August 31, 2026 – December 15, 2026) \\
\textbf{Meeting Time:} & \meetingtime \\
\textbf{Location:} & \classroom \\
\textbf{Delivery Mode:} & In-Person Lecture \\
\end{tabularx}

\vspace{0.3cm}

Office hours are on Mondays and Wednesdays, 11:00 AM -- 12:00 PM and 1:00 -- 3:00 PM, by appointment. Please \href{https://outlook.office.com/bookwithme/user/1c4e4f54e0f0493397e615b2c80a59c7@bentley.edu/meetingtype/rwPThvM0MkK5oGlK_t37nw2?bookingcode=17c9a7b2-69fa-4c48-85ea-ad404642334b&anonymous&ismsaljsauthenabled&ep=mlink}{book a 30-minute slot} to schedule a meeting. If you need to meet urgently or no slots are available, email me directly.

Announcements come by email from me to your Bentley address; check it regularly. Slides, lecture notes, handouts, and mock exams are posted on the course page above after each class. When you email me, put \textbf{EC111} in the subject line.

%-----------------------------------------------------------
\section*{Important Dates and Evaluation}
%-----------------------------------------------------------

\begin{tabularx}{\textwidth}{@{}Xr@{}}
\toprule
\textbf{Component} & \textbf{Weight} \\
\midrule
Weekly assignments + participation & 20\% \\
\textbf{Wednesday, October 7, 2026} – First Midterm & 25\% \\
\textbf{Wednesday, November 11, 2026} – Second Midterm & 25\% \\
\textbf{\finaltime} – Final Exam & 30\% \\
\bottomrule
\end{tabularx}

%-----------------------------------------------------------
\section*{Grading}
%-----------------------------------------------------------

You \textbf{MUST} attend all the midterms and the final as there will be no make-up exams in this course. The midterms are \textbf{NOT} cumulative. If you miss or are likely to miss a midterm due to an emergency, please contact me as soon as possible. You will need to provide supporting documentation/verification of your absence. I will re-weight your final exam accordingly if you have a valid excuse. Please note that family vacations or missing the school shuttle are not valid excuses. If a serious illness, a death in the family, or another difficult circumstance keeps you from an exam or assignment, contact the Office of Student Success (JEN 336, 781.891.2803); they will notify all of your instructors on your behalf.

The final exam \textbf{is} cumulative. If you miss the final exam due to an emergency, you will receive an incomplete for this course. Do not take this class if you know that you will not be able to attend the final exam.

Exams are closed book and closed notes. A basic (non-graphing) calculator is allowed; phones, graphing calculators, and any other device are not, and calculators may not be shared.

You will have \textbf{10 homework assignments} on MindTap, roughly one per week (none in exam weeks), due \textbf{Sunday 11:00 PM}. Your \textbf{two lowest scores are dropped}; the remaining 8 count equally. Each question allows three attempts; an assignment scoring at least \textbf{50\%} receives full credit, otherwise zero. The submission system is automated and will not accept anything after the deadline (even 5 minutes), and there are no exceptions. The questions are algorithmic and slightly different for everyone, so you are encouraged to work together.

\subsection*{Letter grades}
\begin{center}
\begin{tabular}{@{}llllll@{}}
\toprule
A & 94--100 & B+ & 87--89 & C+ & 77--79 \\
A$-$ & 90--93 & B & 83--86 & C & 73--76 \\
 & & B$-$ & 80--82 & C$-$ & 70--72 \\
\midrule
D+ & 67--69 & D & 63--66 & D$-$ & 60--62 \\
F & below 60 & & & & \\
\bottomrule
\end{tabular}
\end{center}
I may lower these cutoffs at the end of the semester; I will not raise them. I do not respond to requests for individual extra credit or for grade adjustments after the fact.

%-----------------------------------------------------------
\section*{Textbook and MindTap}
%-----------------------------------------------------------

The course uses \textbf{MindTap for Mankiw, \textit{Principles of Microeconomics}} (Cengage). The e-book, your homework assignments, and practice questions all live there; a hard copy of the textbook is optional. You \textbf{MUST} register, using your Bentley email, through the link for your section:

\begin{center}
\textbf{\mindtapsection}\\
\url{\mindtaplink}
\end{center}

You have free temporary access until \mindtapaccess; after that you will be prompted to purchase (if your course materials are included in your tuition, no purchase is needed). Registration help: \url{https://startstrong.cengage.com}; technical support: \url{https://support.cengage.com}. Register in the first week; MindTap problems are not an excuse for a missed assignment.

\ifdefined\honors
%-----------------------------------------------------------
\section*{Honors Reading List}
%-----------------------------------------------------------

In addition to the textbook, honors students read and discuss a short set of articles and book chapters over the semester. The list and discussion dates will be posted on the course page in the first weeks of the semester.

\fi
%-----------------------------------------------------------
\section*{Academic Integrity}
%-----------------------------------------------------------

All students are expected to adhere to Bentley's Academic Integrity policy, which includes Bentley's Honor Code (details are in the Undergraduate Student Handbook, the \href{https://catalog.bentley.edu/undergraduate/academic-policies-procedures/}{academic policies section of the catalog}, and the academic integrity course page on Brightspace in which all students and faculty are enrolled). The essence of the policy is that you should not represent someone else's work as your own: no plagiarism, no cheating on exams, no illicit collaboration. Failure to adhere to the policy has serious consequences, including course failure, suspension, or expulsion from the university. The best way to avoid a problem is to consult with me before taking an action that might constitute a violation.

\subsection*{Generative AI}
Nothing in this course is graded on writing, and the homework is algorithmic practice, so the line is simple. You may use generative AI tools (ChatGPT, Claude, Gemini, Copilot, and the like) to study: to explain a concept, to check your reasoning on a practice problem, or to work through a homework question you are stuck on. Treat them as a tutor, not a substitute; the exams are closed book and closed device, and the only way to pass them is to be able to do the work yourself. Any use of AI or any other device during an exam is a violation of the academic integrity policy.

%-----------------------------------------------------------
\section*{Inclusion}
%-----------------------------------------------------------

Everyone in this class has different life experiences and perspectives, and all are valid. As the instructor, I will behave maturely and respectfully in all our class-related engagements, and I expect the same from you. We all adhere to Bentley's standards of conduct, the \href{https://www.bentley.edu/about/mission-and-values}{Bentley Core Values}, in class and in all course work. If you feel that I or anyone in this class has acted outside those values, come and talk to me. If you would rather not, speak with someone in the Office of Student Success (JEN 336, 781.891.2803). My roster shows your name as it appears in Workday; tell me early in the semester if you go by a different name or pronoun.

%-----------------------------------------------------------
\section*{Support and Resources}
%-----------------------------------------------------------

\subsection*{Student Accessibility Services}
Bentley University abides by Section 504 of the Rehabilitation Act of 1973 and the Americans with Disabilities Act of 1990, which stipulate that no student shall be denied the benefits of an education solely by reason of a disability. If you have a hidden or visible disability that may require classroom accommodations, contact \href{https://www.bentley.edu/offices/student-accessibility-services}{Student Accessibility Services} (Office of Student Success, JEN 336, 781.891.2004) within the first four weeks of the semester. Send me your accommodation plan as soon as you have it; accommodations are arranged at least one week before each exam.

\subsection*{The Howard A. Winer '58 Lab for Economics, Accounting and Finance (LEAF)}
The LEAF provides a welcoming and inclusive learning environment where students are encouraged to seek academic support for their Economics, Accounting and Finance courses. Students utilizing the LEAF will find knowledgeable peer tutors ready to help their colleagues to thrive in the Bentley business curriculum. LEAF's free, drop-in tutoring takes place in-person at Morison 220D (rear entrance), opening September 14. For additional information visit \url{https://www.bentley.edu/centers/leaf}. LEAF is a supplement to, not a substitute for, office hours.

\subsection*{Bentley policies and procedures}
All courses at Bentley reflect the university's commitment to a set of core values and practices. You are expected to be familiar with:
\begin{itemize}
    \item \href{https://catalog.bentley.edu/undergraduate/academic-policies-procedures/}{Bentley's honor code and academic integrity system}
    \item \href{https://www.bentley.edu/equity-reporting}{Equity and bias reporting forms and procedures} (Bias Incident Response Team, Office of Institutional Equity)
    \item \href{https://www.bentley.edu/offices/student-accessibility-services}{Student disability accommodations}
    \item \href{https://catalog.bentley.edu/undergraduate/academic-policies-procedures/}{Religious observances policy}
    \item \href{https://www.bentley.edu/about/mission-and-values}{Bentley's core values}
\end{itemize}

\subsection*{Course materials}
Slides, notes, handouts, and exams for this course are for your own use during this semester. They are protected by copyright; do not post or distribute them outside the class without my written consent. The course page is password-protected; do not share the password outside the class.

%-----------------------------------------------------------
\section*{Attendance Policy}
%-----------------------------------------------------------

All students must attend the in-person classes. In-person attendance is essential because it allows for real-time interaction, immediate feedback, and the kind of dialogue that makes learning stick. I do not record lectures, so if you miss a class, you miss the material.

Remote attendance is not an option. This course follows Bentley's \href{https://catalog.bentley.edu/undergraduate/academic-policies-procedures/}{Academic Engagement/Attendance Policy}: students who do not attend or engage during the first week are dropped as no-shows, and two or more consecutive weeks of non-participation can lead to administrative withdrawal. If you must miss class, you are responsible for the material; read the chapter and the posted notes, then bring specific questions to office hours.

\textbf{Weather and closures.} If the university cancels classes, I will email you whether there is an alternative assignment or arrangement for that day.

\textbf{Electronic devices.} Phones away and laptops and tablets closed for the whole class. The work in this course is drawing and reading diagrams, which you learn by doing them by hand, and a screen in your line of sight costs the people around you as much as it costs you. Using a device during class counts against your participation grade. Exceptions are by my permission only; if you have one, you never need to explain it to anyone in the room. Photographing the board at the end of class is fine.

\textbf{Religious observances.} Bentley provides reasonable academic adjustments when a sincerely held religious observance conflicts with a class, exam, or deadline (see the \href{https://downloads.bentley.edu/general/Religious\%20Observances\%202026-2027.pdf}{religious observances calendar} and the \href{https://catalog.bentley.edu/undergraduate/academic-policies-procedures/}{policy}). If you anticipate a conflict, tell me in advance, ideally in the first two weeks of the semester, so we can agree on an adjustment.

%-----------------------------------------------------------
\section*{Tentative Schedule}
%-----------------------------------------------------------

\begin{itemize}
    \item \textbf{Week 1 (Aug 31, Sep 2):} Ten Principles of Economics – Ch. 1 \& Thinking Like an Economist – Ch. 2
    \item \textbf{Week 2 (Sep 9):} Interdependence and the Gains from Trade – Ch. 3 \textit{(no class Sep 7, Labor Day)}
    \item \textbf{Week 3 (Sep 14, 16):} The Market Forces of Supply and Demand – Ch. 4
    \item \textbf{Week 4 (Sep 21, 23):} Elasticity and Its Application – Ch. 5
    \item \textbf{Week 5 (Sep 28, 30):} Supply, Demand, and Government Policies – Ch. 6
    \item \textbf{Week 6 (Oct 5, 7):} Review / \textbf{First Midterm (Oct 7)}
    \item \textbf{Week 7 (Oct 14):} Consumers, Producers, and the Efficiency of Markets – Ch. 7 \textit{(no class Oct 12, Fall Break)}
    \item \textbf{Week 8 (Oct 19, 21):} The Costs of Taxation – Ch. 8 \& International Trade – Ch. 9
    \item \textbf{Week 9 (Oct 26, 28):} Externalities – Ch. 10 \& Public Goods and Common Resources – Ch. 11
    \item \textbf{Week 10 (Nov 2, 4):} The Costs of Production – Ch. 14
    \item \textbf{Week 11 (Nov 9, 11):} Review / \textbf{Second Midterm (Nov 11)}
    \item \textbf{Week 12 (Nov 16, 18):} Firms in Competitive Markets – Ch. 15
    \item \textbf{Week 13 (Nov 23):} Monopoly – Ch. 16 \textit{(no class Nov 25, Thanksgiving Break)}
    \item \textbf{Week 14 (Nov 30, Dec 2):} Monopolistic Competition – Ch. 17 \& Oligopoly – Ch. 18
    \item \textbf{Week 15 (Dec 7):} Review
    \item \textbf{Final Exam:} \finaltime
\end{itemize}

\textit{Note: Chapter references are from Mankiw's \textit{Principles of Microeconomics}, 10th edition (the MindTap e-book). The schedule, and if necessary other parts of this syllabus, may change during the semester; changes are announced in class and by email.}

\end{document}
"
%   xelatex -jobname=EC111_FA2026_Syllabus_20  "\documentclass[11pt,letterpaper]{article}

% Compile with XeLaTeX. Two sections from one source:
%   xelatex -jobname=EC111_FA2026_Syllabus_10H "\def\honors{1}\documentclass[11pt,letterpaper]{article}

% Compile with XeLaTeX. Two sections from one source:
%   xelatex -jobname=EC111_FA2026_Syllabus_10H "\def\honors{1}\documentclass[11pt,letterpaper]{article}

% Compile with XeLaTeX. Two sections from one source:
%   xelatex -jobname=EC111_FA2026_Syllabus_10H "\def\honors{1}\input{EC111_FA2026_Syllabus.tex}"
%   xelatex -jobname=EC111_FA2026_Syllabus_20  "\input{EC111_FA2026_Syllabus.tex}"
\ifdefined\honors
  \newcommand{\sectioncode}{EC 111-10-H}
  \newcommand{\sectionname}{EC 111-10-H Principles of Microeconomics (Honors)}
  \newcommand{\meetingtime}{Monday/Wednesday, 3:30 PM -- 4:50 PM}
  \newcommand{\classroom}{Smith Technology Center 103}
  \newcommand{\finaltime}{Friday, December 11, 2026, 6:00 PM -- 8:00 PM}
  \newcommand{\mindtapsection}{Section 10-H of EC 111 Principles of Microeconomics Fall 2026}
  \newcommand{\mindtaplink}{https://student.cengage.com/course-link/MTPNS57GW6XN}
  \newcommand{\mindtapaccess}{the end of Sunday, September 13, 2026}
\else
  \newcommand{\sectioncode}{EC 111-20}
  \newcommand{\sectionname}{EC 111-20 Principles of Microeconomics}
  \newcommand{\meetingtime}{Monday/Wednesday, 9:30 AM -- 10:50 AM}
  \newcommand{\classroom}{Jennison Hall 412}
  \newcommand{\finaltime}{Friday, December 11, 2026, 8:00 AM -- 10:00 AM}
  \newcommand{\mindtapsection}{Section 20 of EC 111 Principles of Microeconomics Fall 2026}
  \newcommand{\mindtaplink}{https://student.cengage.com/course-link/MTPP8N7GWX7L}
  \newcommand{\mindtapaccess}{the end of Sunday, September 13, 2026}
\fi

% XeLaTeX for Turkish character support
\usepackage{fontspec}
\setmainfont{DejaVuSerif}[Extension=.ttf, BoldFont=DejaVuSerif-Bold, ItalicFont=DejaVuSerif-Italic, BoldItalicFont=DejaVuSerif-BoldItalic]
\setsansfont{DejaVuSans}[Extension=.ttf, BoldFont=DejaVuSans-Bold, ItalicFont=DejaVuSans-Oblique, BoldItalicFont=DejaVuSans-BoldOblique]

% Page layout
\usepackage[margin=1in]{geometry}
\usepackage{setspace}
\onehalfspacing

% Colors and styling
\usepackage[dvipsnames]{xcolor}
\definecolor{bentleyblue}{RGB}{0, 51, 102}
\definecolor{linkblue}{RGB}{0, 102, 204}

% Headers and sections
\usepackage{titlesec}
\titleformat{\section}{\Large\bfseries\color{bentleyblue}}{\thesection}{1em}{}[\titlerule]
\titleformat{\subsection}{\large\bfseries\color{bentleyblue!80}}{\thesubsection}{1em}{}
\titleformat{\subsubsection}{\normalsize\bfseries\color{bentleyblue!60}}{\thesubsubsection}{1em}{}

% Lists
\usepackage{enumitem}
\setlist[itemize]{leftmargin=*, itemsep=2pt, parsep=0pt}
\setlist[enumerate]{leftmargin=*, itemsep=2pt, parsep=0pt}

% Tables
\usepackage{booktabs}
\usepackage{tabularx}

% Links
\usepackage{hyperref}
\hypersetup{
    colorlinks=true,
    linkcolor=linkblue,
    urlcolor=linkblue,
    pdfauthor={Onur Altındağ},
    pdftitle={EC111 Principles of Microeconomics - Fall 2026 Syllabus}
}

% Header/Footer
\usepackage{fancyhdr}
\pagestyle{fancy}
\fancyhf{}
\fancyhead[L]{\small\color{gray}\sectioncode{} Principles of Microeconomics}
\fancyhead[R]{\small\color{gray}Fall 2026}
\fancyfoot[C]{\thepage}
\renewcommand{\headrulewidth}{0.4pt}

% Remove paragraph indent, add space between paragraphs
\setlength{\parindent}{0pt}
\setlength{\parskip}{6pt}

\begin{document}

% Title
\begin{center}
{\LARGE\bfseries\color{bentleyblue} EC 111: Principles of Microeconomics}\\[0.3cm]
{\Large Fall 2026 Syllabus}\\[0.5cm]
{\large Bentley University}
\end{center}

\vspace{0.5cm}

%-----------------------------------------------------------
\section*{Course Description}
%-----------------------------------------------------------

This course introduces the fundamental principles and tools of microeconomics: how individuals and firms make decisions when resources are scarce, and how markets coordinate those decisions. We study demand, supply, and market equilibrium; elasticity; the effects of government policies such as taxes and price controls; the efficiency of markets and the sources of market failure; the costs of production; and firm behavior in different market structures, from perfect competition to monopoly, monopolistic competition, and oligopoly.

The emphasis throughout is on economic reasoning: thinking at the margin, recognizing trade-offs, and using simple graphical and numerical models to make sense of real-world questions in business and policy.

\subsection*{Knowledge}
\begin{itemize}
    \item Scarcity, opportunity cost, trade-offs, and the gains from trade.
    \item Demand, supply, market equilibrium, and how markets respond to shocks.
    \item Elasticity and its role in pricing, revenue, and tax incidence.
    \item Consumer surplus, producer surplus, market efficiency, and the deadweight loss of taxes and price controls.
    \item Externalities, public goods, and the case for government intervention.
    \item Costs of production and firm behavior under perfect competition, monopoly, monopolistic competition, and oligopoly.
\end{itemize}

\subsection*{Skills}
\begin{itemize}
    \item Use supply and demand diagrams to predict the effects of changes in market conditions and policy.
    \item Compute and interpret elasticities, surpluses, costs, and profit-maximizing output and price.
    \item Read and evaluate economic arguments critically and communicate them clearly.
\end{itemize}

\subsection*{Perspectives}
\begin{itemize}
    \item Understand marginal decision making as the core of economic reasoning.
    \item Understand when markets work well, when they fail, and what policy can and cannot fix.
\end{itemize}

%-----------------------------------------------------------
\section*{Class Information}
%-----------------------------------------------------------

\begin{tabularx}{\textwidth}{@{}lX@{}}
\textbf{Instructor:} & Onur Altındağ \\
\textbf{Office:} & RAU 019F \\
\textbf{Email:} & \href{mailto:oaltindag@bentley.edu}{oaltindag@bentley.edu} \\
\textbf{Web:} & \url{https://www.onuraltindag.info/} \\
\textbf{Course page:} & \url{https://www.onuraltindag.info/teaching/ec111/} \\
\textbf{Password:} & \texttt{GoFalcons!} \\[0.3cm]
\textbf{Course:} & \sectionname \\
\textbf{Semester:} & Fall 2026 (August 31, 2026 – December 15, 2026) \\
\textbf{Meeting Time:} & \meetingtime \\
\textbf{Location:} & \classroom \\
\textbf{Delivery Mode:} & In-Person Lecture \\
\end{tabularx}

\vspace{0.3cm}

Office hours are on Mondays and Wednesdays, 11:00 AM -- 12:00 PM and 1:00 -- 3:00 PM, by appointment. Please \href{https://outlook.office.com/bookwithme/user/1c4e4f54e0f0493397e615b2c80a59c7@bentley.edu/meetingtype/rwPThvM0MkK5oGlK_t37nw2?bookingcode=17c9a7b2-69fa-4c48-85ea-ad404642334b&anonymous&ismsaljsauthenabled&ep=mlink}{book a 30-minute slot} to schedule a meeting. If you need to meet urgently or no slots are available, email me directly.

Announcements come by email from me to your Bentley address; check it regularly. Slides, lecture notes, handouts, and mock exams are posted on the course page above after each class. When you email me, put \textbf{EC111} in the subject line.

%-----------------------------------------------------------
\section*{Important Dates and Evaluation}
%-----------------------------------------------------------

\begin{tabularx}{\textwidth}{@{}Xr@{}}
\toprule
\textbf{Component} & \textbf{Weight} \\
\midrule
Weekly assignments + participation & 20\% \\
\textbf{Wednesday, October 7, 2026} – First Midterm & 25\% \\
\textbf{Wednesday, November 11, 2026} – Second Midterm & 25\% \\
\textbf{\finaltime} – Final Exam & 30\% \\
\bottomrule
\end{tabularx}

%-----------------------------------------------------------
\section*{Grading}
%-----------------------------------------------------------

You \textbf{MUST} attend all the midterms and the final as there will be no make-up exams in this course. The midterms are \textbf{NOT} cumulative. If you miss or are likely to miss a midterm due to an emergency, please contact me as soon as possible. You will need to provide supporting documentation/verification of your absence. I will re-weight your final exam accordingly if you have a valid excuse. Please note that family vacations or missing the school shuttle are not valid excuses. If a serious illness, a death in the family, or another difficult circumstance keeps you from an exam or assignment, contact the Office of Student Success (JEN 336, 781.891.2803); they will notify all of your instructors on your behalf.

The final exam \textbf{is} cumulative. If you miss the final exam due to an emergency, you will receive an incomplete for this course. Do not take this class if you know that you will not be able to attend the final exam.

Exams are closed book and closed notes. A basic (non-graphing) calculator is allowed; phones, graphing calculators, and any other device are not, and calculators may not be shared.

You will have \textbf{10 homework assignments} on MindTap, roughly one per week (none in exam weeks), due \textbf{Sunday 11:00 PM}. Your \textbf{two lowest scores are dropped}; the remaining 8 count equally. Each question allows three attempts; an assignment scoring at least \textbf{50\%} receives full credit, otherwise zero. The submission system is automated and will not accept anything after the deadline (even 5 minutes), and there are no exceptions. The questions are algorithmic and slightly different for everyone, so you are encouraged to work together.

\subsection*{Letter grades}
\begin{center}
\begin{tabular}{@{}llllll@{}}
\toprule
A & 94--100 & B+ & 87--89 & C+ & 77--79 \\
A$-$ & 90--93 & B & 83--86 & C & 73--76 \\
 & & B$-$ & 80--82 & C$-$ & 70--72 \\
\midrule
D+ & 67--69 & D & 63--66 & D$-$ & 60--62 \\
F & below 60 & & & & \\
\bottomrule
\end{tabular}
\end{center}
I may lower these cutoffs at the end of the semester; I will not raise them. I do not respond to requests for individual extra credit or for grade adjustments after the fact.

%-----------------------------------------------------------
\section*{Textbook and MindTap}
%-----------------------------------------------------------

The course uses \textbf{MindTap for Mankiw, \textit{Principles of Microeconomics}} (Cengage). The e-book, your homework assignments, and practice questions all live there; a hard copy of the textbook is optional. You \textbf{MUST} register, using your Bentley email, through the link for your section:

\begin{center}
\textbf{\mindtapsection}\\
\url{\mindtaplink}
\end{center}

You have free temporary access until \mindtapaccess; after that you will be prompted to purchase (if your course materials are included in your tuition, no purchase is needed). Registration help: \url{https://startstrong.cengage.com}; technical support: \url{https://support.cengage.com}. Register in the first week; MindTap problems are not an excuse for a missed assignment.

\ifdefined\honors
%-----------------------------------------------------------
\section*{Honors Reading List}
%-----------------------------------------------------------

In addition to the textbook, honors students read and discuss a short set of articles and book chapters over the semester. The list and discussion dates will be posted on the course page in the first weeks of the semester.

\fi
%-----------------------------------------------------------
\section*{Academic Integrity}
%-----------------------------------------------------------

All students are expected to adhere to Bentley's Academic Integrity policy, which includes Bentley's Honor Code (details are in the Undergraduate Student Handbook, the \href{https://catalog.bentley.edu/undergraduate/academic-policies-procedures/}{academic policies section of the catalog}, and the academic integrity course page on Brightspace in which all students and faculty are enrolled). The essence of the policy is that you should not represent someone else's work as your own: no plagiarism, no cheating on exams, no illicit collaboration. Failure to adhere to the policy has serious consequences, including course failure, suspension, or expulsion from the university. The best way to avoid a problem is to consult with me before taking an action that might constitute a violation.

\subsection*{Generative AI}
Nothing in this course is graded on writing, and the homework is algorithmic practice, so the line is simple. You may use generative AI tools (ChatGPT, Claude, Gemini, Copilot, and the like) to study: to explain a concept, to check your reasoning on a practice problem, or to work through a homework question you are stuck on. Treat them as a tutor, not a substitute; the exams are closed book and closed device, and the only way to pass them is to be able to do the work yourself. Any use of AI or any other device during an exam is a violation of the academic integrity policy.

%-----------------------------------------------------------
\section*{Inclusion}
%-----------------------------------------------------------

Everyone in this class has different life experiences and perspectives, and all are valid. As the instructor, I will behave maturely and respectfully in all our class-related engagements, and I expect the same from you. We all adhere to Bentley's standards of conduct, the \href{https://www.bentley.edu/about/mission-and-values}{Bentley Core Values}, in class and in all course work. If you feel that I or anyone in this class has acted outside those values, come and talk to me. If you would rather not, speak with someone in the Office of Student Success (JEN 336, 781.891.2803). My roster shows your name as it appears in Workday; tell me early in the semester if you go by a different name or pronoun.

%-----------------------------------------------------------
\section*{Support and Resources}
%-----------------------------------------------------------

\subsection*{Student Accessibility Services}
Bentley University abides by Section 504 of the Rehabilitation Act of 1973 and the Americans with Disabilities Act of 1990, which stipulate that no student shall be denied the benefits of an education solely by reason of a disability. If you have a hidden or visible disability that may require classroom accommodations, contact \href{https://www.bentley.edu/offices/student-accessibility-services}{Student Accessibility Services} (Office of Student Success, JEN 336, 781.891.2004) within the first four weeks of the semester. Send me your accommodation plan as soon as you have it; accommodations are arranged at least one week before each exam.

\subsection*{The Howard A. Winer '58 Lab for Economics, Accounting and Finance (LEAF)}
The LEAF provides a welcoming and inclusive learning environment where students are encouraged to seek academic support for their Economics, Accounting and Finance courses. Students utilizing the LEAF will find knowledgeable peer tutors ready to help their colleagues to thrive in the Bentley business curriculum. LEAF's free, drop-in tutoring takes place in-person at Morison 220D (rear entrance), opening September 14. For additional information visit \url{https://www.bentley.edu/centers/leaf}. LEAF is a supplement to, not a substitute for, office hours.

\subsection*{Bentley policies and procedures}
All courses at Bentley reflect the university's commitment to a set of core values and practices. You are expected to be familiar with:
\begin{itemize}
    \item \href{https://catalog.bentley.edu/undergraduate/academic-policies-procedures/}{Bentley's honor code and academic integrity system}
    \item \href{https://www.bentley.edu/equity-reporting}{Equity and bias reporting forms and procedures} (Bias Incident Response Team, Office of Institutional Equity)
    \item \href{https://www.bentley.edu/offices/student-accessibility-services}{Student disability accommodations}
    \item \href{https://catalog.bentley.edu/undergraduate/academic-policies-procedures/}{Religious observances policy}
    \item \href{https://www.bentley.edu/about/mission-and-values}{Bentley's core values}
\end{itemize}

\subsection*{Course materials}
Slides, notes, handouts, and exams for this course are for your own use during this semester. They are protected by copyright; do not post or distribute them outside the class without my written consent. The course page is password-protected; do not share the password outside the class.

%-----------------------------------------------------------
\section*{Attendance Policy}
%-----------------------------------------------------------

All students must attend the in-person classes. In-person attendance is essential because it allows for real-time interaction, immediate feedback, and the kind of dialogue that makes learning stick. I do not record lectures, so if you miss a class, you miss the material.

Remote attendance is not an option. This course follows Bentley's \href{https://catalog.bentley.edu/undergraduate/academic-policies-procedures/}{Academic Engagement/Attendance Policy}: students who do not attend or engage during the first week are dropped as no-shows, and two or more consecutive weeks of non-participation can lead to administrative withdrawal. If you must miss class, you are responsible for the material; read the chapter and the posted notes, then bring specific questions to office hours.

\textbf{Weather and closures.} If the university cancels classes, I will email you whether there is an alternative assignment or arrangement for that day.

\textbf{Electronic devices.} Phones away and laptops and tablets closed for the whole class. The work in this course is drawing and reading diagrams, which you learn by doing them by hand, and a screen in your line of sight costs the people around you as much as it costs you. Using a device during class counts against your participation grade. Exceptions are by my permission only; if you have one, you never need to explain it to anyone in the room. Photographing the board at the end of class is fine.

\textbf{Religious observances.} Bentley provides reasonable academic adjustments when a sincerely held religious observance conflicts with a class, exam, or deadline (see the \href{https://downloads.bentley.edu/general/Religious\%20Observances\%202026-2027.pdf}{religious observances calendar} and the \href{https://catalog.bentley.edu/undergraduate/academic-policies-procedures/}{policy}). If you anticipate a conflict, tell me in advance, ideally in the first two weeks of the semester, so we can agree on an adjustment.

%-----------------------------------------------------------
\section*{Tentative Schedule}
%-----------------------------------------------------------

\begin{itemize}
    \item \textbf{Week 1 (Aug 31, Sep 2):} Ten Principles of Economics – Ch. 1 \& Thinking Like an Economist – Ch. 2
    \item \textbf{Week 2 (Sep 9):} Interdependence and the Gains from Trade – Ch. 3 \textit{(no class Sep 7, Labor Day)}
    \item \textbf{Week 3 (Sep 14, 16):} The Market Forces of Supply and Demand – Ch. 4
    \item \textbf{Week 4 (Sep 21, 23):} Elasticity and Its Application – Ch. 5
    \item \textbf{Week 5 (Sep 28, 30):} Supply, Demand, and Government Policies – Ch. 6
    \item \textbf{Week 6 (Oct 5, 7):} Review / \textbf{First Midterm (Oct 7)}
    \item \textbf{Week 7 (Oct 14):} Consumers, Producers, and the Efficiency of Markets – Ch. 7 \textit{(no class Oct 12, Fall Break)}
    \item \textbf{Week 8 (Oct 19, 21):} The Costs of Taxation – Ch. 8 \& International Trade – Ch. 9
    \item \textbf{Week 9 (Oct 26, 28):} Externalities – Ch. 10 \& Public Goods and Common Resources – Ch. 11
    \item \textbf{Week 10 (Nov 2, 4):} The Costs of Production – Ch. 14
    \item \textbf{Week 11 (Nov 9, 11):} Review / \textbf{Second Midterm (Nov 11)}
    \item \textbf{Week 12 (Nov 16, 18):} Firms in Competitive Markets – Ch. 15
    \item \textbf{Week 13 (Nov 23):} Monopoly – Ch. 16 \textit{(no class Nov 25, Thanksgiving Break)}
    \item \textbf{Week 14 (Nov 30, Dec 2):} Monopolistic Competition – Ch. 17 \& Oligopoly – Ch. 18
    \item \textbf{Week 15 (Dec 7):} Review
    \item \textbf{Final Exam:} \finaltime
\end{itemize}

\textit{Note: Chapter references are from Mankiw's \textit{Principles of Microeconomics}, 10th edition (the MindTap e-book). The schedule, and if necessary other parts of this syllabus, may change during the semester; changes are announced in class and by email.}

\end{document}
"
%   xelatex -jobname=EC111_FA2026_Syllabus_20  "\documentclass[11pt,letterpaper]{article}

% Compile with XeLaTeX. Two sections from one source:
%   xelatex -jobname=EC111_FA2026_Syllabus_10H "\def\honors{1}\input{EC111_FA2026_Syllabus.tex}"
%   xelatex -jobname=EC111_FA2026_Syllabus_20  "\input{EC111_FA2026_Syllabus.tex}"
\ifdefined\honors
  \newcommand{\sectioncode}{EC 111-10-H}
  \newcommand{\sectionname}{EC 111-10-H Principles of Microeconomics (Honors)}
  \newcommand{\meetingtime}{Monday/Wednesday, 3:30 PM -- 4:50 PM}
  \newcommand{\classroom}{Smith Technology Center 103}
  \newcommand{\finaltime}{Friday, December 11, 2026, 6:00 PM -- 8:00 PM}
  \newcommand{\mindtapsection}{Section 10-H of EC 111 Principles of Microeconomics Fall 2026}
  \newcommand{\mindtaplink}{https://student.cengage.com/course-link/MTPNS57GW6XN}
  \newcommand{\mindtapaccess}{the end of Sunday, September 13, 2026}
\else
  \newcommand{\sectioncode}{EC 111-20}
  \newcommand{\sectionname}{EC 111-20 Principles of Microeconomics}
  \newcommand{\meetingtime}{Monday/Wednesday, 9:30 AM -- 10:50 AM}
  \newcommand{\classroom}{Jennison Hall 412}
  \newcommand{\finaltime}{Friday, December 11, 2026, 8:00 AM -- 10:00 AM}
  \newcommand{\mindtapsection}{Section 20 of EC 111 Principles of Microeconomics Fall 2026}
  \newcommand{\mindtaplink}{https://student.cengage.com/course-link/MTPP8N7GWX7L}
  \newcommand{\mindtapaccess}{the end of Sunday, September 13, 2026}
\fi

% XeLaTeX for Turkish character support
\usepackage{fontspec}
\setmainfont{DejaVuSerif}[Extension=.ttf, BoldFont=DejaVuSerif-Bold, ItalicFont=DejaVuSerif-Italic, BoldItalicFont=DejaVuSerif-BoldItalic]
\setsansfont{DejaVuSans}[Extension=.ttf, BoldFont=DejaVuSans-Bold, ItalicFont=DejaVuSans-Oblique, BoldItalicFont=DejaVuSans-BoldOblique]

% Page layout
\usepackage[margin=1in]{geometry}
\usepackage{setspace}
\onehalfspacing

% Colors and styling
\usepackage[dvipsnames]{xcolor}
\definecolor{bentleyblue}{RGB}{0, 51, 102}
\definecolor{linkblue}{RGB}{0, 102, 204}

% Headers and sections
\usepackage{titlesec}
\titleformat{\section}{\Large\bfseries\color{bentleyblue}}{\thesection}{1em}{}[\titlerule]
\titleformat{\subsection}{\large\bfseries\color{bentleyblue!80}}{\thesubsection}{1em}{}
\titleformat{\subsubsection}{\normalsize\bfseries\color{bentleyblue!60}}{\thesubsubsection}{1em}{}

% Lists
\usepackage{enumitem}
\setlist[itemize]{leftmargin=*, itemsep=2pt, parsep=0pt}
\setlist[enumerate]{leftmargin=*, itemsep=2pt, parsep=0pt}

% Tables
\usepackage{booktabs}
\usepackage{tabularx}

% Links
\usepackage{hyperref}
\hypersetup{
    colorlinks=true,
    linkcolor=linkblue,
    urlcolor=linkblue,
    pdfauthor={Onur Altındağ},
    pdftitle={EC111 Principles of Microeconomics - Fall 2026 Syllabus}
}

% Header/Footer
\usepackage{fancyhdr}
\pagestyle{fancy}
\fancyhf{}
\fancyhead[L]{\small\color{gray}\sectioncode{} Principles of Microeconomics}
\fancyhead[R]{\small\color{gray}Fall 2026}
\fancyfoot[C]{\thepage}
\renewcommand{\headrulewidth}{0.4pt}

% Remove paragraph indent, add space between paragraphs
\setlength{\parindent}{0pt}
\setlength{\parskip}{6pt}

\begin{document}

% Title
\begin{center}
{\LARGE\bfseries\color{bentleyblue} EC 111: Principles of Microeconomics}\\[0.3cm]
{\Large Fall 2026 Syllabus}\\[0.5cm]
{\large Bentley University}
\end{center}

\vspace{0.5cm}

%-----------------------------------------------------------
\section*{Course Description}
%-----------------------------------------------------------

This course introduces the fundamental principles and tools of microeconomics: how individuals and firms make decisions when resources are scarce, and how markets coordinate those decisions. We study demand, supply, and market equilibrium; elasticity; the effects of government policies such as taxes and price controls; the efficiency of markets and the sources of market failure; the costs of production; and firm behavior in different market structures, from perfect competition to monopoly, monopolistic competition, and oligopoly.

The emphasis throughout is on economic reasoning: thinking at the margin, recognizing trade-offs, and using simple graphical and numerical models to make sense of real-world questions in business and policy.

\subsection*{Knowledge}
\begin{itemize}
    \item Scarcity, opportunity cost, trade-offs, and the gains from trade.
    \item Demand, supply, market equilibrium, and how markets respond to shocks.
    \item Elasticity and its role in pricing, revenue, and tax incidence.
    \item Consumer surplus, producer surplus, market efficiency, and the deadweight loss of taxes and price controls.
    \item Externalities, public goods, and the case for government intervention.
    \item Costs of production and firm behavior under perfect competition, monopoly, monopolistic competition, and oligopoly.
\end{itemize}

\subsection*{Skills}
\begin{itemize}
    \item Use supply and demand diagrams to predict the effects of changes in market conditions and policy.
    \item Compute and interpret elasticities, surpluses, costs, and profit-maximizing output and price.
    \item Read and evaluate economic arguments critically and communicate them clearly.
\end{itemize}

\subsection*{Perspectives}
\begin{itemize}
    \item Understand marginal decision making as the core of economic reasoning.
    \item Understand when markets work well, when they fail, and what policy can and cannot fix.
\end{itemize}

%-----------------------------------------------------------
\section*{Class Information}
%-----------------------------------------------------------

\begin{tabularx}{\textwidth}{@{}lX@{}}
\textbf{Instructor:} & Onur Altındağ \\
\textbf{Office:} & RAU 019F \\
\textbf{Email:} & \href{mailto:oaltindag@bentley.edu}{oaltindag@bentley.edu} \\
\textbf{Web:} & \url{https://www.onuraltindag.info/} \\
\textbf{Course page:} & \url{https://www.onuraltindag.info/teaching/ec111/} \\
\textbf{Password:} & \texttt{GoFalcons!} \\[0.3cm]
\textbf{Course:} & \sectionname \\
\textbf{Semester:} & Fall 2026 (August 31, 2026 – December 15, 2026) \\
\textbf{Meeting Time:} & \meetingtime \\
\textbf{Location:} & \classroom \\
\textbf{Delivery Mode:} & In-Person Lecture \\
\end{tabularx}

\vspace{0.3cm}

Office hours are on Mondays and Wednesdays, 11:00 AM -- 12:00 PM and 1:00 -- 3:00 PM, by appointment. Please \href{https://outlook.office.com/bookwithme/user/1c4e4f54e0f0493397e615b2c80a59c7@bentley.edu/meetingtype/rwPThvM0MkK5oGlK_t37nw2?bookingcode=17c9a7b2-69fa-4c48-85ea-ad404642334b&anonymous&ismsaljsauthenabled&ep=mlink}{book a 30-minute slot} to schedule a meeting. If you need to meet urgently or no slots are available, email me directly.

Announcements come by email from me to your Bentley address; check it regularly. Slides, lecture notes, handouts, and mock exams are posted on the course page above after each class. When you email me, put \textbf{EC111} in the subject line.

%-----------------------------------------------------------
\section*{Important Dates and Evaluation}
%-----------------------------------------------------------

\begin{tabularx}{\textwidth}{@{}Xr@{}}
\toprule
\textbf{Component} & \textbf{Weight} \\
\midrule
Weekly assignments + participation & 20\% \\
\textbf{Wednesday, October 7, 2026} – First Midterm & 25\% \\
\textbf{Wednesday, November 11, 2026} – Second Midterm & 25\% \\
\textbf{\finaltime} – Final Exam & 30\% \\
\bottomrule
\end{tabularx}

%-----------------------------------------------------------
\section*{Grading}
%-----------------------------------------------------------

You \textbf{MUST} attend all the midterms and the final as there will be no make-up exams in this course. The midterms are \textbf{NOT} cumulative. If you miss or are likely to miss a midterm due to an emergency, please contact me as soon as possible. You will need to provide supporting documentation/verification of your absence. I will re-weight your final exam accordingly if you have a valid excuse. Please note that family vacations or missing the school shuttle are not valid excuses. If a serious illness, a death in the family, or another difficult circumstance keeps you from an exam or assignment, contact the Office of Student Success (JEN 336, 781.891.2803); they will notify all of your instructors on your behalf.

The final exam \textbf{is} cumulative. If you miss the final exam due to an emergency, you will receive an incomplete for this course. Do not take this class if you know that you will not be able to attend the final exam.

Exams are closed book and closed notes. A basic (non-graphing) calculator is allowed; phones, graphing calculators, and any other device are not, and calculators may not be shared.

You will have \textbf{10 homework assignments} on MindTap, roughly one per week (none in exam weeks), due \textbf{Sunday 11:00 PM}. Your \textbf{two lowest scores are dropped}; the remaining 8 count equally. Each question allows three attempts; an assignment scoring at least \textbf{50\%} receives full credit, otherwise zero. The submission system is automated and will not accept anything after the deadline (even 5 minutes), and there are no exceptions. The questions are algorithmic and slightly different for everyone, so you are encouraged to work together.

\subsection*{Letter grades}
\begin{center}
\begin{tabular}{@{}llllll@{}}
\toprule
A & 94--100 & B+ & 87--89 & C+ & 77--79 \\
A$-$ & 90--93 & B & 83--86 & C & 73--76 \\
 & & B$-$ & 80--82 & C$-$ & 70--72 \\
\midrule
D+ & 67--69 & D & 63--66 & D$-$ & 60--62 \\
F & below 60 & & & & \\
\bottomrule
\end{tabular}
\end{center}
I may lower these cutoffs at the end of the semester; I will not raise them. I do not respond to requests for individual extra credit or for grade adjustments after the fact.

%-----------------------------------------------------------
\section*{Textbook and MindTap}
%-----------------------------------------------------------

The course uses \textbf{MindTap for Mankiw, \textit{Principles of Microeconomics}} (Cengage). The e-book, your homework assignments, and practice questions all live there; a hard copy of the textbook is optional. You \textbf{MUST} register, using your Bentley email, through the link for your section:

\begin{center}
\textbf{\mindtapsection}\\
\url{\mindtaplink}
\end{center}

You have free temporary access until \mindtapaccess; after that you will be prompted to purchase (if your course materials are included in your tuition, no purchase is needed). Registration help: \url{https://startstrong.cengage.com}; technical support: \url{https://support.cengage.com}. Register in the first week; MindTap problems are not an excuse for a missed assignment.

\ifdefined\honors
%-----------------------------------------------------------
\section*{Honors Reading List}
%-----------------------------------------------------------

In addition to the textbook, honors students read and discuss a short set of articles and book chapters over the semester. The list and discussion dates will be posted on the course page in the first weeks of the semester.

\fi
%-----------------------------------------------------------
\section*{Academic Integrity}
%-----------------------------------------------------------

All students are expected to adhere to Bentley's Academic Integrity policy, which includes Bentley's Honor Code (details are in the Undergraduate Student Handbook, the \href{https://catalog.bentley.edu/undergraduate/academic-policies-procedures/}{academic policies section of the catalog}, and the academic integrity course page on Brightspace in which all students and faculty are enrolled). The essence of the policy is that you should not represent someone else's work as your own: no plagiarism, no cheating on exams, no illicit collaboration. Failure to adhere to the policy has serious consequences, including course failure, suspension, or expulsion from the university. The best way to avoid a problem is to consult with me before taking an action that might constitute a violation.

\subsection*{Generative AI}
Nothing in this course is graded on writing, and the homework is algorithmic practice, so the line is simple. You may use generative AI tools (ChatGPT, Claude, Gemini, Copilot, and the like) to study: to explain a concept, to check your reasoning on a practice problem, or to work through a homework question you are stuck on. Treat them as a tutor, not a substitute; the exams are closed book and closed device, and the only way to pass them is to be able to do the work yourself. Any use of AI or any other device during an exam is a violation of the academic integrity policy.

%-----------------------------------------------------------
\section*{Inclusion}
%-----------------------------------------------------------

Everyone in this class has different life experiences and perspectives, and all are valid. As the instructor, I will behave maturely and respectfully in all our class-related engagements, and I expect the same from you. We all adhere to Bentley's standards of conduct, the \href{https://www.bentley.edu/about/mission-and-values}{Bentley Core Values}, in class and in all course work. If you feel that I or anyone in this class has acted outside those values, come and talk to me. If you would rather not, speak with someone in the Office of Student Success (JEN 336, 781.891.2803). My roster shows your name as it appears in Workday; tell me early in the semester if you go by a different name or pronoun.

%-----------------------------------------------------------
\section*{Support and Resources}
%-----------------------------------------------------------

\subsection*{Student Accessibility Services}
Bentley University abides by Section 504 of the Rehabilitation Act of 1973 and the Americans with Disabilities Act of 1990, which stipulate that no student shall be denied the benefits of an education solely by reason of a disability. If you have a hidden or visible disability that may require classroom accommodations, contact \href{https://www.bentley.edu/offices/student-accessibility-services}{Student Accessibility Services} (Office of Student Success, JEN 336, 781.891.2004) within the first four weeks of the semester. Send me your accommodation plan as soon as you have it; accommodations are arranged at least one week before each exam.

\subsection*{The Howard A. Winer '58 Lab for Economics, Accounting and Finance (LEAF)}
The LEAF provides a welcoming and inclusive learning environment where students are encouraged to seek academic support for their Economics, Accounting and Finance courses. Students utilizing the LEAF will find knowledgeable peer tutors ready to help their colleagues to thrive in the Bentley business curriculum. LEAF's free, drop-in tutoring takes place in-person at Morison 220D (rear entrance), opening September 14. For additional information visit \url{https://www.bentley.edu/centers/leaf}. LEAF is a supplement to, not a substitute for, office hours.

\subsection*{Bentley policies and procedures}
All courses at Bentley reflect the university's commitment to a set of core values and practices. You are expected to be familiar with:
\begin{itemize}
    \item \href{https://catalog.bentley.edu/undergraduate/academic-policies-procedures/}{Bentley's honor code and academic integrity system}
    \item \href{https://www.bentley.edu/equity-reporting}{Equity and bias reporting forms and procedures} (Bias Incident Response Team, Office of Institutional Equity)
    \item \href{https://www.bentley.edu/offices/student-accessibility-services}{Student disability accommodations}
    \item \href{https://catalog.bentley.edu/undergraduate/academic-policies-procedures/}{Religious observances policy}
    \item \href{https://www.bentley.edu/about/mission-and-values}{Bentley's core values}
\end{itemize}

\subsection*{Course materials}
Slides, notes, handouts, and exams for this course are for your own use during this semester. They are protected by copyright; do not post or distribute them outside the class without my written consent. The course page is password-protected; do not share the password outside the class.

%-----------------------------------------------------------
\section*{Attendance Policy}
%-----------------------------------------------------------

All students must attend the in-person classes. In-person attendance is essential because it allows for real-time interaction, immediate feedback, and the kind of dialogue that makes learning stick. I do not record lectures, so if you miss a class, you miss the material.

Remote attendance is not an option. This course follows Bentley's \href{https://catalog.bentley.edu/undergraduate/academic-policies-procedures/}{Academic Engagement/Attendance Policy}: students who do not attend or engage during the first week are dropped as no-shows, and two or more consecutive weeks of non-participation can lead to administrative withdrawal. If you must miss class, you are responsible for the material; read the chapter and the posted notes, then bring specific questions to office hours.

\textbf{Weather and closures.} If the university cancels classes, I will email you whether there is an alternative assignment or arrangement for that day.

\textbf{Electronic devices.} Phones away and laptops and tablets closed for the whole class. The work in this course is drawing and reading diagrams, which you learn by doing them by hand, and a screen in your line of sight costs the people around you as much as it costs you. Using a device during class counts against your participation grade. Exceptions are by my permission only; if you have one, you never need to explain it to anyone in the room. Photographing the board at the end of class is fine.

\textbf{Religious observances.} Bentley provides reasonable academic adjustments when a sincerely held religious observance conflicts with a class, exam, or deadline (see the \href{https://downloads.bentley.edu/general/Religious\%20Observances\%202026-2027.pdf}{religious observances calendar} and the \href{https://catalog.bentley.edu/undergraduate/academic-policies-procedures/}{policy}). If you anticipate a conflict, tell me in advance, ideally in the first two weeks of the semester, so we can agree on an adjustment.

%-----------------------------------------------------------
\section*{Tentative Schedule}
%-----------------------------------------------------------

\begin{itemize}
    \item \textbf{Week 1 (Aug 31, Sep 2):} Ten Principles of Economics – Ch. 1 \& Thinking Like an Economist – Ch. 2
    \item \textbf{Week 2 (Sep 9):} Interdependence and the Gains from Trade – Ch. 3 \textit{(no class Sep 7, Labor Day)}
    \item \textbf{Week 3 (Sep 14, 16):} The Market Forces of Supply and Demand – Ch. 4
    \item \textbf{Week 4 (Sep 21, 23):} Elasticity and Its Application – Ch. 5
    \item \textbf{Week 5 (Sep 28, 30):} Supply, Demand, and Government Policies – Ch. 6
    \item \textbf{Week 6 (Oct 5, 7):} Review / \textbf{First Midterm (Oct 7)}
    \item \textbf{Week 7 (Oct 14):} Consumers, Producers, and the Efficiency of Markets – Ch. 7 \textit{(no class Oct 12, Fall Break)}
    \item \textbf{Week 8 (Oct 19, 21):} The Costs of Taxation – Ch. 8 \& International Trade – Ch. 9
    \item \textbf{Week 9 (Oct 26, 28):} Externalities – Ch. 10 \& Public Goods and Common Resources – Ch. 11
    \item \textbf{Week 10 (Nov 2, 4):} The Costs of Production – Ch. 14
    \item \textbf{Week 11 (Nov 9, 11):} Review / \textbf{Second Midterm (Nov 11)}
    \item \textbf{Week 12 (Nov 16, 18):} Firms in Competitive Markets – Ch. 15
    \item \textbf{Week 13 (Nov 23):} Monopoly – Ch. 16 \textit{(no class Nov 25, Thanksgiving Break)}
    \item \textbf{Week 14 (Nov 30, Dec 2):} Monopolistic Competition – Ch. 17 \& Oligopoly – Ch. 18
    \item \textbf{Week 15 (Dec 7):} Review
    \item \textbf{Final Exam:} \finaltime
\end{itemize}

\textit{Note: Chapter references are from Mankiw's \textit{Principles of Microeconomics}, 10th edition (the MindTap e-book). The schedule, and if necessary other parts of this syllabus, may change during the semester; changes are announced in class and by email.}

\end{document}
"
\ifdefined\honors
  \newcommand{\sectioncode}{EC 111-10-H}
  \newcommand{\sectionname}{EC 111-10-H Principles of Microeconomics (Honors)}
  \newcommand{\meetingtime}{Monday/Wednesday, 3:30 PM -- 4:50 PM}
  \newcommand{\classroom}{Smith Technology Center 103}
  \newcommand{\finaltime}{Friday, December 11, 2026, 6:00 PM -- 8:00 PM}
  \newcommand{\mindtapsection}{Section 10-H of EC 111 Principles of Microeconomics Fall 2026}
  \newcommand{\mindtaplink}{https://student.cengage.com/course-link/MTPNS57GW6XN}
  \newcommand{\mindtapaccess}{the end of Sunday, September 13, 2026}
\else
  \newcommand{\sectioncode}{EC 111-20}
  \newcommand{\sectionname}{EC 111-20 Principles of Microeconomics}
  \newcommand{\meetingtime}{Monday/Wednesday, 9:30 AM -- 10:50 AM}
  \newcommand{\classroom}{Jennison Hall 412}
  \newcommand{\finaltime}{Friday, December 11, 2026, 8:00 AM -- 10:00 AM}
  \newcommand{\mindtapsection}{Section 20 of EC 111 Principles of Microeconomics Fall 2026}
  \newcommand{\mindtaplink}{https://student.cengage.com/course-link/MTPP8N7GWX7L}
  \newcommand{\mindtapaccess}{the end of Sunday, September 13, 2026}
\fi

% XeLaTeX for Turkish character support
\usepackage{fontspec}
\setmainfont{DejaVuSerif}[Extension=.ttf, BoldFont=DejaVuSerif-Bold, ItalicFont=DejaVuSerif-Italic, BoldItalicFont=DejaVuSerif-BoldItalic]
\setsansfont{DejaVuSans}[Extension=.ttf, BoldFont=DejaVuSans-Bold, ItalicFont=DejaVuSans-Oblique, BoldItalicFont=DejaVuSans-BoldOblique]

% Page layout
\usepackage[margin=1in]{geometry}
\usepackage{setspace}
\onehalfspacing

% Colors and styling
\usepackage[dvipsnames]{xcolor}
\definecolor{bentleyblue}{RGB}{0, 51, 102}
\definecolor{linkblue}{RGB}{0, 102, 204}

% Headers and sections
\usepackage{titlesec}
\titleformat{\section}{\Large\bfseries\color{bentleyblue}}{\thesection}{1em}{}[\titlerule]
\titleformat{\subsection}{\large\bfseries\color{bentleyblue!80}}{\thesubsection}{1em}{}
\titleformat{\subsubsection}{\normalsize\bfseries\color{bentleyblue!60}}{\thesubsubsection}{1em}{}

% Lists
\usepackage{enumitem}
\setlist[itemize]{leftmargin=*, itemsep=2pt, parsep=0pt}
\setlist[enumerate]{leftmargin=*, itemsep=2pt, parsep=0pt}

% Tables
\usepackage{booktabs}
\usepackage{tabularx}

% Links
\usepackage{hyperref}
\hypersetup{
    colorlinks=true,
    linkcolor=linkblue,
    urlcolor=linkblue,
    pdfauthor={Onur Altındağ},
    pdftitle={EC111 Principles of Microeconomics - Fall 2026 Syllabus}
}

% Header/Footer
\usepackage{fancyhdr}
\pagestyle{fancy}
\fancyhf{}
\fancyhead[L]{\small\color{gray}\sectioncode{} Principles of Microeconomics}
\fancyhead[R]{\small\color{gray}Fall 2026}
\fancyfoot[C]{\thepage}
\renewcommand{\headrulewidth}{0.4pt}

% Remove paragraph indent, add space between paragraphs
\setlength{\parindent}{0pt}
\setlength{\parskip}{6pt}

\begin{document}

% Title
\begin{center}
{\LARGE\bfseries\color{bentleyblue} EC 111: Principles of Microeconomics}\\[0.3cm]
{\Large Fall 2026 Syllabus}\\[0.5cm]
{\large Bentley University}
\end{center}

\vspace{0.5cm}

%-----------------------------------------------------------
\section*{Course Description}
%-----------------------------------------------------------

This course introduces the fundamental principles and tools of microeconomics: how individuals and firms make decisions when resources are scarce, and how markets coordinate those decisions. We study demand, supply, and market equilibrium; elasticity; the effects of government policies such as taxes and price controls; the efficiency of markets and the sources of market failure; the costs of production; and firm behavior in different market structures, from perfect competition to monopoly, monopolistic competition, and oligopoly.

The emphasis throughout is on economic reasoning: thinking at the margin, recognizing trade-offs, and using simple graphical and numerical models to make sense of real-world questions in business and policy.

\subsection*{Knowledge}
\begin{itemize}
    \item Scarcity, opportunity cost, trade-offs, and the gains from trade.
    \item Demand, supply, market equilibrium, and how markets respond to shocks.
    \item Elasticity and its role in pricing, revenue, and tax incidence.
    \item Consumer surplus, producer surplus, market efficiency, and the deadweight loss of taxes and price controls.
    \item Externalities, public goods, and the case for government intervention.
    \item Costs of production and firm behavior under perfect competition, monopoly, monopolistic competition, and oligopoly.
\end{itemize}

\subsection*{Skills}
\begin{itemize}
    \item Use supply and demand diagrams to predict the effects of changes in market conditions and policy.
    \item Compute and interpret elasticities, surpluses, costs, and profit-maximizing output and price.
    \item Read and evaluate economic arguments critically and communicate them clearly.
\end{itemize}

\subsection*{Perspectives}
\begin{itemize}
    \item Understand marginal decision making as the core of economic reasoning.
    \item Understand when markets work well, when they fail, and what policy can and cannot fix.
\end{itemize}

%-----------------------------------------------------------
\section*{Class Information}
%-----------------------------------------------------------

\begin{tabularx}{\textwidth}{@{}lX@{}}
\textbf{Instructor:} & Onur Altındağ \\
\textbf{Office:} & RAU 019F \\
\textbf{Email:} & \href{mailto:oaltindag@bentley.edu}{oaltindag@bentley.edu} \\
\textbf{Web:} & \url{https://www.onuraltindag.info/} \\
\textbf{Course page:} & \url{https://www.onuraltindag.info/teaching/ec111/} \\
\textbf{Password:} & \texttt{GoFalcons!} \\[0.3cm]
\textbf{Course:} & \sectionname \\
\textbf{Semester:} & Fall 2026 (August 31, 2026 – December 15, 2026) \\
\textbf{Meeting Time:} & \meetingtime \\
\textbf{Location:} & \classroom \\
\textbf{Delivery Mode:} & In-Person Lecture \\
\end{tabularx}

\vspace{0.3cm}

Office hours are on Mondays and Wednesdays, 11:00 AM -- 12:00 PM and 1:00 -- 3:00 PM, by appointment. Please \href{https://outlook.office.com/bookwithme/user/1c4e4f54e0f0493397e615b2c80a59c7@bentley.edu/meetingtype/rwPThvM0MkK5oGlK_t37nw2?bookingcode=17c9a7b2-69fa-4c48-85ea-ad404642334b&anonymous&ismsaljsauthenabled&ep=mlink}{book a 30-minute slot} to schedule a meeting. If you need to meet urgently or no slots are available, email me directly.

Announcements come by email from me to your Bentley address; check it regularly. Slides, lecture notes, handouts, and mock exams are posted on the course page above after each class. When you email me, put \textbf{EC111} in the subject line.

%-----------------------------------------------------------
\section*{Important Dates and Evaluation}
%-----------------------------------------------------------

\begin{tabularx}{\textwidth}{@{}Xr@{}}
\toprule
\textbf{Component} & \textbf{Weight} \\
\midrule
Weekly assignments + participation & 20\% \\
\textbf{Wednesday, October 7, 2026} – First Midterm & 25\% \\
\textbf{Wednesday, November 11, 2026} – Second Midterm & 25\% \\
\textbf{\finaltime} – Final Exam & 30\% \\
\bottomrule
\end{tabularx}

%-----------------------------------------------------------
\section*{Grading}
%-----------------------------------------------------------

You \textbf{MUST} attend all the midterms and the final as there will be no make-up exams in this course. The midterms are \textbf{NOT} cumulative. If you miss or are likely to miss a midterm due to an emergency, please contact me as soon as possible. You will need to provide supporting documentation/verification of your absence. I will re-weight your final exam accordingly if you have a valid excuse. Please note that family vacations or missing the school shuttle are not valid excuses. If a serious illness, a death in the family, or another difficult circumstance keeps you from an exam or assignment, contact the Office of Student Success (JEN 336, 781.891.2803); they will notify all of your instructors on your behalf.

The final exam \textbf{is} cumulative. If you miss the final exam due to an emergency, you will receive an incomplete for this course. Do not take this class if you know that you will not be able to attend the final exam.

Exams are closed book and closed notes. A basic (non-graphing) calculator is allowed; phones, graphing calculators, and any other device are not, and calculators may not be shared.

You will have \textbf{10 homework assignments} on MindTap, roughly one per week (none in exam weeks), due \textbf{Sunday 11:00 PM}. Your \textbf{two lowest scores are dropped}; the remaining 8 count equally. Each question allows three attempts; an assignment scoring at least \textbf{50\%} receives full credit, otherwise zero. The submission system is automated and will not accept anything after the deadline (even 5 minutes), and there are no exceptions. The questions are algorithmic and slightly different for everyone, so you are encouraged to work together.

\subsection*{Letter grades}
\begin{center}
\begin{tabular}{@{}llllll@{}}
\toprule
A & 94--100 & B+ & 87--89 & C+ & 77--79 \\
A$-$ & 90--93 & B & 83--86 & C & 73--76 \\
 & & B$-$ & 80--82 & C$-$ & 70--72 \\
\midrule
D+ & 67--69 & D & 63--66 & D$-$ & 60--62 \\
F & below 60 & & & & \\
\bottomrule
\end{tabular}
\end{center}
I may lower these cutoffs at the end of the semester; I will not raise them. I do not respond to requests for individual extra credit or for grade adjustments after the fact.

%-----------------------------------------------------------
\section*{Textbook and MindTap}
%-----------------------------------------------------------

The course uses \textbf{MindTap for Mankiw, \textit{Principles of Microeconomics}} (Cengage). The e-book, your homework assignments, and practice questions all live there; a hard copy of the textbook is optional. You \textbf{MUST} register, using your Bentley email, through the link for your section:

\begin{center}
\textbf{\mindtapsection}\\
\url{\mindtaplink}
\end{center}

You have free temporary access until \mindtapaccess; after that you will be prompted to purchase (if your course materials are included in your tuition, no purchase is needed). Registration help: \url{https://startstrong.cengage.com}; technical support: \url{https://support.cengage.com}. Register in the first week; MindTap problems are not an excuse for a missed assignment.

\ifdefined\honors
%-----------------------------------------------------------
\section*{Honors Reading List}
%-----------------------------------------------------------

In addition to the textbook, honors students read and discuss a short set of articles and book chapters over the semester. The list and discussion dates will be posted on the course page in the first weeks of the semester.

\fi
%-----------------------------------------------------------
\section*{Academic Integrity}
%-----------------------------------------------------------

All students are expected to adhere to Bentley's Academic Integrity policy, which includes Bentley's Honor Code (details are in the Undergraduate Student Handbook, the \href{https://catalog.bentley.edu/undergraduate/academic-policies-procedures/}{academic policies section of the catalog}, and the academic integrity course page on Brightspace in which all students and faculty are enrolled). The essence of the policy is that you should not represent someone else's work as your own: no plagiarism, no cheating on exams, no illicit collaboration. Failure to adhere to the policy has serious consequences, including course failure, suspension, or expulsion from the university. The best way to avoid a problem is to consult with me before taking an action that might constitute a violation.

\subsection*{Generative AI}
Nothing in this course is graded on writing, and the homework is algorithmic practice, so the line is simple. You may use generative AI tools (ChatGPT, Claude, Gemini, Copilot, and the like) to study: to explain a concept, to check your reasoning on a practice problem, or to work through a homework question you are stuck on. Treat them as a tutor, not a substitute; the exams are closed book and closed device, and the only way to pass them is to be able to do the work yourself. Any use of AI or any other device during an exam is a violation of the academic integrity policy.

%-----------------------------------------------------------
\section*{Inclusion}
%-----------------------------------------------------------

Everyone in this class has different life experiences and perspectives, and all are valid. As the instructor, I will behave maturely and respectfully in all our class-related engagements, and I expect the same from you. We all adhere to Bentley's standards of conduct, the \href{https://www.bentley.edu/about/mission-and-values}{Bentley Core Values}, in class and in all course work. If you feel that I or anyone in this class has acted outside those values, come and talk to me. If you would rather not, speak with someone in the Office of Student Success (JEN 336, 781.891.2803). My roster shows your name as it appears in Workday; tell me early in the semester if you go by a different name or pronoun.

%-----------------------------------------------------------
\section*{Support and Resources}
%-----------------------------------------------------------

\subsection*{Student Accessibility Services}
Bentley University abides by Section 504 of the Rehabilitation Act of 1973 and the Americans with Disabilities Act of 1990, which stipulate that no student shall be denied the benefits of an education solely by reason of a disability. If you have a hidden or visible disability that may require classroom accommodations, contact \href{https://www.bentley.edu/offices/student-accessibility-services}{Student Accessibility Services} (Office of Student Success, JEN 336, 781.891.2004) within the first four weeks of the semester. Send me your accommodation plan as soon as you have it; accommodations are arranged at least one week before each exam.

\subsection*{The Howard A. Winer '58 Lab for Economics, Accounting and Finance (LEAF)}
The LEAF provides a welcoming and inclusive learning environment where students are encouraged to seek academic support for their Economics, Accounting and Finance courses. Students utilizing the LEAF will find knowledgeable peer tutors ready to help their colleagues to thrive in the Bentley business curriculum. LEAF's free, drop-in tutoring takes place in-person at Morison 220D (rear entrance), opening September 14. For additional information visit \url{https://www.bentley.edu/centers/leaf}. LEAF is a supplement to, not a substitute for, office hours.

\subsection*{Bentley policies and procedures}
All courses at Bentley reflect the university's commitment to a set of core values and practices. You are expected to be familiar with:
\begin{itemize}
    \item \href{https://catalog.bentley.edu/undergraduate/academic-policies-procedures/}{Bentley's honor code and academic integrity system}
    \item \href{https://www.bentley.edu/equity-reporting}{Equity and bias reporting forms and procedures} (Bias Incident Response Team, Office of Institutional Equity)
    \item \href{https://www.bentley.edu/offices/student-accessibility-services}{Student disability accommodations}
    \item \href{https://catalog.bentley.edu/undergraduate/academic-policies-procedures/}{Religious observances policy}
    \item \href{https://www.bentley.edu/about/mission-and-values}{Bentley's core values}
\end{itemize}

\subsection*{Course materials}
Slides, notes, handouts, and exams for this course are for your own use during this semester. They are protected by copyright; do not post or distribute them outside the class without my written consent. The course page is password-protected; do not share the password outside the class.

%-----------------------------------------------------------
\section*{Attendance Policy}
%-----------------------------------------------------------

All students must attend the in-person classes. In-person attendance is essential because it allows for real-time interaction, immediate feedback, and the kind of dialogue that makes learning stick. I do not record lectures, so if you miss a class, you miss the material.

Remote attendance is not an option. This course follows Bentley's \href{https://catalog.bentley.edu/undergraduate/academic-policies-procedures/}{Academic Engagement/Attendance Policy}: students who do not attend or engage during the first week are dropped as no-shows, and two or more consecutive weeks of non-participation can lead to administrative withdrawal. If you must miss class, you are responsible for the material; read the chapter and the posted notes, then bring specific questions to office hours.

\textbf{Weather and closures.} If the university cancels classes, I will email you whether there is an alternative assignment or arrangement for that day.

\textbf{Electronic devices.} Phones away and laptops and tablets closed for the whole class. The work in this course is drawing and reading diagrams, which you learn by doing them by hand, and a screen in your line of sight costs the people around you as much as it costs you. Using a device during class counts against your participation grade. Exceptions are by my permission only; if you have one, you never need to explain it to anyone in the room. Photographing the board at the end of class is fine.

\textbf{Religious observances.} Bentley provides reasonable academic adjustments when a sincerely held religious observance conflicts with a class, exam, or deadline (see the \href{https://downloads.bentley.edu/general/Religious\%20Observances\%202026-2027.pdf}{religious observances calendar} and the \href{https://catalog.bentley.edu/undergraduate/academic-policies-procedures/}{policy}). If you anticipate a conflict, tell me in advance, ideally in the first two weeks of the semester, so we can agree on an adjustment.

%-----------------------------------------------------------
\section*{Tentative Schedule}
%-----------------------------------------------------------

\begin{itemize}
    \item \textbf{Week 1 (Aug 31, Sep 2):} Ten Principles of Economics – Ch. 1 \& Thinking Like an Economist – Ch. 2
    \item \textbf{Week 2 (Sep 9):} Interdependence and the Gains from Trade – Ch. 3 \textit{(no class Sep 7, Labor Day)}
    \item \textbf{Week 3 (Sep 14, 16):} The Market Forces of Supply and Demand – Ch. 4
    \item \textbf{Week 4 (Sep 21, 23):} Elasticity and Its Application – Ch. 5
    \item \textbf{Week 5 (Sep 28, 30):} Supply, Demand, and Government Policies – Ch. 6
    \item \textbf{Week 6 (Oct 5, 7):} Review / \textbf{First Midterm (Oct 7)}
    \item \textbf{Week 7 (Oct 14):} Consumers, Producers, and the Efficiency of Markets – Ch. 7 \textit{(no class Oct 12, Fall Break)}
    \item \textbf{Week 8 (Oct 19, 21):} The Costs of Taxation – Ch. 8 \& International Trade – Ch. 9
    \item \textbf{Week 9 (Oct 26, 28):} Externalities – Ch. 10 \& Public Goods and Common Resources – Ch. 11
    \item \textbf{Week 10 (Nov 2, 4):} The Costs of Production – Ch. 14
    \item \textbf{Week 11 (Nov 9, 11):} Review / \textbf{Second Midterm (Nov 11)}
    \item \textbf{Week 12 (Nov 16, 18):} Firms in Competitive Markets – Ch. 15
    \item \textbf{Week 13 (Nov 23):} Monopoly – Ch. 16 \textit{(no class Nov 25, Thanksgiving Break)}
    \item \textbf{Week 14 (Nov 30, Dec 2):} Monopolistic Competition – Ch. 17 \& Oligopoly – Ch. 18
    \item \textbf{Week 15 (Dec 7):} Review
    \item \textbf{Final Exam:} \finaltime
\end{itemize}

\textit{Note: Chapter references are from Mankiw's \textit{Principles of Microeconomics}, 10th edition (the MindTap e-book). The schedule, and if necessary other parts of this syllabus, may change during the semester; changes are announced in class and by email.}

\end{document}
"
%   xelatex -jobname=EC111_FA2026_Syllabus_20  "\documentclass[11pt,letterpaper]{article}

% Compile with XeLaTeX. Two sections from one source:
%   xelatex -jobname=EC111_FA2026_Syllabus_10H "\def\honors{1}\documentclass[11pt,letterpaper]{article}

% Compile with XeLaTeX. Two sections from one source:
%   xelatex -jobname=EC111_FA2026_Syllabus_10H "\def\honors{1}\input{EC111_FA2026_Syllabus.tex}"
%   xelatex -jobname=EC111_FA2026_Syllabus_20  "\input{EC111_FA2026_Syllabus.tex}"
\ifdefined\honors
  \newcommand{\sectioncode}{EC 111-10-H}
  \newcommand{\sectionname}{EC 111-10-H Principles of Microeconomics (Honors)}
  \newcommand{\meetingtime}{Monday/Wednesday, 3:30 PM -- 4:50 PM}
  \newcommand{\classroom}{Smith Technology Center 103}
  \newcommand{\finaltime}{Friday, December 11, 2026, 6:00 PM -- 8:00 PM}
  \newcommand{\mindtapsection}{Section 10-H of EC 111 Principles of Microeconomics Fall 2026}
  \newcommand{\mindtaplink}{https://student.cengage.com/course-link/MTPNS57GW6XN}
  \newcommand{\mindtapaccess}{the end of Sunday, September 13, 2026}
\else
  \newcommand{\sectioncode}{EC 111-20}
  \newcommand{\sectionname}{EC 111-20 Principles of Microeconomics}
  \newcommand{\meetingtime}{Monday/Wednesday, 9:30 AM -- 10:50 AM}
  \newcommand{\classroom}{Jennison Hall 412}
  \newcommand{\finaltime}{Friday, December 11, 2026, 8:00 AM -- 10:00 AM}
  \newcommand{\mindtapsection}{Section 20 of EC 111 Principles of Microeconomics Fall 2026}
  \newcommand{\mindtaplink}{https://student.cengage.com/course-link/MTPP8N7GWX7L}
  \newcommand{\mindtapaccess}{the end of Sunday, September 13, 2026}
\fi

% XeLaTeX for Turkish character support
\usepackage{fontspec}
\setmainfont{DejaVuSerif}[Extension=.ttf, BoldFont=DejaVuSerif-Bold, ItalicFont=DejaVuSerif-Italic, BoldItalicFont=DejaVuSerif-BoldItalic]
\setsansfont{DejaVuSans}[Extension=.ttf, BoldFont=DejaVuSans-Bold, ItalicFont=DejaVuSans-Oblique, BoldItalicFont=DejaVuSans-BoldOblique]

% Page layout
\usepackage[margin=1in]{geometry}
\usepackage{setspace}
\onehalfspacing

% Colors and styling
\usepackage[dvipsnames]{xcolor}
\definecolor{bentleyblue}{RGB}{0, 51, 102}
\definecolor{linkblue}{RGB}{0, 102, 204}

% Headers and sections
\usepackage{titlesec}
\titleformat{\section}{\Large\bfseries\color{bentleyblue}}{\thesection}{1em}{}[\titlerule]
\titleformat{\subsection}{\large\bfseries\color{bentleyblue!80}}{\thesubsection}{1em}{}
\titleformat{\subsubsection}{\normalsize\bfseries\color{bentleyblue!60}}{\thesubsubsection}{1em}{}

% Lists
\usepackage{enumitem}
\setlist[itemize]{leftmargin=*, itemsep=2pt, parsep=0pt}
\setlist[enumerate]{leftmargin=*, itemsep=2pt, parsep=0pt}

% Tables
\usepackage{booktabs}
\usepackage{tabularx}

% Links
\usepackage{hyperref}
\hypersetup{
    colorlinks=true,
    linkcolor=linkblue,
    urlcolor=linkblue,
    pdfauthor={Onur Altındağ},
    pdftitle={EC111 Principles of Microeconomics - Fall 2026 Syllabus}
}

% Header/Footer
\usepackage{fancyhdr}
\pagestyle{fancy}
\fancyhf{}
\fancyhead[L]{\small\color{gray}\sectioncode{} Principles of Microeconomics}
\fancyhead[R]{\small\color{gray}Fall 2026}
\fancyfoot[C]{\thepage}
\renewcommand{\headrulewidth}{0.4pt}

% Remove paragraph indent, add space between paragraphs
\setlength{\parindent}{0pt}
\setlength{\parskip}{6pt}

\begin{document}

% Title
\begin{center}
{\LARGE\bfseries\color{bentleyblue} EC 111: Principles of Microeconomics}\\[0.3cm]
{\Large Fall 2026 Syllabus}\\[0.5cm]
{\large Bentley University}
\end{center}

\vspace{0.5cm}

%-----------------------------------------------------------
\section*{Course Description}
%-----------------------------------------------------------

This course introduces the fundamental principles and tools of microeconomics: how individuals and firms make decisions when resources are scarce, and how markets coordinate those decisions. We study demand, supply, and market equilibrium; elasticity; the effects of government policies such as taxes and price controls; the efficiency of markets and the sources of market failure; the costs of production; and firm behavior in different market structures, from perfect competition to monopoly, monopolistic competition, and oligopoly.

The emphasis throughout is on economic reasoning: thinking at the margin, recognizing trade-offs, and using simple graphical and numerical models to make sense of real-world questions in business and policy.

\subsection*{Knowledge}
\begin{itemize}
    \item Scarcity, opportunity cost, trade-offs, and the gains from trade.
    \item Demand, supply, market equilibrium, and how markets respond to shocks.
    \item Elasticity and its role in pricing, revenue, and tax incidence.
    \item Consumer surplus, producer surplus, market efficiency, and the deadweight loss of taxes and price controls.
    \item Externalities, public goods, and the case for government intervention.
    \item Costs of production and firm behavior under perfect competition, monopoly, monopolistic competition, and oligopoly.
\end{itemize}

\subsection*{Skills}
\begin{itemize}
    \item Use supply and demand diagrams to predict the effects of changes in market conditions and policy.
    \item Compute and interpret elasticities, surpluses, costs, and profit-maximizing output and price.
    \item Read and evaluate economic arguments critically and communicate them clearly.
\end{itemize}

\subsection*{Perspectives}
\begin{itemize}
    \item Understand marginal decision making as the core of economic reasoning.
    \item Understand when markets work well, when they fail, and what policy can and cannot fix.
\end{itemize}

%-----------------------------------------------------------
\section*{Class Information}
%-----------------------------------------------------------

\begin{tabularx}{\textwidth}{@{}lX@{}}
\textbf{Instructor:} & Onur Altındağ \\
\textbf{Office:} & RAU 019F \\
\textbf{Email:} & \href{mailto:oaltindag@bentley.edu}{oaltindag@bentley.edu} \\
\textbf{Web:} & \url{https://www.onuraltindag.info/} \\
\textbf{Course page:} & \url{https://www.onuraltindag.info/teaching/ec111/} \\
\textbf{Password:} & \texttt{GoFalcons!} \\[0.3cm]
\textbf{Course:} & \sectionname \\
\textbf{Semester:} & Fall 2026 (August 31, 2026 – December 15, 2026) \\
\textbf{Meeting Time:} & \meetingtime \\
\textbf{Location:} & \classroom \\
\textbf{Delivery Mode:} & In-Person Lecture \\
\end{tabularx}

\vspace{0.3cm}

Office hours are on Mondays and Wednesdays, 11:00 AM -- 12:00 PM and 1:00 -- 3:00 PM, by appointment. Please \href{https://outlook.office.com/bookwithme/user/1c4e4f54e0f0493397e615b2c80a59c7@bentley.edu/meetingtype/rwPThvM0MkK5oGlK_t37nw2?bookingcode=17c9a7b2-69fa-4c48-85ea-ad404642334b&anonymous&ismsaljsauthenabled&ep=mlink}{book a 30-minute slot} to schedule a meeting. If you need to meet urgently or no slots are available, email me directly.

Announcements come by email from me to your Bentley address; check it regularly. Slides, lecture notes, handouts, and mock exams are posted on the course page above after each class. When you email me, put \textbf{EC111} in the subject line.

%-----------------------------------------------------------
\section*{Important Dates and Evaluation}
%-----------------------------------------------------------

\begin{tabularx}{\textwidth}{@{}Xr@{}}
\toprule
\textbf{Component} & \textbf{Weight} \\
\midrule
Weekly assignments + participation & 20\% \\
\textbf{Wednesday, October 7, 2026} – First Midterm & 25\% \\
\textbf{Wednesday, November 11, 2026} – Second Midterm & 25\% \\
\textbf{\finaltime} – Final Exam & 30\% \\
\bottomrule
\end{tabularx}

%-----------------------------------------------------------
\section*{Grading}
%-----------------------------------------------------------

You \textbf{MUST} attend all the midterms and the final as there will be no make-up exams in this course. The midterms are \textbf{NOT} cumulative. If you miss or are likely to miss a midterm due to an emergency, please contact me as soon as possible. You will need to provide supporting documentation/verification of your absence. I will re-weight your final exam accordingly if you have a valid excuse. Please note that family vacations or missing the school shuttle are not valid excuses. If a serious illness, a death in the family, or another difficult circumstance keeps you from an exam or assignment, contact the Office of Student Success (JEN 336, 781.891.2803); they will notify all of your instructors on your behalf.

The final exam \textbf{is} cumulative. If you miss the final exam due to an emergency, you will receive an incomplete for this course. Do not take this class if you know that you will not be able to attend the final exam.

Exams are closed book and closed notes. A basic (non-graphing) calculator is allowed; phones, graphing calculators, and any other device are not, and calculators may not be shared.

You will have \textbf{10 homework assignments} on MindTap, roughly one per week (none in exam weeks), due \textbf{Sunday 11:00 PM}. Your \textbf{two lowest scores are dropped}; the remaining 8 count equally. Each question allows three attempts; an assignment scoring at least \textbf{50\%} receives full credit, otherwise zero. The submission system is automated and will not accept anything after the deadline (even 5 minutes), and there are no exceptions. The questions are algorithmic and slightly different for everyone, so you are encouraged to work together.

\subsection*{Letter grades}
\begin{center}
\begin{tabular}{@{}llllll@{}}
\toprule
A & 94--100 & B+ & 87--89 & C+ & 77--79 \\
A$-$ & 90--93 & B & 83--86 & C & 73--76 \\
 & & B$-$ & 80--82 & C$-$ & 70--72 \\
\midrule
D+ & 67--69 & D & 63--66 & D$-$ & 60--62 \\
F & below 60 & & & & \\
\bottomrule
\end{tabular}
\end{center}
I may lower these cutoffs at the end of the semester; I will not raise them. I do not respond to requests for individual extra credit or for grade adjustments after the fact.

%-----------------------------------------------------------
\section*{Textbook and MindTap}
%-----------------------------------------------------------

The course uses \textbf{MindTap for Mankiw, \textit{Principles of Microeconomics}} (Cengage). The e-book, your homework assignments, and practice questions all live there; a hard copy of the textbook is optional. You \textbf{MUST} register, using your Bentley email, through the link for your section:

\begin{center}
\textbf{\mindtapsection}\\
\url{\mindtaplink}
\end{center}

You have free temporary access until \mindtapaccess; after that you will be prompted to purchase (if your course materials are included in your tuition, no purchase is needed). Registration help: \url{https://startstrong.cengage.com}; technical support: \url{https://support.cengage.com}. Register in the first week; MindTap problems are not an excuse for a missed assignment.

\ifdefined\honors
%-----------------------------------------------------------
\section*{Honors Reading List}
%-----------------------------------------------------------

In addition to the textbook, honors students read and discuss a short set of articles and book chapters over the semester. The list and discussion dates will be posted on the course page in the first weeks of the semester.

\fi
%-----------------------------------------------------------
\section*{Academic Integrity}
%-----------------------------------------------------------

All students are expected to adhere to Bentley's Academic Integrity policy, which includes Bentley's Honor Code (details are in the Undergraduate Student Handbook, the \href{https://catalog.bentley.edu/undergraduate/academic-policies-procedures/}{academic policies section of the catalog}, and the academic integrity course page on Brightspace in which all students and faculty are enrolled). The essence of the policy is that you should not represent someone else's work as your own: no plagiarism, no cheating on exams, no illicit collaboration. Failure to adhere to the policy has serious consequences, including course failure, suspension, or expulsion from the university. The best way to avoid a problem is to consult with me before taking an action that might constitute a violation.

\subsection*{Generative AI}
Nothing in this course is graded on writing, and the homework is algorithmic practice, so the line is simple. You may use generative AI tools (ChatGPT, Claude, Gemini, Copilot, and the like) to study: to explain a concept, to check your reasoning on a practice problem, or to work through a homework question you are stuck on. Treat them as a tutor, not a substitute; the exams are closed book and closed device, and the only way to pass them is to be able to do the work yourself. Any use of AI or any other device during an exam is a violation of the academic integrity policy.

%-----------------------------------------------------------
\section*{Inclusion}
%-----------------------------------------------------------

Everyone in this class has different life experiences and perspectives, and all are valid. As the instructor, I will behave maturely and respectfully in all our class-related engagements, and I expect the same from you. We all adhere to Bentley's standards of conduct, the \href{https://www.bentley.edu/about/mission-and-values}{Bentley Core Values}, in class and in all course work. If you feel that I or anyone in this class has acted outside those values, come and talk to me. If you would rather not, speak with someone in the Office of Student Success (JEN 336, 781.891.2803). My roster shows your name as it appears in Workday; tell me early in the semester if you go by a different name or pronoun.

%-----------------------------------------------------------
\section*{Support and Resources}
%-----------------------------------------------------------

\subsection*{Student Accessibility Services}
Bentley University abides by Section 504 of the Rehabilitation Act of 1973 and the Americans with Disabilities Act of 1990, which stipulate that no student shall be denied the benefits of an education solely by reason of a disability. If you have a hidden or visible disability that may require classroom accommodations, contact \href{https://www.bentley.edu/offices/student-accessibility-services}{Student Accessibility Services} (Office of Student Success, JEN 336, 781.891.2004) within the first four weeks of the semester. Send me your accommodation plan as soon as you have it; accommodations are arranged at least one week before each exam.

\subsection*{The Howard A. Winer '58 Lab for Economics, Accounting and Finance (LEAF)}
The LEAF provides a welcoming and inclusive learning environment where students are encouraged to seek academic support for their Economics, Accounting and Finance courses. Students utilizing the LEAF will find knowledgeable peer tutors ready to help their colleagues to thrive in the Bentley business curriculum. LEAF's free, drop-in tutoring takes place in-person at Morison 220D (rear entrance), opening September 14. For additional information visit \url{https://www.bentley.edu/centers/leaf}. LEAF is a supplement to, not a substitute for, office hours.

\subsection*{Bentley policies and procedures}
All courses at Bentley reflect the university's commitment to a set of core values and practices. You are expected to be familiar with:
\begin{itemize}
    \item \href{https://catalog.bentley.edu/undergraduate/academic-policies-procedures/}{Bentley's honor code and academic integrity system}
    \item \href{https://www.bentley.edu/equity-reporting}{Equity and bias reporting forms and procedures} (Bias Incident Response Team, Office of Institutional Equity)
    \item \href{https://www.bentley.edu/offices/student-accessibility-services}{Student disability accommodations}
    \item \href{https://catalog.bentley.edu/undergraduate/academic-policies-procedures/}{Religious observances policy}
    \item \href{https://www.bentley.edu/about/mission-and-values}{Bentley's core values}
\end{itemize}

\subsection*{Course materials}
Slides, notes, handouts, and exams for this course are for your own use during this semester. They are protected by copyright; do not post or distribute them outside the class without my written consent. The course page is password-protected; do not share the password outside the class.

%-----------------------------------------------------------
\section*{Attendance Policy}
%-----------------------------------------------------------

All students must attend the in-person classes. In-person attendance is essential because it allows for real-time interaction, immediate feedback, and the kind of dialogue that makes learning stick. I do not record lectures, so if you miss a class, you miss the material.

Remote attendance is not an option. This course follows Bentley's \href{https://catalog.bentley.edu/undergraduate/academic-policies-procedures/}{Academic Engagement/Attendance Policy}: students who do not attend or engage during the first week are dropped as no-shows, and two or more consecutive weeks of non-participation can lead to administrative withdrawal. If you must miss class, you are responsible for the material; read the chapter and the posted notes, then bring specific questions to office hours.

\textbf{Weather and closures.} If the university cancels classes, I will email you whether there is an alternative assignment or arrangement for that day.

\textbf{Electronic devices.} Phones away and laptops and tablets closed for the whole class. The work in this course is drawing and reading diagrams, which you learn by doing them by hand, and a screen in your line of sight costs the people around you as much as it costs you. Using a device during class counts against your participation grade. Exceptions are by my permission only; if you have one, you never need to explain it to anyone in the room. Photographing the board at the end of class is fine.

\textbf{Religious observances.} Bentley provides reasonable academic adjustments when a sincerely held religious observance conflicts with a class, exam, or deadline (see the \href{https://downloads.bentley.edu/general/Religious\%20Observances\%202026-2027.pdf}{religious observances calendar} and the \href{https://catalog.bentley.edu/undergraduate/academic-policies-procedures/}{policy}). If you anticipate a conflict, tell me in advance, ideally in the first two weeks of the semester, so we can agree on an adjustment.

%-----------------------------------------------------------
\section*{Tentative Schedule}
%-----------------------------------------------------------

\begin{itemize}
    \item \textbf{Week 1 (Aug 31, Sep 2):} Ten Principles of Economics – Ch. 1 \& Thinking Like an Economist – Ch. 2
    \item \textbf{Week 2 (Sep 9):} Interdependence and the Gains from Trade – Ch. 3 \textit{(no class Sep 7, Labor Day)}
    \item \textbf{Week 3 (Sep 14, 16):} The Market Forces of Supply and Demand – Ch. 4
    \item \textbf{Week 4 (Sep 21, 23):} Elasticity and Its Application – Ch. 5
    \item \textbf{Week 5 (Sep 28, 30):} Supply, Demand, and Government Policies – Ch. 6
    \item \textbf{Week 6 (Oct 5, 7):} Review / \textbf{First Midterm (Oct 7)}
    \item \textbf{Week 7 (Oct 14):} Consumers, Producers, and the Efficiency of Markets – Ch. 7 \textit{(no class Oct 12, Fall Break)}
    \item \textbf{Week 8 (Oct 19, 21):} The Costs of Taxation – Ch. 8 \& International Trade – Ch. 9
    \item \textbf{Week 9 (Oct 26, 28):} Externalities – Ch. 10 \& Public Goods and Common Resources – Ch. 11
    \item \textbf{Week 10 (Nov 2, 4):} The Costs of Production – Ch. 14
    \item \textbf{Week 11 (Nov 9, 11):} Review / \textbf{Second Midterm (Nov 11)}
    \item \textbf{Week 12 (Nov 16, 18):} Firms in Competitive Markets – Ch. 15
    \item \textbf{Week 13 (Nov 23):} Monopoly – Ch. 16 \textit{(no class Nov 25, Thanksgiving Break)}
    \item \textbf{Week 14 (Nov 30, Dec 2):} Monopolistic Competition – Ch. 17 \& Oligopoly – Ch. 18
    \item \textbf{Week 15 (Dec 7):} Review
    \item \textbf{Final Exam:} \finaltime
\end{itemize}

\textit{Note: Chapter references are from Mankiw's \textit{Principles of Microeconomics}, 10th edition (the MindTap e-book). The schedule, and if necessary other parts of this syllabus, may change during the semester; changes are announced in class and by email.}

\end{document}
"
%   xelatex -jobname=EC111_FA2026_Syllabus_20  "\documentclass[11pt,letterpaper]{article}

% Compile with XeLaTeX. Two sections from one source:
%   xelatex -jobname=EC111_FA2026_Syllabus_10H "\def\honors{1}\input{EC111_FA2026_Syllabus.tex}"
%   xelatex -jobname=EC111_FA2026_Syllabus_20  "\input{EC111_FA2026_Syllabus.tex}"
\ifdefined\honors
  \newcommand{\sectioncode}{EC 111-10-H}
  \newcommand{\sectionname}{EC 111-10-H Principles of Microeconomics (Honors)}
  \newcommand{\meetingtime}{Monday/Wednesday, 3:30 PM -- 4:50 PM}
  \newcommand{\classroom}{Smith Technology Center 103}
  \newcommand{\finaltime}{Friday, December 11, 2026, 6:00 PM -- 8:00 PM}
  \newcommand{\mindtapsection}{Section 10-H of EC 111 Principles of Microeconomics Fall 2026}
  \newcommand{\mindtaplink}{https://student.cengage.com/course-link/MTPNS57GW6XN}
  \newcommand{\mindtapaccess}{the end of Sunday, September 13, 2026}
\else
  \newcommand{\sectioncode}{EC 111-20}
  \newcommand{\sectionname}{EC 111-20 Principles of Microeconomics}
  \newcommand{\meetingtime}{Monday/Wednesday, 9:30 AM -- 10:50 AM}
  \newcommand{\classroom}{Jennison Hall 412}
  \newcommand{\finaltime}{Friday, December 11, 2026, 8:00 AM -- 10:00 AM}
  \newcommand{\mindtapsection}{Section 20 of EC 111 Principles of Microeconomics Fall 2026}
  \newcommand{\mindtaplink}{https://student.cengage.com/course-link/MTPP8N7GWX7L}
  \newcommand{\mindtapaccess}{the end of Sunday, September 13, 2026}
\fi

% XeLaTeX for Turkish character support
\usepackage{fontspec}
\setmainfont{DejaVuSerif}[Extension=.ttf, BoldFont=DejaVuSerif-Bold, ItalicFont=DejaVuSerif-Italic, BoldItalicFont=DejaVuSerif-BoldItalic]
\setsansfont{DejaVuSans}[Extension=.ttf, BoldFont=DejaVuSans-Bold, ItalicFont=DejaVuSans-Oblique, BoldItalicFont=DejaVuSans-BoldOblique]

% Page layout
\usepackage[margin=1in]{geometry}
\usepackage{setspace}
\onehalfspacing

% Colors and styling
\usepackage[dvipsnames]{xcolor}
\definecolor{bentleyblue}{RGB}{0, 51, 102}
\definecolor{linkblue}{RGB}{0, 102, 204}

% Headers and sections
\usepackage{titlesec}
\titleformat{\section}{\Large\bfseries\color{bentleyblue}}{\thesection}{1em}{}[\titlerule]
\titleformat{\subsection}{\large\bfseries\color{bentleyblue!80}}{\thesubsection}{1em}{}
\titleformat{\subsubsection}{\normalsize\bfseries\color{bentleyblue!60}}{\thesubsubsection}{1em}{}

% Lists
\usepackage{enumitem}
\setlist[itemize]{leftmargin=*, itemsep=2pt, parsep=0pt}
\setlist[enumerate]{leftmargin=*, itemsep=2pt, parsep=0pt}

% Tables
\usepackage{booktabs}
\usepackage{tabularx}

% Links
\usepackage{hyperref}
\hypersetup{
    colorlinks=true,
    linkcolor=linkblue,
    urlcolor=linkblue,
    pdfauthor={Onur Altındağ},
    pdftitle={EC111 Principles of Microeconomics - Fall 2026 Syllabus}
}

% Header/Footer
\usepackage{fancyhdr}
\pagestyle{fancy}
\fancyhf{}
\fancyhead[L]{\small\color{gray}\sectioncode{} Principles of Microeconomics}
\fancyhead[R]{\small\color{gray}Fall 2026}
\fancyfoot[C]{\thepage}
\renewcommand{\headrulewidth}{0.4pt}

% Remove paragraph indent, add space between paragraphs
\setlength{\parindent}{0pt}
\setlength{\parskip}{6pt}

\begin{document}

% Title
\begin{center}
{\LARGE\bfseries\color{bentleyblue} EC 111: Principles of Microeconomics}\\[0.3cm]
{\Large Fall 2026 Syllabus}\\[0.5cm]
{\large Bentley University}
\end{center}

\vspace{0.5cm}

%-----------------------------------------------------------
\section*{Course Description}
%-----------------------------------------------------------

This course introduces the fundamental principles and tools of microeconomics: how individuals and firms make decisions when resources are scarce, and how markets coordinate those decisions. We study demand, supply, and market equilibrium; elasticity; the effects of government policies such as taxes and price controls; the efficiency of markets and the sources of market failure; the costs of production; and firm behavior in different market structures, from perfect competition to monopoly, monopolistic competition, and oligopoly.

The emphasis throughout is on economic reasoning: thinking at the margin, recognizing trade-offs, and using simple graphical and numerical models to make sense of real-world questions in business and policy.

\subsection*{Knowledge}
\begin{itemize}
    \item Scarcity, opportunity cost, trade-offs, and the gains from trade.
    \item Demand, supply, market equilibrium, and how markets respond to shocks.
    \item Elasticity and its role in pricing, revenue, and tax incidence.
    \item Consumer surplus, producer surplus, market efficiency, and the deadweight loss of taxes and price controls.
    \item Externalities, public goods, and the case for government intervention.
    \item Costs of production and firm behavior under perfect competition, monopoly, monopolistic competition, and oligopoly.
\end{itemize}

\subsection*{Skills}
\begin{itemize}
    \item Use supply and demand diagrams to predict the effects of changes in market conditions and policy.
    \item Compute and interpret elasticities, surpluses, costs, and profit-maximizing output and price.
    \item Read and evaluate economic arguments critically and communicate them clearly.
\end{itemize}

\subsection*{Perspectives}
\begin{itemize}
    \item Understand marginal decision making as the core of economic reasoning.
    \item Understand when markets work well, when they fail, and what policy can and cannot fix.
\end{itemize}

%-----------------------------------------------------------
\section*{Class Information}
%-----------------------------------------------------------

\begin{tabularx}{\textwidth}{@{}lX@{}}
\textbf{Instructor:} & Onur Altındağ \\
\textbf{Office:} & RAU 019F \\
\textbf{Email:} & \href{mailto:oaltindag@bentley.edu}{oaltindag@bentley.edu} \\
\textbf{Web:} & \url{https://www.onuraltindag.info/} \\
\textbf{Course page:} & \url{https://www.onuraltindag.info/teaching/ec111/} \\
\textbf{Password:} & \texttt{GoFalcons!} \\[0.3cm]
\textbf{Course:} & \sectionname \\
\textbf{Semester:} & Fall 2026 (August 31, 2026 – December 15, 2026) \\
\textbf{Meeting Time:} & \meetingtime \\
\textbf{Location:} & \classroom \\
\textbf{Delivery Mode:} & In-Person Lecture \\
\end{tabularx}

\vspace{0.3cm}

Office hours are on Mondays and Wednesdays, 11:00 AM -- 12:00 PM and 1:00 -- 3:00 PM, by appointment. Please \href{https://outlook.office.com/bookwithme/user/1c4e4f54e0f0493397e615b2c80a59c7@bentley.edu/meetingtype/rwPThvM0MkK5oGlK_t37nw2?bookingcode=17c9a7b2-69fa-4c48-85ea-ad404642334b&anonymous&ismsaljsauthenabled&ep=mlink}{book a 30-minute slot} to schedule a meeting. If you need to meet urgently or no slots are available, email me directly.

Announcements come by email from me to your Bentley address; check it regularly. Slides, lecture notes, handouts, and mock exams are posted on the course page above after each class. When you email me, put \textbf{EC111} in the subject line.

%-----------------------------------------------------------
\section*{Important Dates and Evaluation}
%-----------------------------------------------------------

\begin{tabularx}{\textwidth}{@{}Xr@{}}
\toprule
\textbf{Component} & \textbf{Weight} \\
\midrule
Weekly assignments + participation & 20\% \\
\textbf{Wednesday, October 7, 2026} – First Midterm & 25\% \\
\textbf{Wednesday, November 11, 2026} – Second Midterm & 25\% \\
\textbf{\finaltime} – Final Exam & 30\% \\
\bottomrule
\end{tabularx}

%-----------------------------------------------------------
\section*{Grading}
%-----------------------------------------------------------

You \textbf{MUST} attend all the midterms and the final as there will be no make-up exams in this course. The midterms are \textbf{NOT} cumulative. If you miss or are likely to miss a midterm due to an emergency, please contact me as soon as possible. You will need to provide supporting documentation/verification of your absence. I will re-weight your final exam accordingly if you have a valid excuse. Please note that family vacations or missing the school shuttle are not valid excuses. If a serious illness, a death in the family, or another difficult circumstance keeps you from an exam or assignment, contact the Office of Student Success (JEN 336, 781.891.2803); they will notify all of your instructors on your behalf.

The final exam \textbf{is} cumulative. If you miss the final exam due to an emergency, you will receive an incomplete for this course. Do not take this class if you know that you will not be able to attend the final exam.

Exams are closed book and closed notes. A basic (non-graphing) calculator is allowed; phones, graphing calculators, and any other device are not, and calculators may not be shared.

You will have \textbf{10 homework assignments} on MindTap, roughly one per week (none in exam weeks), due \textbf{Sunday 11:00 PM}. Your \textbf{two lowest scores are dropped}; the remaining 8 count equally. Each question allows three attempts; an assignment scoring at least \textbf{50\%} receives full credit, otherwise zero. The submission system is automated and will not accept anything after the deadline (even 5 minutes), and there are no exceptions. The questions are algorithmic and slightly different for everyone, so you are encouraged to work together.

\subsection*{Letter grades}
\begin{center}
\begin{tabular}{@{}llllll@{}}
\toprule
A & 94--100 & B+ & 87--89 & C+ & 77--79 \\
A$-$ & 90--93 & B & 83--86 & C & 73--76 \\
 & & B$-$ & 80--82 & C$-$ & 70--72 \\
\midrule
D+ & 67--69 & D & 63--66 & D$-$ & 60--62 \\
F & below 60 & & & & \\
\bottomrule
\end{tabular}
\end{center}
I may lower these cutoffs at the end of the semester; I will not raise them. I do not respond to requests for individual extra credit or for grade adjustments after the fact.

%-----------------------------------------------------------
\section*{Textbook and MindTap}
%-----------------------------------------------------------

The course uses \textbf{MindTap for Mankiw, \textit{Principles of Microeconomics}} (Cengage). The e-book, your homework assignments, and practice questions all live there; a hard copy of the textbook is optional. You \textbf{MUST} register, using your Bentley email, through the link for your section:

\begin{center}
\textbf{\mindtapsection}\\
\url{\mindtaplink}
\end{center}

You have free temporary access until \mindtapaccess; after that you will be prompted to purchase (if your course materials are included in your tuition, no purchase is needed). Registration help: \url{https://startstrong.cengage.com}; technical support: \url{https://support.cengage.com}. Register in the first week; MindTap problems are not an excuse for a missed assignment.

\ifdefined\honors
%-----------------------------------------------------------
\section*{Honors Reading List}
%-----------------------------------------------------------

In addition to the textbook, honors students read and discuss a short set of articles and book chapters over the semester. The list and discussion dates will be posted on the course page in the first weeks of the semester.

\fi
%-----------------------------------------------------------
\section*{Academic Integrity}
%-----------------------------------------------------------

All students are expected to adhere to Bentley's Academic Integrity policy, which includes Bentley's Honor Code (details are in the Undergraduate Student Handbook, the \href{https://catalog.bentley.edu/undergraduate/academic-policies-procedures/}{academic policies section of the catalog}, and the academic integrity course page on Brightspace in which all students and faculty are enrolled). The essence of the policy is that you should not represent someone else's work as your own: no plagiarism, no cheating on exams, no illicit collaboration. Failure to adhere to the policy has serious consequences, including course failure, suspension, or expulsion from the university. The best way to avoid a problem is to consult with me before taking an action that might constitute a violation.

\subsection*{Generative AI}
Nothing in this course is graded on writing, and the homework is algorithmic practice, so the line is simple. You may use generative AI tools (ChatGPT, Claude, Gemini, Copilot, and the like) to study: to explain a concept, to check your reasoning on a practice problem, or to work through a homework question you are stuck on. Treat them as a tutor, not a substitute; the exams are closed book and closed device, and the only way to pass them is to be able to do the work yourself. Any use of AI or any other device during an exam is a violation of the academic integrity policy.

%-----------------------------------------------------------
\section*{Inclusion}
%-----------------------------------------------------------

Everyone in this class has different life experiences and perspectives, and all are valid. As the instructor, I will behave maturely and respectfully in all our class-related engagements, and I expect the same from you. We all adhere to Bentley's standards of conduct, the \href{https://www.bentley.edu/about/mission-and-values}{Bentley Core Values}, in class and in all course work. If you feel that I or anyone in this class has acted outside those values, come and talk to me. If you would rather not, speak with someone in the Office of Student Success (JEN 336, 781.891.2803). My roster shows your name as it appears in Workday; tell me early in the semester if you go by a different name or pronoun.

%-----------------------------------------------------------
\section*{Support and Resources}
%-----------------------------------------------------------

\subsection*{Student Accessibility Services}
Bentley University abides by Section 504 of the Rehabilitation Act of 1973 and the Americans with Disabilities Act of 1990, which stipulate that no student shall be denied the benefits of an education solely by reason of a disability. If you have a hidden or visible disability that may require classroom accommodations, contact \href{https://www.bentley.edu/offices/student-accessibility-services}{Student Accessibility Services} (Office of Student Success, JEN 336, 781.891.2004) within the first four weeks of the semester. Send me your accommodation plan as soon as you have it; accommodations are arranged at least one week before each exam.

\subsection*{The Howard A. Winer '58 Lab for Economics, Accounting and Finance (LEAF)}
The LEAF provides a welcoming and inclusive learning environment where students are encouraged to seek academic support for their Economics, Accounting and Finance courses. Students utilizing the LEAF will find knowledgeable peer tutors ready to help their colleagues to thrive in the Bentley business curriculum. LEAF's free, drop-in tutoring takes place in-person at Morison 220D (rear entrance), opening September 14. For additional information visit \url{https://www.bentley.edu/centers/leaf}. LEAF is a supplement to, not a substitute for, office hours.

\subsection*{Bentley policies and procedures}
All courses at Bentley reflect the university's commitment to a set of core values and practices. You are expected to be familiar with:
\begin{itemize}
    \item \href{https://catalog.bentley.edu/undergraduate/academic-policies-procedures/}{Bentley's honor code and academic integrity system}
    \item \href{https://www.bentley.edu/equity-reporting}{Equity and bias reporting forms and procedures} (Bias Incident Response Team, Office of Institutional Equity)
    \item \href{https://www.bentley.edu/offices/student-accessibility-services}{Student disability accommodations}
    \item \href{https://catalog.bentley.edu/undergraduate/academic-policies-procedures/}{Religious observances policy}
    \item \href{https://www.bentley.edu/about/mission-and-values}{Bentley's core values}
\end{itemize}

\subsection*{Course materials}
Slides, notes, handouts, and exams for this course are for your own use during this semester. They are protected by copyright; do not post or distribute them outside the class without my written consent. The course page is password-protected; do not share the password outside the class.

%-----------------------------------------------------------
\section*{Attendance Policy}
%-----------------------------------------------------------

All students must attend the in-person classes. In-person attendance is essential because it allows for real-time interaction, immediate feedback, and the kind of dialogue that makes learning stick. I do not record lectures, so if you miss a class, you miss the material.

Remote attendance is not an option. This course follows Bentley's \href{https://catalog.bentley.edu/undergraduate/academic-policies-procedures/}{Academic Engagement/Attendance Policy}: students who do not attend or engage during the first week are dropped as no-shows, and two or more consecutive weeks of non-participation can lead to administrative withdrawal. If you must miss class, you are responsible for the material; read the chapter and the posted notes, then bring specific questions to office hours.

\textbf{Weather and closures.} If the university cancels classes, I will email you whether there is an alternative assignment or arrangement for that day.

\textbf{Electronic devices.} Phones away and laptops and tablets closed for the whole class. The work in this course is drawing and reading diagrams, which you learn by doing them by hand, and a screen in your line of sight costs the people around you as much as it costs you. Using a device during class counts against your participation grade. Exceptions are by my permission only; if you have one, you never need to explain it to anyone in the room. Photographing the board at the end of class is fine.

\textbf{Religious observances.} Bentley provides reasonable academic adjustments when a sincerely held religious observance conflicts with a class, exam, or deadline (see the \href{https://downloads.bentley.edu/general/Religious\%20Observances\%202026-2027.pdf}{religious observances calendar} and the \href{https://catalog.bentley.edu/undergraduate/academic-policies-procedures/}{policy}). If you anticipate a conflict, tell me in advance, ideally in the first two weeks of the semester, so we can agree on an adjustment.

%-----------------------------------------------------------
\section*{Tentative Schedule}
%-----------------------------------------------------------

\begin{itemize}
    \item \textbf{Week 1 (Aug 31, Sep 2):} Ten Principles of Economics – Ch. 1 \& Thinking Like an Economist – Ch. 2
    \item \textbf{Week 2 (Sep 9):} Interdependence and the Gains from Trade – Ch. 3 \textit{(no class Sep 7, Labor Day)}
    \item \textbf{Week 3 (Sep 14, 16):} The Market Forces of Supply and Demand – Ch. 4
    \item \textbf{Week 4 (Sep 21, 23):} Elasticity and Its Application – Ch. 5
    \item \textbf{Week 5 (Sep 28, 30):} Supply, Demand, and Government Policies – Ch. 6
    \item \textbf{Week 6 (Oct 5, 7):} Review / \textbf{First Midterm (Oct 7)}
    \item \textbf{Week 7 (Oct 14):} Consumers, Producers, and the Efficiency of Markets – Ch. 7 \textit{(no class Oct 12, Fall Break)}
    \item \textbf{Week 8 (Oct 19, 21):} The Costs of Taxation – Ch. 8 \& International Trade – Ch. 9
    \item \textbf{Week 9 (Oct 26, 28):} Externalities – Ch. 10 \& Public Goods and Common Resources – Ch. 11
    \item \textbf{Week 10 (Nov 2, 4):} The Costs of Production – Ch. 14
    \item \textbf{Week 11 (Nov 9, 11):} Review / \textbf{Second Midterm (Nov 11)}
    \item \textbf{Week 12 (Nov 16, 18):} Firms in Competitive Markets – Ch. 15
    \item \textbf{Week 13 (Nov 23):} Monopoly – Ch. 16 \textit{(no class Nov 25, Thanksgiving Break)}
    \item \textbf{Week 14 (Nov 30, Dec 2):} Monopolistic Competition – Ch. 17 \& Oligopoly – Ch. 18
    \item \textbf{Week 15 (Dec 7):} Review
    \item \textbf{Final Exam:} \finaltime
\end{itemize}

\textit{Note: Chapter references are from Mankiw's \textit{Principles of Microeconomics}, 10th edition (the MindTap e-book). The schedule, and if necessary other parts of this syllabus, may change during the semester; changes are announced in class and by email.}

\end{document}
"
\ifdefined\honors
  \newcommand{\sectioncode}{EC 111-10-H}
  \newcommand{\sectionname}{EC 111-10-H Principles of Microeconomics (Honors)}
  \newcommand{\meetingtime}{Monday/Wednesday, 3:30 PM -- 4:50 PM}
  \newcommand{\classroom}{Smith Technology Center 103}
  \newcommand{\finaltime}{Friday, December 11, 2026, 6:00 PM -- 8:00 PM}
  \newcommand{\mindtapsection}{Section 10-H of EC 111 Principles of Microeconomics Fall 2026}
  \newcommand{\mindtaplink}{https://student.cengage.com/course-link/MTPNS57GW6XN}
  \newcommand{\mindtapaccess}{the end of Sunday, September 13, 2026}
\else
  \newcommand{\sectioncode}{EC 111-20}
  \newcommand{\sectionname}{EC 111-20 Principles of Microeconomics}
  \newcommand{\meetingtime}{Monday/Wednesday, 9:30 AM -- 10:50 AM}
  \newcommand{\classroom}{Jennison Hall 412}
  \newcommand{\finaltime}{Friday, December 11, 2026, 8:00 AM -- 10:00 AM}
  \newcommand{\mindtapsection}{Section 20 of EC 111 Principles of Microeconomics Fall 2026}
  \newcommand{\mindtaplink}{https://student.cengage.com/course-link/MTPP8N7GWX7L}
  \newcommand{\mindtapaccess}{the end of Sunday, September 13, 2026}
\fi

% XeLaTeX for Turkish character support
\usepackage{fontspec}
\setmainfont{DejaVuSerif}[Extension=.ttf, BoldFont=DejaVuSerif-Bold, ItalicFont=DejaVuSerif-Italic, BoldItalicFont=DejaVuSerif-BoldItalic]
\setsansfont{DejaVuSans}[Extension=.ttf, BoldFont=DejaVuSans-Bold, ItalicFont=DejaVuSans-Oblique, BoldItalicFont=DejaVuSans-BoldOblique]

% Page layout
\usepackage[margin=1in]{geometry}
\usepackage{setspace}
\onehalfspacing

% Colors and styling
\usepackage[dvipsnames]{xcolor}
\definecolor{bentleyblue}{RGB}{0, 51, 102}
\definecolor{linkblue}{RGB}{0, 102, 204}

% Headers and sections
\usepackage{titlesec}
\titleformat{\section}{\Large\bfseries\color{bentleyblue}}{\thesection}{1em}{}[\titlerule]
\titleformat{\subsection}{\large\bfseries\color{bentleyblue!80}}{\thesubsection}{1em}{}
\titleformat{\subsubsection}{\normalsize\bfseries\color{bentleyblue!60}}{\thesubsubsection}{1em}{}

% Lists
\usepackage{enumitem}
\setlist[itemize]{leftmargin=*, itemsep=2pt, parsep=0pt}
\setlist[enumerate]{leftmargin=*, itemsep=2pt, parsep=0pt}

% Tables
\usepackage{booktabs}
\usepackage{tabularx}

% Links
\usepackage{hyperref}
\hypersetup{
    colorlinks=true,
    linkcolor=linkblue,
    urlcolor=linkblue,
    pdfauthor={Onur Altındağ},
    pdftitle={EC111 Principles of Microeconomics - Fall 2026 Syllabus}
}

% Header/Footer
\usepackage{fancyhdr}
\pagestyle{fancy}
\fancyhf{}
\fancyhead[L]{\small\color{gray}\sectioncode{} Principles of Microeconomics}
\fancyhead[R]{\small\color{gray}Fall 2026}
\fancyfoot[C]{\thepage}
\renewcommand{\headrulewidth}{0.4pt}

% Remove paragraph indent, add space between paragraphs
\setlength{\parindent}{0pt}
\setlength{\parskip}{6pt}

\begin{document}

% Title
\begin{center}
{\LARGE\bfseries\color{bentleyblue} EC 111: Principles of Microeconomics}\\[0.3cm]
{\Large Fall 2026 Syllabus}\\[0.5cm]
{\large Bentley University}
\end{center}

\vspace{0.5cm}

%-----------------------------------------------------------
\section*{Course Description}
%-----------------------------------------------------------

This course introduces the fundamental principles and tools of microeconomics: how individuals and firms make decisions when resources are scarce, and how markets coordinate those decisions. We study demand, supply, and market equilibrium; elasticity; the effects of government policies such as taxes and price controls; the efficiency of markets and the sources of market failure; the costs of production; and firm behavior in different market structures, from perfect competition to monopoly, monopolistic competition, and oligopoly.

The emphasis throughout is on economic reasoning: thinking at the margin, recognizing trade-offs, and using simple graphical and numerical models to make sense of real-world questions in business and policy.

\subsection*{Knowledge}
\begin{itemize}
    \item Scarcity, opportunity cost, trade-offs, and the gains from trade.
    \item Demand, supply, market equilibrium, and how markets respond to shocks.
    \item Elasticity and its role in pricing, revenue, and tax incidence.
    \item Consumer surplus, producer surplus, market efficiency, and the deadweight loss of taxes and price controls.
    \item Externalities, public goods, and the case for government intervention.
    \item Costs of production and firm behavior under perfect competition, monopoly, monopolistic competition, and oligopoly.
\end{itemize}

\subsection*{Skills}
\begin{itemize}
    \item Use supply and demand diagrams to predict the effects of changes in market conditions and policy.
    \item Compute and interpret elasticities, surpluses, costs, and profit-maximizing output and price.
    \item Read and evaluate economic arguments critically and communicate them clearly.
\end{itemize}

\subsection*{Perspectives}
\begin{itemize}
    \item Understand marginal decision making as the core of economic reasoning.
    \item Understand when markets work well, when they fail, and what policy can and cannot fix.
\end{itemize}

%-----------------------------------------------------------
\section*{Class Information}
%-----------------------------------------------------------

\begin{tabularx}{\textwidth}{@{}lX@{}}
\textbf{Instructor:} & Onur Altındağ \\
\textbf{Office:} & RAU 019F \\
\textbf{Email:} & \href{mailto:oaltindag@bentley.edu}{oaltindag@bentley.edu} \\
\textbf{Web:} & \url{https://www.onuraltindag.info/} \\
\textbf{Course page:} & \url{https://www.onuraltindag.info/teaching/ec111/} \\
\textbf{Password:} & \texttt{GoFalcons!} \\[0.3cm]
\textbf{Course:} & \sectionname \\
\textbf{Semester:} & Fall 2026 (August 31, 2026 – December 15, 2026) \\
\textbf{Meeting Time:} & \meetingtime \\
\textbf{Location:} & \classroom \\
\textbf{Delivery Mode:} & In-Person Lecture \\
\end{tabularx}

\vspace{0.3cm}

Office hours are on Mondays and Wednesdays, 11:00 AM -- 12:00 PM and 1:00 -- 3:00 PM, by appointment. Please \href{https://outlook.office.com/bookwithme/user/1c4e4f54e0f0493397e615b2c80a59c7@bentley.edu/meetingtype/rwPThvM0MkK5oGlK_t37nw2?bookingcode=17c9a7b2-69fa-4c48-85ea-ad404642334b&anonymous&ismsaljsauthenabled&ep=mlink}{book a 30-minute slot} to schedule a meeting. If you need to meet urgently or no slots are available, email me directly.

Announcements come by email from me to your Bentley address; check it regularly. Slides, lecture notes, handouts, and mock exams are posted on the course page above after each class. When you email me, put \textbf{EC111} in the subject line.

%-----------------------------------------------------------
\section*{Important Dates and Evaluation}
%-----------------------------------------------------------

\begin{tabularx}{\textwidth}{@{}Xr@{}}
\toprule
\textbf{Component} & \textbf{Weight} \\
\midrule
Weekly assignments + participation & 20\% \\
\textbf{Wednesday, October 7, 2026} – First Midterm & 25\% \\
\textbf{Wednesday, November 11, 2026} – Second Midterm & 25\% \\
\textbf{\finaltime} – Final Exam & 30\% \\
\bottomrule
\end{tabularx}

%-----------------------------------------------------------
\section*{Grading}
%-----------------------------------------------------------

You \textbf{MUST} attend all the midterms and the final as there will be no make-up exams in this course. The midterms are \textbf{NOT} cumulative. If you miss or are likely to miss a midterm due to an emergency, please contact me as soon as possible. You will need to provide supporting documentation/verification of your absence. I will re-weight your final exam accordingly if you have a valid excuse. Please note that family vacations or missing the school shuttle are not valid excuses. If a serious illness, a death in the family, or another difficult circumstance keeps you from an exam or assignment, contact the Office of Student Success (JEN 336, 781.891.2803); they will notify all of your instructors on your behalf.

The final exam \textbf{is} cumulative. If you miss the final exam due to an emergency, you will receive an incomplete for this course. Do not take this class if you know that you will not be able to attend the final exam.

Exams are closed book and closed notes. A basic (non-graphing) calculator is allowed; phones, graphing calculators, and any other device are not, and calculators may not be shared.

You will have \textbf{10 homework assignments} on MindTap, roughly one per week (none in exam weeks), due \textbf{Sunday 11:00 PM}. Your \textbf{two lowest scores are dropped}; the remaining 8 count equally. Each question allows three attempts; an assignment scoring at least \textbf{50\%} receives full credit, otherwise zero. The submission system is automated and will not accept anything after the deadline (even 5 minutes), and there are no exceptions. The questions are algorithmic and slightly different for everyone, so you are encouraged to work together.

\subsection*{Letter grades}
\begin{center}
\begin{tabular}{@{}llllll@{}}
\toprule
A & 94--100 & B+ & 87--89 & C+ & 77--79 \\
A$-$ & 90--93 & B & 83--86 & C & 73--76 \\
 & & B$-$ & 80--82 & C$-$ & 70--72 \\
\midrule
D+ & 67--69 & D & 63--66 & D$-$ & 60--62 \\
F & below 60 & & & & \\
\bottomrule
\end{tabular}
\end{center}
I may lower these cutoffs at the end of the semester; I will not raise them. I do not respond to requests for individual extra credit or for grade adjustments after the fact.

%-----------------------------------------------------------
\section*{Textbook and MindTap}
%-----------------------------------------------------------

The course uses \textbf{MindTap for Mankiw, \textit{Principles of Microeconomics}} (Cengage). The e-book, your homework assignments, and practice questions all live there; a hard copy of the textbook is optional. You \textbf{MUST} register, using your Bentley email, through the link for your section:

\begin{center}
\textbf{\mindtapsection}\\
\url{\mindtaplink}
\end{center}

You have free temporary access until \mindtapaccess; after that you will be prompted to purchase (if your course materials are included in your tuition, no purchase is needed). Registration help: \url{https://startstrong.cengage.com}; technical support: \url{https://support.cengage.com}. Register in the first week; MindTap problems are not an excuse for a missed assignment.

\ifdefined\honors
%-----------------------------------------------------------
\section*{Honors Reading List}
%-----------------------------------------------------------

In addition to the textbook, honors students read and discuss a short set of articles and book chapters over the semester. The list and discussion dates will be posted on the course page in the first weeks of the semester.

\fi
%-----------------------------------------------------------
\section*{Academic Integrity}
%-----------------------------------------------------------

All students are expected to adhere to Bentley's Academic Integrity policy, which includes Bentley's Honor Code (details are in the Undergraduate Student Handbook, the \href{https://catalog.bentley.edu/undergraduate/academic-policies-procedures/}{academic policies section of the catalog}, and the academic integrity course page on Brightspace in which all students and faculty are enrolled). The essence of the policy is that you should not represent someone else's work as your own: no plagiarism, no cheating on exams, no illicit collaboration. Failure to adhere to the policy has serious consequences, including course failure, suspension, or expulsion from the university. The best way to avoid a problem is to consult with me before taking an action that might constitute a violation.

\subsection*{Generative AI}
Nothing in this course is graded on writing, and the homework is algorithmic practice, so the line is simple. You may use generative AI tools (ChatGPT, Claude, Gemini, Copilot, and the like) to study: to explain a concept, to check your reasoning on a practice problem, or to work through a homework question you are stuck on. Treat them as a tutor, not a substitute; the exams are closed book and closed device, and the only way to pass them is to be able to do the work yourself. Any use of AI or any other device during an exam is a violation of the academic integrity policy.

%-----------------------------------------------------------
\section*{Inclusion}
%-----------------------------------------------------------

Everyone in this class has different life experiences and perspectives, and all are valid. As the instructor, I will behave maturely and respectfully in all our class-related engagements, and I expect the same from you. We all adhere to Bentley's standards of conduct, the \href{https://www.bentley.edu/about/mission-and-values}{Bentley Core Values}, in class and in all course work. If you feel that I or anyone in this class has acted outside those values, come and talk to me. If you would rather not, speak with someone in the Office of Student Success (JEN 336, 781.891.2803). My roster shows your name as it appears in Workday; tell me early in the semester if you go by a different name or pronoun.

%-----------------------------------------------------------
\section*{Support and Resources}
%-----------------------------------------------------------

\subsection*{Student Accessibility Services}
Bentley University abides by Section 504 of the Rehabilitation Act of 1973 and the Americans with Disabilities Act of 1990, which stipulate that no student shall be denied the benefits of an education solely by reason of a disability. If you have a hidden or visible disability that may require classroom accommodations, contact \href{https://www.bentley.edu/offices/student-accessibility-services}{Student Accessibility Services} (Office of Student Success, JEN 336, 781.891.2004) within the first four weeks of the semester. Send me your accommodation plan as soon as you have it; accommodations are arranged at least one week before each exam.

\subsection*{The Howard A. Winer '58 Lab for Economics, Accounting and Finance (LEAF)}
The LEAF provides a welcoming and inclusive learning environment where students are encouraged to seek academic support for their Economics, Accounting and Finance courses. Students utilizing the LEAF will find knowledgeable peer tutors ready to help their colleagues to thrive in the Bentley business curriculum. LEAF's free, drop-in tutoring takes place in-person at Morison 220D (rear entrance), opening September 14. For additional information visit \url{https://www.bentley.edu/centers/leaf}. LEAF is a supplement to, not a substitute for, office hours.

\subsection*{Bentley policies and procedures}
All courses at Bentley reflect the university's commitment to a set of core values and practices. You are expected to be familiar with:
\begin{itemize}
    \item \href{https://catalog.bentley.edu/undergraduate/academic-policies-procedures/}{Bentley's honor code and academic integrity system}
    \item \href{https://www.bentley.edu/equity-reporting}{Equity and bias reporting forms and procedures} (Bias Incident Response Team, Office of Institutional Equity)
    \item \href{https://www.bentley.edu/offices/student-accessibility-services}{Student disability accommodations}
    \item \href{https://catalog.bentley.edu/undergraduate/academic-policies-procedures/}{Religious observances policy}
    \item \href{https://www.bentley.edu/about/mission-and-values}{Bentley's core values}
\end{itemize}

\subsection*{Course materials}
Slides, notes, handouts, and exams for this course are for your own use during this semester. They are protected by copyright; do not post or distribute them outside the class without my written consent. The course page is password-protected; do not share the password outside the class.

%-----------------------------------------------------------
\section*{Attendance Policy}
%-----------------------------------------------------------

All students must attend the in-person classes. In-person attendance is essential because it allows for real-time interaction, immediate feedback, and the kind of dialogue that makes learning stick. I do not record lectures, so if you miss a class, you miss the material.

Remote attendance is not an option. This course follows Bentley's \href{https://catalog.bentley.edu/undergraduate/academic-policies-procedures/}{Academic Engagement/Attendance Policy}: students who do not attend or engage during the first week are dropped as no-shows, and two or more consecutive weeks of non-participation can lead to administrative withdrawal. If you must miss class, you are responsible for the material; read the chapter and the posted notes, then bring specific questions to office hours.

\textbf{Weather and closures.} If the university cancels classes, I will email you whether there is an alternative assignment or arrangement for that day.

\textbf{Electronic devices.} Phones away and laptops and tablets closed for the whole class. The work in this course is drawing and reading diagrams, which you learn by doing them by hand, and a screen in your line of sight costs the people around you as much as it costs you. Using a device during class counts against your participation grade. Exceptions are by my permission only; if you have one, you never need to explain it to anyone in the room. Photographing the board at the end of class is fine.

\textbf{Religious observances.} Bentley provides reasonable academic adjustments when a sincerely held religious observance conflicts with a class, exam, or deadline (see the \href{https://downloads.bentley.edu/general/Religious\%20Observances\%202026-2027.pdf}{religious observances calendar} and the \href{https://catalog.bentley.edu/undergraduate/academic-policies-procedures/}{policy}). If you anticipate a conflict, tell me in advance, ideally in the first two weeks of the semester, so we can agree on an adjustment.

%-----------------------------------------------------------
\section*{Tentative Schedule}
%-----------------------------------------------------------

\begin{itemize}
    \item \textbf{Week 1 (Aug 31, Sep 2):} Ten Principles of Economics – Ch. 1 \& Thinking Like an Economist – Ch. 2
    \item \textbf{Week 2 (Sep 9):} Interdependence and the Gains from Trade – Ch. 3 \textit{(no class Sep 7, Labor Day)}
    \item \textbf{Week 3 (Sep 14, 16):} The Market Forces of Supply and Demand – Ch. 4
    \item \textbf{Week 4 (Sep 21, 23):} Elasticity and Its Application – Ch. 5
    \item \textbf{Week 5 (Sep 28, 30):} Supply, Demand, and Government Policies – Ch. 6
    \item \textbf{Week 6 (Oct 5, 7):} Review / \textbf{First Midterm (Oct 7)}
    \item \textbf{Week 7 (Oct 14):} Consumers, Producers, and the Efficiency of Markets – Ch. 7 \textit{(no class Oct 12, Fall Break)}
    \item \textbf{Week 8 (Oct 19, 21):} The Costs of Taxation – Ch. 8 \& International Trade – Ch. 9
    \item \textbf{Week 9 (Oct 26, 28):} Externalities – Ch. 10 \& Public Goods and Common Resources – Ch. 11
    \item \textbf{Week 10 (Nov 2, 4):} The Costs of Production – Ch. 14
    \item \textbf{Week 11 (Nov 9, 11):} Review / \textbf{Second Midterm (Nov 11)}
    \item \textbf{Week 12 (Nov 16, 18):} Firms in Competitive Markets – Ch. 15
    \item \textbf{Week 13 (Nov 23):} Monopoly – Ch. 16 \textit{(no class Nov 25, Thanksgiving Break)}
    \item \textbf{Week 14 (Nov 30, Dec 2):} Monopolistic Competition – Ch. 17 \& Oligopoly – Ch. 18
    \item \textbf{Week 15 (Dec 7):} Review
    \item \textbf{Final Exam:} \finaltime
\end{itemize}

\textit{Note: Chapter references are from Mankiw's \textit{Principles of Microeconomics}, 10th edition (the MindTap e-book). The schedule, and if necessary other parts of this syllabus, may change during the semester; changes are announced in class and by email.}

\end{document}
"
\ifdefined\honors
  \newcommand{\sectioncode}{EC 111-10-H}
  \newcommand{\sectionname}{EC 111-10-H Principles of Microeconomics (Honors)}
  \newcommand{\meetingtime}{Monday/Wednesday, 3:30 PM -- 4:50 PM}
  \newcommand{\classroom}{Smith Technology Center 103}
  \newcommand{\finaltime}{Friday, December 11, 2026, 6:00 PM -- 8:00 PM}
  \newcommand{\mindtapsection}{Section 10-H of EC 111 Principles of Microeconomics Fall 2026}
  \newcommand{\mindtaplink}{https://student.cengage.com/course-link/MTPNS57GW6XN}
  \newcommand{\mindtapaccess}{the end of Sunday, September 13, 2026}
\else
  \newcommand{\sectioncode}{EC 111-20}
  \newcommand{\sectionname}{EC 111-20 Principles of Microeconomics}
  \newcommand{\meetingtime}{Monday/Wednesday, 9:30 AM -- 10:50 AM}
  \newcommand{\classroom}{Jennison Hall 412}
  \newcommand{\finaltime}{Friday, December 11, 2026, 8:00 AM -- 10:00 AM}
  \newcommand{\mindtapsection}{Section 20 of EC 111 Principles of Microeconomics Fall 2026}
  \newcommand{\mindtaplink}{https://student.cengage.com/course-link/MTPP8N7GWX7L}
  \newcommand{\mindtapaccess}{the end of Sunday, September 13, 2026}
\fi

% XeLaTeX for Turkish character support
\usepackage{fontspec}
\setmainfont{DejaVuSerif}[Extension=.ttf, BoldFont=DejaVuSerif-Bold, ItalicFont=DejaVuSerif-Italic, BoldItalicFont=DejaVuSerif-BoldItalic]
\setsansfont{DejaVuSans}[Extension=.ttf, BoldFont=DejaVuSans-Bold, ItalicFont=DejaVuSans-Oblique, BoldItalicFont=DejaVuSans-BoldOblique]

% Page layout
\usepackage[margin=1in]{geometry}
\usepackage{setspace}
\onehalfspacing

% Colors and styling
\usepackage[dvipsnames]{xcolor}
\definecolor{bentleyblue}{RGB}{0, 51, 102}
\definecolor{linkblue}{RGB}{0, 102, 204}

% Headers and sections
\usepackage{titlesec}
\titleformat{\section}{\Large\bfseries\color{bentleyblue}}{\thesection}{1em}{}[\titlerule]
\titleformat{\subsection}{\large\bfseries\color{bentleyblue!80}}{\thesubsection}{1em}{}
\titleformat{\subsubsection}{\normalsize\bfseries\color{bentleyblue!60}}{\thesubsubsection}{1em}{}

% Lists
\usepackage{enumitem}
\setlist[itemize]{leftmargin=*, itemsep=2pt, parsep=0pt}
\setlist[enumerate]{leftmargin=*, itemsep=2pt, parsep=0pt}

% Tables
\usepackage{booktabs}
\usepackage{tabularx}

% Links
\usepackage{hyperref}
\hypersetup{
    colorlinks=true,
    linkcolor=linkblue,
    urlcolor=linkblue,
    pdfauthor={Onur Altındağ},
    pdftitle={EC111 Principles of Microeconomics - Fall 2026 Syllabus}
}

% Header/Footer
\usepackage{fancyhdr}
\pagestyle{fancy}
\fancyhf{}
\fancyhead[L]{\small\color{gray}\sectioncode{} Principles of Microeconomics}
\fancyhead[R]{\small\color{gray}Fall 2026}
\fancyfoot[C]{\thepage}
\renewcommand{\headrulewidth}{0.4pt}

% Remove paragraph indent, add space between paragraphs
\setlength{\parindent}{0pt}
\setlength{\parskip}{6pt}

\begin{document}

% Title
\begin{center}
{\LARGE\bfseries\color{bentleyblue} EC 111: Principles of Microeconomics}\\[0.3cm]
{\Large Fall 2026 Syllabus}\\[0.5cm]
{\large Bentley University}
\end{center}

\vspace{0.5cm}

%-----------------------------------------------------------
\section*{Course Description}
%-----------------------------------------------------------

This course introduces the fundamental principles and tools of microeconomics: how individuals and firms make decisions when resources are scarce, and how markets coordinate those decisions. We study demand, supply, and market equilibrium; elasticity; the effects of government policies such as taxes and price controls; the efficiency of markets and the sources of market failure; the costs of production; and firm behavior in different market structures, from perfect competition to monopoly, monopolistic competition, and oligopoly.

The emphasis throughout is on economic reasoning: thinking at the margin, recognizing trade-offs, and using simple graphical and numerical models to make sense of real-world questions in business and policy.

\subsection*{Knowledge}
\begin{itemize}
    \item Scarcity, opportunity cost, trade-offs, and the gains from trade.
    \item Demand, supply, market equilibrium, and how markets respond to shocks.
    \item Elasticity and its role in pricing, revenue, and tax incidence.
    \item Consumer surplus, producer surplus, market efficiency, and the deadweight loss of taxes and price controls.
    \item Externalities, public goods, and the case for government intervention.
    \item Costs of production and firm behavior under perfect competition, monopoly, monopolistic competition, and oligopoly.
\end{itemize}

\subsection*{Skills}
\begin{itemize}
    \item Use supply and demand diagrams to predict the effects of changes in market conditions and policy.
    \item Compute and interpret elasticities, surpluses, costs, and profit-maximizing output and price.
    \item Read and evaluate economic arguments critically and communicate them clearly.
\end{itemize}

\subsection*{Perspectives}
\begin{itemize}
    \item Understand marginal decision making as the core of economic reasoning.
    \item Understand when markets work well, when they fail, and what policy can and cannot fix.
\end{itemize}

%-----------------------------------------------------------
\section*{Class Information}
%-----------------------------------------------------------

\begin{tabularx}{\textwidth}{@{}lX@{}}
\textbf{Instructor:} & Onur Altındağ \\
\textbf{Office:} & RAU 019F \\
\textbf{Email:} & \href{mailto:oaltindag@bentley.edu}{oaltindag@bentley.edu} \\
\textbf{Web:} & \url{https://www.onuraltindag.info/} \\
\textbf{Course page:} & \url{https://www.onuraltindag.info/teaching/ec111/} \\
\textbf{Password:} & \texttt{GoFalcons!} \\[0.3cm]
\textbf{Course:} & \sectionname \\
\textbf{Semester:} & Fall 2026 (August 31, 2026 – December 15, 2026) \\
\textbf{Meeting Time:} & \meetingtime \\
\textbf{Location:} & \classroom \\
\textbf{Delivery Mode:} & In-Person Lecture \\
\end{tabularx}

\vspace{0.3cm}

Office hours are on Mondays and Wednesdays, 11:00 AM -- 12:00 PM and 1:00 -- 3:00 PM, by appointment. Please \href{https://outlook.office.com/bookwithme/user/1c4e4f54e0f0493397e615b2c80a59c7@bentley.edu/meetingtype/rwPThvM0MkK5oGlK_t37nw2?bookingcode=17c9a7b2-69fa-4c48-85ea-ad404642334b&anonymous&ismsaljsauthenabled&ep=mlink}{book a 30-minute slot} to schedule a meeting. If you need to meet urgently or no slots are available, email me directly.

Announcements come by email from me to your Bentley address; check it regularly. Slides, lecture notes, handouts, and mock exams are posted on the course page above after each class. When you email me, put \textbf{EC111} in the subject line.

%-----------------------------------------------------------
\section*{Important Dates and Evaluation}
%-----------------------------------------------------------

\begin{tabularx}{\textwidth}{@{}Xr@{}}
\toprule
\textbf{Component} & \textbf{Weight} \\
\midrule
Weekly assignments + participation & 20\% \\
\textbf{Wednesday, October 7, 2026} – First Midterm & 25\% \\
\textbf{Wednesday, November 11, 2026} – Second Midterm & 25\% \\
\textbf{\finaltime} – Final Exam & 30\% \\
\bottomrule
\end{tabularx}

%-----------------------------------------------------------
\section*{Grading}
%-----------------------------------------------------------

You \textbf{MUST} attend all the midterms and the final as there will be no make-up exams in this course. The midterms are \textbf{NOT} cumulative. If you miss or are likely to miss a midterm due to an emergency, please contact me as soon as possible. You will need to provide supporting documentation/verification of your absence. I will re-weight your final exam accordingly if you have a valid excuse. Please note that family vacations or missing the school shuttle are not valid excuses. If a serious illness, a death in the family, or another difficult circumstance keeps you from an exam or assignment, contact the Office of Student Success (JEN 336, 781.891.2803); they will notify all of your instructors on your behalf.

The final exam \textbf{is} cumulative. If you miss the final exam due to an emergency, you will receive an incomplete for this course. Do not take this class if you know that you will not be able to attend the final exam.

Exams are closed book and closed notes. A basic (non-graphing) calculator is allowed; phones, graphing calculators, and any other device are not, and calculators may not be shared.

You will have \textbf{10 homework assignments} on MindTap, roughly one per week (none in exam weeks), due \textbf{Sunday 11:00 PM}. Your \textbf{two lowest scores are dropped}; the remaining 8 count equally. Each question allows three attempts; an assignment scoring at least \textbf{50\%} receives full credit, otherwise zero. The submission system is automated and will not accept anything after the deadline (even 5 minutes), and there are no exceptions. The questions are algorithmic and slightly different for everyone, so you are encouraged to work together.

\subsection*{Letter grades}
\begin{center}
\begin{tabular}{@{}llllll@{}}
\toprule
A & 94--100 & B+ & 87--89 & C+ & 77--79 \\
A$-$ & 90--93 & B & 83--86 & C & 73--76 \\
 & & B$-$ & 80--82 & C$-$ & 70--72 \\
\midrule
D+ & 67--69 & D & 63--66 & D$-$ & 60--62 \\
F & below 60 & & & & \\
\bottomrule
\end{tabular}
\end{center}
I may lower these cutoffs at the end of the semester; I will not raise them. I do not respond to requests for individual extra credit or for grade adjustments after the fact.

%-----------------------------------------------------------
\section*{Textbook and MindTap}
%-----------------------------------------------------------

The course uses \textbf{MindTap for Mankiw, \textit{Principles of Microeconomics}} (Cengage). The e-book, your homework assignments, and practice questions all live there; a hard copy of the textbook is optional. You \textbf{MUST} register, using your Bentley email, through the link for your section:

\begin{center}
\textbf{\mindtapsection}\\
\url{\mindtaplink}
\end{center}

You have free temporary access until \mindtapaccess; after that you will be prompted to purchase (if your course materials are included in your tuition, no purchase is needed). Registration help: \url{https://startstrong.cengage.com}; technical support: \url{https://support.cengage.com}. Register in the first week; MindTap problems are not an excuse for a missed assignment.

\ifdefined\honors
%-----------------------------------------------------------
\section*{Honors Reading List}
%-----------------------------------------------------------

In addition to the textbook, honors students read and discuss a short set of articles and book chapters over the semester. The list and discussion dates will be posted on the course page in the first weeks of the semester.

\fi
%-----------------------------------------------------------
\section*{Academic Integrity}
%-----------------------------------------------------------

All students are expected to adhere to Bentley's Academic Integrity policy, which includes Bentley's Honor Code (details are in the Undergraduate Student Handbook, the \href{https://catalog.bentley.edu/undergraduate/academic-policies-procedures/}{academic policies section of the catalog}, and the academic integrity course page on Brightspace in which all students and faculty are enrolled). The essence of the policy is that you should not represent someone else's work as your own: no plagiarism, no cheating on exams, no illicit collaboration. Failure to adhere to the policy has serious consequences, including course failure, suspension, or expulsion from the university. The best way to avoid a problem is to consult with me before taking an action that might constitute a violation.

\subsection*{Generative AI}
Nothing in this course is graded on writing, and the homework is algorithmic practice, so the line is simple. You may use generative AI tools (ChatGPT, Claude, Gemini, Copilot, and the like) to study: to explain a concept, to check your reasoning on a practice problem, or to work through a homework question you are stuck on. Treat them as a tutor, not a substitute; the exams are closed book and closed device, and the only way to pass them is to be able to do the work yourself. Any use of AI or any other device during an exam is a violation of the academic integrity policy.

%-----------------------------------------------------------
\section*{Inclusion}
%-----------------------------------------------------------

Everyone in this class has different life experiences and perspectives, and all are valid. As the instructor, I will behave maturely and respectfully in all our class-related engagements, and I expect the same from you. We all adhere to Bentley's standards of conduct, the \href{https://www.bentley.edu/about/mission-and-values}{Bentley Core Values}, in class and in all course work. If you feel that I or anyone in this class has acted outside those values, come and talk to me. If you would rather not, speak with someone in the Office of Student Success (JEN 336, 781.891.2803). My roster shows your name as it appears in Workday; tell me early in the semester if you go by a different name or pronoun.

%-----------------------------------------------------------
\section*{Support and Resources}
%-----------------------------------------------------------

\subsection*{Student Accessibility Services}
Bentley University abides by Section 504 of the Rehabilitation Act of 1973 and the Americans with Disabilities Act of 1990, which stipulate that no student shall be denied the benefits of an education solely by reason of a disability. If you have a hidden or visible disability that may require classroom accommodations, contact \href{https://www.bentley.edu/offices/student-accessibility-services}{Student Accessibility Services} (Office of Student Success, JEN 336, 781.891.2004) within the first four weeks of the semester. Send me your accommodation plan as soon as you have it; accommodations are arranged at least one week before each exam.

\subsection*{The Howard A. Winer '58 Lab for Economics, Accounting and Finance (LEAF)}
The LEAF provides a welcoming and inclusive learning environment where students are encouraged to seek academic support for their Economics, Accounting and Finance courses. Students utilizing the LEAF will find knowledgeable peer tutors ready to help their colleagues to thrive in the Bentley business curriculum. LEAF's free, drop-in tutoring takes place in-person at Morison 220D (rear entrance), opening September 14. For additional information visit \url{https://www.bentley.edu/centers/leaf}. LEAF is a supplement to, not a substitute for, office hours.

\subsection*{Bentley policies and procedures}
All courses at Bentley reflect the university's commitment to a set of core values and practices. You are expected to be familiar with:
\begin{itemize}
    \item \href{https://catalog.bentley.edu/undergraduate/academic-policies-procedures/}{Bentley's honor code and academic integrity system}
    \item \href{https://www.bentley.edu/equity-reporting}{Equity and bias reporting forms and procedures} (Bias Incident Response Team, Office of Institutional Equity)
    \item \href{https://www.bentley.edu/offices/student-accessibility-services}{Student disability accommodations}
    \item \href{https://catalog.bentley.edu/undergraduate/academic-policies-procedures/}{Religious observances policy}
    \item \href{https://www.bentley.edu/about/mission-and-values}{Bentley's core values}
\end{itemize}

\subsection*{Course materials}
Slides, notes, handouts, and exams for this course are for your own use during this semester. They are protected by copyright; do not post or distribute them outside the class without my written consent. The course page is password-protected; do not share the password outside the class.

%-----------------------------------------------------------
\section*{Attendance Policy}
%-----------------------------------------------------------

All students must attend the in-person classes. In-person attendance is essential because it allows for real-time interaction, immediate feedback, and the kind of dialogue that makes learning stick. I do not record lectures, so if you miss a class, you miss the material.

Remote attendance is not an option. This course follows Bentley's \href{https://catalog.bentley.edu/undergraduate/academic-policies-procedures/}{Academic Engagement/Attendance Policy}: students who do not attend or engage during the first week are dropped as no-shows, and two or more consecutive weeks of non-participation can lead to administrative withdrawal. If you must miss class, you are responsible for the material; read the chapter and the posted notes, then bring specific questions to office hours.

\textbf{Weather and closures.} If the university cancels classes, I will email you whether there is an alternative assignment or arrangement for that day.

\textbf{Electronic devices.} Phones away and laptops and tablets closed for the whole class. The work in this course is drawing and reading diagrams, which you learn by doing them by hand, and a screen in your line of sight costs the people around you as much as it costs you. Using a device during class counts against your participation grade. Exceptions are by my permission only; if you have one, you never need to explain it to anyone in the room. Photographing the board at the end of class is fine.

\textbf{Religious observances.} Bentley provides reasonable academic adjustments when a sincerely held religious observance conflicts with a class, exam, or deadline (see the \href{https://downloads.bentley.edu/general/Religious\%20Observances\%202026-2027.pdf}{religious observances calendar} and the \href{https://catalog.bentley.edu/undergraduate/academic-policies-procedures/}{policy}). If you anticipate a conflict, tell me in advance, ideally in the first two weeks of the semester, so we can agree on an adjustment.

%-----------------------------------------------------------
\section*{Tentative Schedule}
%-----------------------------------------------------------

\begin{itemize}
    \item \textbf{Week 1 (Aug 31, Sep 2):} Ten Principles of Economics – Ch. 1 \& Thinking Like an Economist – Ch. 2
    \item \textbf{Week 2 (Sep 9):} Interdependence and the Gains from Trade – Ch. 3 \textit{(no class Sep 7, Labor Day)}
    \item \textbf{Week 3 (Sep 14, 16):} The Market Forces of Supply and Demand – Ch. 4
    \item \textbf{Week 4 (Sep 21, 23):} Elasticity and Its Application – Ch. 5
    \item \textbf{Week 5 (Sep 28, 30):} Supply, Demand, and Government Policies – Ch. 6
    \item \textbf{Week 6 (Oct 5, 7):} Review / \textbf{First Midterm (Oct 7)}
    \item \textbf{Week 7 (Oct 14):} Consumers, Producers, and the Efficiency of Markets – Ch. 7 \textit{(no class Oct 12, Fall Break)}
    \item \textbf{Week 8 (Oct 19, 21):} The Costs of Taxation – Ch. 8 \& International Trade – Ch. 9
    \item \textbf{Week 9 (Oct 26, 28):} Externalities – Ch. 10 \& Public Goods and Common Resources – Ch. 11
    \item \textbf{Week 10 (Nov 2, 4):} The Costs of Production – Ch. 14
    \item \textbf{Week 11 (Nov 9, 11):} Review / \textbf{Second Midterm (Nov 11)}
    \item \textbf{Week 12 (Nov 16, 18):} Firms in Competitive Markets – Ch. 15
    \item \textbf{Week 13 (Nov 23):} Monopoly – Ch. 16 \textit{(no class Nov 25, Thanksgiving Break)}
    \item \textbf{Week 14 (Nov 30, Dec 2):} Monopolistic Competition – Ch. 17 \& Oligopoly – Ch. 18
    \item \textbf{Week 15 (Dec 7):} Review
    \item \textbf{Final Exam:} \finaltime
\end{itemize}

\textit{Note: Chapter references are from Mankiw's \textit{Principles of Microeconomics}, 10th edition (the MindTap e-book). The schedule, and if necessary other parts of this syllabus, may change during the semester; changes are announced in class and by email.}

\end{document}
"
\ifdefined\honors
  \newcommand{\sectioncode}{EC 111-10-H}
  \newcommand{\sectionname}{EC 111-10-H Principles of Microeconomics (Honors)}
  \newcommand{\meetingtime}{Monday/Wednesday, 3:30 PM -- 4:50 PM}
  \newcommand{\classroom}{Smith Technology Center 103}
  \newcommand{\finaltime}{Friday, December 11, 2026, 6:00 PM -- 8:00 PM}
  \newcommand{\mindtapsection}{Section 10-H of EC 111 Principles of Microeconomics Fall 2026}
  \newcommand{\mindtaplink}{https://student.cengage.com/course-link/MTPNS57GW6XN}
  \newcommand{\mindtapaccess}{the end of Sunday, September 13, 2026}
\else
  \newcommand{\sectioncode}{EC 111-20}
  \newcommand{\sectionname}{EC 111-20 Principles of Microeconomics}
  \newcommand{\meetingtime}{Monday/Wednesday, 9:30 AM -- 10:50 AM}
  \newcommand{\classroom}{Jennison Hall 412}
  \newcommand{\finaltime}{Friday, December 11, 2026, 8:00 AM -- 10:00 AM}
  \newcommand{\mindtapsection}{Section 20 of EC 111 Principles of Microeconomics Fall 2026}
  \newcommand{\mindtaplink}{https://student.cengage.com/course-link/MTPP8N7GWX7L}
  \newcommand{\mindtapaccess}{the end of Sunday, September 13, 2026}
\fi

% XeLaTeX for Turkish character support
\usepackage{fontspec}
\setmainfont{DejaVuSerif}[Extension=.ttf, BoldFont=DejaVuSerif-Bold, ItalicFont=DejaVuSerif-Italic, BoldItalicFont=DejaVuSerif-BoldItalic]
\setsansfont{DejaVuSans}[Extension=.ttf, BoldFont=DejaVuSans-Bold, ItalicFont=DejaVuSans-Oblique, BoldItalicFont=DejaVuSans-BoldOblique]

% Page layout
\usepackage[margin=1in]{geometry}
\usepackage{setspace}
\onehalfspacing

% Colors and styling
\usepackage[dvipsnames]{xcolor}
\definecolor{bentleyblue}{RGB}{0, 51, 102}
\definecolor{linkblue}{RGB}{0, 102, 204}

% Headers and sections
\usepackage{titlesec}
\titleformat{\section}{\Large\bfseries\color{bentleyblue}}{\thesection}{1em}{}[\titlerule]
\titleformat{\subsection}{\large\bfseries\color{bentleyblue!80}}{\thesubsection}{1em}{}
\titleformat{\subsubsection}{\normalsize\bfseries\color{bentleyblue!60}}{\thesubsubsection}{1em}{}

% Lists
\usepackage{enumitem}
\setlist[itemize]{leftmargin=*, itemsep=2pt, parsep=0pt}
\setlist[enumerate]{leftmargin=*, itemsep=2pt, parsep=0pt}

% Tables
\usepackage{booktabs}
\usepackage{tabularx}

% Links
\usepackage{hyperref}
\hypersetup{
    colorlinks=true,
    linkcolor=linkblue,
    urlcolor=linkblue,
    pdfauthor={Onur Altındağ},
    pdftitle={EC111 Principles of Microeconomics - Fall 2026 Syllabus}
}

% Header/Footer
\usepackage{fancyhdr}
\pagestyle{fancy}
\fancyhf{}
\fancyhead[L]{\small\color{gray}\sectioncode{} Principles of Microeconomics}
\fancyhead[R]{\small\color{gray}Fall 2026}
\fancyfoot[C]{\thepage}
\renewcommand{\headrulewidth}{0.4pt}

% Remove paragraph indent, add space between paragraphs
\setlength{\parindent}{0pt}
\setlength{\parskip}{6pt}

\begin{document}

% Title
\begin{center}
{\LARGE\bfseries\color{bentleyblue} EC 111: Principles of Microeconomics}\\[0.3cm]
{\Large Fall 2026 Syllabus}\\[0.5cm]
{\large Bentley University}
\end{center}

\vspace{0.5cm}

%-----------------------------------------------------------
\section*{Course Description}
%-----------------------------------------------------------

This course introduces the fundamental principles and tools of microeconomics: how individuals and firms make decisions when resources are scarce, and how markets coordinate those decisions. We study demand, supply, and market equilibrium; elasticity; the effects of government policies such as taxes and price controls; the efficiency of markets and the sources of market failure; the costs of production; and firm behavior in different market structures, from perfect competition to monopoly, monopolistic competition, and oligopoly.

The emphasis throughout is on economic reasoning: thinking at the margin, recognizing trade-offs, and using simple graphical and numerical models to make sense of real-world questions in business and policy.

\subsection*{Knowledge}
\begin{itemize}
    \item Scarcity, opportunity cost, trade-offs, and the gains from trade.
    \item Demand, supply, market equilibrium, and how markets respond to shocks.
    \item Elasticity and its role in pricing, revenue, and tax incidence.
    \item Consumer surplus, producer surplus, market efficiency, and the deadweight loss of taxes and price controls.
    \item Externalities, public goods, and the case for government intervention.
    \item Costs of production and firm behavior under perfect competition, monopoly, monopolistic competition, and oligopoly.
\end{itemize}

\subsection*{Skills}
\begin{itemize}
    \item Use supply and demand diagrams to predict the effects of changes in market conditions and policy.
    \item Compute and interpret elasticities, surpluses, costs, and profit-maximizing output and price.
    \item Read and evaluate economic arguments critically and communicate them clearly.
\end{itemize}

\subsection*{Perspectives}
\begin{itemize}
    \item Understand marginal decision making as the core of economic reasoning.
    \item Understand when markets work well, when they fail, and what policy can and cannot fix.
\end{itemize}

%-----------------------------------------------------------
\section*{Class Information}
%-----------------------------------------------------------

\begin{tabularx}{\textwidth}{@{}lX@{}}
\textbf{Instructor:} & Onur Altındağ \\
\textbf{Office:} & RAU 019F \\
\textbf{Email:} & \href{mailto:oaltindag@bentley.edu}{oaltindag@bentley.edu} \\
\textbf{Web:} & \url{https://www.onuraltindag.info/} \\
\textbf{Course page:} & \url{https://www.onuraltindag.info/teaching/ec111/} \\
\textbf{Password:} & \texttt{GoFalcons!} \\[0.3cm]
\textbf{Course:} & \sectionname \\
\textbf{Semester:} & Fall 2026 (August 31, 2026 – December 15, 2026) \\
\textbf{Meeting Time:} & \meetingtime \\
\textbf{Location:} & \classroom \\
\textbf{Delivery Mode:} & In-Person Lecture \\
\end{tabularx}

\vspace{0.3cm}

Office hours are on Mondays and Wednesdays, 11:00 AM -- 12:00 PM and 1:00 -- 3:00 PM, by appointment. Please \href{https://outlook.office.com/bookwithme/user/1c4e4f54e0f0493397e615b2c80a59c7@bentley.edu/meetingtype/rwPThvM0MkK5oGlK_t37nw2?bookingcode=17c9a7b2-69fa-4c48-85ea-ad404642334b&anonymous&ismsaljsauthenabled&ep=mlink}{book a 30-minute slot} to schedule a meeting. If you need to meet urgently or no slots are available, email me directly.

Announcements come by email from me to your Bentley address; check it regularly. Slides, lecture notes, handouts, and mock exams are posted on the course page above after each class. When you email me, put \textbf{EC111} in the subject line.

%-----------------------------------------------------------
\section*{Important Dates and Evaluation}
%-----------------------------------------------------------

\begin{tabularx}{\textwidth}{@{}Xr@{}}
\toprule
\textbf{Component} & \textbf{Weight} \\
\midrule
Weekly assignments + participation & 20\% \\
\textbf{Wednesday, October 7, 2026} – First Midterm & 25\% \\
\textbf{Wednesday, November 11, 2026} – Second Midterm & 25\% \\
\textbf{\finaltime} – Final Exam & 30\% \\
\bottomrule
\end{tabularx}

%-----------------------------------------------------------
\section*{Grading}
%-----------------------------------------------------------

You \textbf{MUST} attend all the midterms and the final as there will be no make-up exams in this course. The midterms are \textbf{NOT} cumulative. If you miss or are likely to miss a midterm due to an emergency, please contact me as soon as possible. You will need to provide supporting documentation/verification of your absence. I will re-weight your final exam accordingly if you have a valid excuse. Please note that family vacations or missing the school shuttle are not valid excuses. If a serious illness, a death in the family, or another difficult circumstance keeps you from an exam or assignment, contact the Office of Student Success (JEN 336, 781.891.2803); they will notify all of your instructors on your behalf.

The final exam \textbf{is} cumulative. If you miss the final exam due to an emergency, you will receive an incomplete for this course. Do not take this class if you know that you will not be able to attend the final exam.

Exams are closed book and closed notes. A basic (non-graphing) calculator is allowed; phones, graphing calculators, and any other device are not, and calculators may not be shared.

You will have \textbf{10 homework assignments} on MindTap, roughly one per week (none in exam weeks), due \textbf{Sunday 11:00 PM}. Your \textbf{two lowest scores are dropped}; the remaining 8 count equally. Each question allows three attempts; an assignment scoring at least \textbf{50\%} receives full credit, otherwise zero. The submission system is automated and will not accept anything after the deadline (even 5 minutes), and there are no exceptions. The questions are algorithmic and slightly different for everyone, so you are encouraged to work together.

\subsection*{Letter grades}
\begin{center}
\begin{tabular}{@{}llllll@{}}
\toprule
A & 94--100 & B+ & 87--89 & C+ & 77--79 \\
A$-$ & 90--93 & B & 83--86 & C & 73--76 \\
 & & B$-$ & 80--82 & C$-$ & 70--72 \\
\midrule
D+ & 67--69 & D & 63--66 & D$-$ & 60--62 \\
F & below 60 & & & & \\
\bottomrule
\end{tabular}
\end{center}
I may lower these cutoffs at the end of the semester; I will not raise them. I do not respond to requests for individual extra credit or for grade adjustments after the fact.

%-----------------------------------------------------------
\section*{Textbook and MindTap}
%-----------------------------------------------------------

The course uses \textbf{MindTap for Mankiw, \textit{Principles of Microeconomics}} (Cengage). The e-book, your homework assignments, and practice questions all live there; a hard copy of the textbook is optional. You \textbf{MUST} register, using your Bentley email, through the link for your section:

\begin{center}
\textbf{\mindtapsection}\\
\url{\mindtaplink}
\end{center}

You have free temporary access until \mindtapaccess; after that you will be prompted to purchase (if your course materials are included in your tuition, no purchase is needed). Registration help: \url{https://startstrong.cengage.com}; technical support: \url{https://support.cengage.com}. Register in the first week; MindTap problems are not an excuse for a missed assignment.

\ifdefined\honors
%-----------------------------------------------------------
\section*{Honors Reading List}
%-----------------------------------------------------------

In addition to the textbook, honors students read and discuss a short set of articles and book chapters over the semester. The list and discussion dates will be posted on the course page in the first weeks of the semester.

\fi
%-----------------------------------------------------------
\section*{Academic Integrity}
%-----------------------------------------------------------

All students are expected to adhere to Bentley's Academic Integrity policy, which includes Bentley's Honor Code (details are in the Undergraduate Student Handbook, the \href{https://catalog.bentley.edu/undergraduate/academic-policies-procedures/}{academic policies section of the catalog}, and the academic integrity course page on Brightspace in which all students and faculty are enrolled). The essence of the policy is that you should not represent someone else's work as your own: no plagiarism, no cheating on exams, no illicit collaboration. Failure to adhere to the policy has serious consequences, including course failure, suspension, or expulsion from the university. The best way to avoid a problem is to consult with me before taking an action that might constitute a violation.

\subsection*{Generative AI}
Nothing in this course is graded on writing, and the homework is algorithmic practice, so the line is simple. You may use generative AI tools (ChatGPT, Claude, Gemini, Copilot, and the like) to study: to explain a concept, to check your reasoning on a practice problem, or to work through a homework question you are stuck on. Treat them as a tutor, not a substitute; the exams are closed book and closed device, and the only way to pass them is to be able to do the work yourself. Any use of AI or any other device during an exam is a violation of the academic integrity policy.

%-----------------------------------------------------------
\section*{Inclusion}
%-----------------------------------------------------------

Everyone in this class has different life experiences and perspectives, and all are valid. As the instructor, I will behave maturely and respectfully in all our class-related engagements, and I expect the same from you. We all adhere to Bentley's standards of conduct, the \href{https://www.bentley.edu/about/mission-and-values}{Bentley Core Values}, in class and in all course work. If you feel that I or anyone in this class has acted outside those values, come and talk to me. If you would rather not, speak with someone in the Office of Student Success (JEN 336, 781.891.2803). My roster shows your name as it appears in Workday; tell me early in the semester if you go by a different name or pronoun.

%-----------------------------------------------------------
\section*{Support and Resources}
%-----------------------------------------------------------

\subsection*{Student Accessibility Services}
Bentley University abides by Section 504 of the Rehabilitation Act of 1973 and the Americans with Disabilities Act of 1990, which stipulate that no student shall be denied the benefits of an education solely by reason of a disability. If you have a hidden or visible disability that may require classroom accommodations, contact \href{https://www.bentley.edu/offices/student-accessibility-services}{Student Accessibility Services} (Office of Student Success, JEN 336, 781.891.2004) within the first four weeks of the semester. Send me your accommodation plan as soon as you have it; accommodations are arranged at least one week before each exam.

\subsection*{The Howard A. Winer '58 Lab for Economics, Accounting and Finance (LEAF)}
The LEAF provides a welcoming and inclusive learning environment where students are encouraged to seek academic support for their Economics, Accounting and Finance courses. Students utilizing the LEAF will find knowledgeable peer tutors ready to help their colleagues to thrive in the Bentley business curriculum. LEAF's free, drop-in tutoring takes place in-person at Morison 220D (rear entrance), opening September 14. For additional information visit \url{https://www.bentley.edu/centers/leaf}. LEAF is a supplement to, not a substitute for, office hours.

\subsection*{Bentley policies and procedures}
All courses at Bentley reflect the university's commitment to a set of core values and practices. You are expected to be familiar with:
\begin{itemize}
    \item \href{https://catalog.bentley.edu/undergraduate/academic-policies-procedures/}{Bentley's honor code and academic integrity system}
    \item \href{https://www.bentley.edu/equity-reporting}{Equity and bias reporting forms and procedures} (Bias Incident Response Team, Office of Institutional Equity)
    \item \href{https://www.bentley.edu/offices/student-accessibility-services}{Student disability accommodations}
    \item \href{https://catalog.bentley.edu/undergraduate/academic-policies-procedures/}{Religious observances policy}
    \item \href{https://www.bentley.edu/about/mission-and-values}{Bentley's core values}
\end{itemize}

\subsection*{Course materials}
Slides, notes, handouts, and exams for this course are for your own use during this semester. They are protected by copyright; do not post or distribute them outside the class without my written consent. The course page is password-protected; do not share the password outside the class.

%-----------------------------------------------------------
\section*{Attendance Policy}
%-----------------------------------------------------------

All students must attend the in-person classes. In-person attendance is essential because it allows for real-time interaction, immediate feedback, and the kind of dialogue that makes learning stick. I do not record lectures, so if you miss a class, you miss the material.

Remote attendance is not an option. This course follows Bentley's \href{https://catalog.bentley.edu/undergraduate/academic-policies-procedures/}{Academic Engagement/Attendance Policy}: students who do not attend or engage during the first week are dropped as no-shows, and two or more consecutive weeks of non-participation can lead to administrative withdrawal. If you must miss class, you are responsible for the material; read the chapter and the posted notes, then bring specific questions to office hours.

\textbf{Weather and closures.} If the university cancels classes, I will email you whether there is an alternative assignment or arrangement for that day.

\textbf{Electronic devices.} Phones away and laptops and tablets closed for the whole class. The work in this course is drawing and reading diagrams, which you learn by doing them by hand, and a screen in your line of sight costs the people around you as much as it costs you. Using a device during class counts against your participation grade. Exceptions are by my permission only; if you have one, you never need to explain it to anyone in the room. Photographing the board at the end of class is fine.

\textbf{Religious observances.} Bentley provides reasonable academic adjustments when a sincerely held religious observance conflicts with a class, exam, or deadline (see the \href{https://downloads.bentley.edu/general/Religious\%20Observances\%202026-2027.pdf}{religious observances calendar} and the \href{https://catalog.bentley.edu/undergraduate/academic-policies-procedures/}{policy}). If you anticipate a conflict, tell me in advance, ideally in the first two weeks of the semester, so we can agree on an adjustment.

%-----------------------------------------------------------
\section*{Tentative Schedule}
%-----------------------------------------------------------

\begin{itemize}
    \item \textbf{Week 1 (Aug 31, Sep 2):} Ten Principles of Economics – Ch. 1 \& Thinking Like an Economist – Ch. 2
    \item \textbf{Week 2 (Sep 9):} Interdependence and the Gains from Trade – Ch. 3 \textit{(no class Sep 7, Labor Day)}
    \item \textbf{Week 3 (Sep 14, 16):} The Market Forces of Supply and Demand – Ch. 4
    \item \textbf{Week 4 (Sep 21, 23):} Elasticity and Its Application – Ch. 5
    \item \textbf{Week 5 (Sep 28, 30):} Supply, Demand, and Government Policies – Ch. 6
    \item \textbf{Week 6 (Oct 5, 7):} Review / \textbf{First Midterm (Oct 7)}
    \item \textbf{Week 7 (Oct 14):} Consumers, Producers, and the Efficiency of Markets – Ch. 7 \textit{(no class Oct 12, Fall Break)}
    \item \textbf{Week 8 (Oct 19, 21):} The Costs of Taxation – Ch. 8 \& International Trade – Ch. 9
    \item \textbf{Week 9 (Oct 26, 28):} Externalities – Ch. 10 \& Public Goods and Common Resources – Ch. 11
    \item \textbf{Week 10 (Nov 2, 4):} The Costs of Production – Ch. 14
    \item \textbf{Week 11 (Nov 9, 11):} Review / \textbf{Second Midterm (Nov 11)}
    \item \textbf{Week 12 (Nov 16, 18):} Firms in Competitive Markets – Ch. 15
    \item \textbf{Week 13 (Nov 23):} Monopoly – Ch. 16 \textit{(no class Nov 25, Thanksgiving Break)}
    \item \textbf{Week 14 (Nov 30, Dec 2):} Monopolistic Competition – Ch. 17 \& Oligopoly – Ch. 18
    \item \textbf{Week 15 (Dec 7):} Review
    \item \textbf{Final Exam:} \finaltime
\end{itemize}

\textit{Note: Chapter references are from Mankiw's \textit{Principles of Microeconomics}, 10th edition (the MindTap e-book). The schedule, and if necessary other parts of this syllabus, may change during the semester; changes are announced in class and by email.}

\end{document}
